% Énergie cinétique et théorème de l'énergie cinétique — Physique 1ère C — Cours signé OnBuch+
% Programme officiel MINESEC. Compiler avec : tectonic N04-energie-cinetique-theoreme.tex
\newcommand{\DOCMATIERE}{Physique}
\newcommand{\DOCNIVEAU}{1ère C}
\newcommand{\DOCMODULE}{Module 2 : Mouvements et interactions mécaniques}
\newcommand{\DOCLECON}{N04}
\newcommand{\DOCTITRE}{Énergie cinétique et théorème de l'énergie cinétique}
% OnBuch+ — préambule « cours signés » ENRICHI (compilé avec tectonic / XeLaTeX)
% Système visuel « Pop » d'OnBuch+ : fond crème (jamais blanc), bordures encre
% épaisses, ombres « dures » décalées, titres Archivo Black en capitales.
% Version MATHÉMATIQUES : graphes (pgfplots/tikz), boîtes pédagogiques
% (définition, propriété, méthode, attention, le savais-tu, exercice résolu,
% exercice, TP, bilan), page de garde et signature OnBuch+ ancrée en bas.
%
% Macros attendues dans main.tex AVANT \input{preamble} :
%   \DOCMATIERE  \DOCNIVEAU  \DOCMODULE  \DOCLECON (numéro)  \DOCTITRE
\documentclass[11pt]{article}

\usepackage{fontspec}
\usepackage[french]{babel}
\usepackage[a4paper,margin=2cm,top=3.4cm,bottom=2.8cm,headheight=1.4cm,footskip=1.3cm]{geometry}
\usepackage{amsmath,amssymb,amsfonts}
\usepackage{mathtools} % \xmapsto, \coloneqq... souvent utilisés par les modèles
\usepackage{cancel} % simplifications barrées (bilans de forces)
% \abs{z} (module, valeur absolue) et \norm{u} (norme), utilisés par les modèles
\providecommand{\abs}{}\let\abs\relax\DeclarePairedDelimiter{\abs}{\lvert}{\rvert}
\providecommand{\norm}{}\let\norm\relax\DeclarePairedDelimiter{\norm}{\lVert}{\rVert}
\usepackage[table]{xcolor} % [table] : fournit aussi \rowcolors (alternance de lignes)
\usepackage{pagecolor}
\usepackage{tikz}
\usetikzlibrary{arrows.meta,positioning,calc,shapes.geometric,shapes.misc,shapes.arrows,decorations.pathmorphing,decorations.pathreplacing,decorations.markings,patterns,shadows,mindmap,trees,fit,backgrounds,matrix,intersections,angles,quotes}
\usepackage{pgfplots}
\usepackage[version=4]{mhchem} % équations nucléaires \ce{^{235}_{92}U}
\usepackage[european,siunitx]{circuitikz} % schémas électriques
\pgfplotsset{compat=1.18}
\usepgfplotslibrary{fillbetween,groupplots} % aires entre courbes (calcul intégral)
\usepackage{enumitem}
\usepackage{fancyhdr}
% [most] charge toutes les bibliothèques tcolorbox (listings, minted, raster,
% documentation...) et gonfle inutilement la mémoire TeX sur un document long
% (~300 boîtes) : on ne charge que ce qui est réellement utilisé.
\usepackage[skins,breakable]{tcolorbox}
\usepackage{graphicx}
\usepackage{colortbl}
\usepackage{booktabs}
\usepackage{tabularx}
\usepackage{diagbox} % case d'en-tête diagonale des tableaux à double entrée
% Colonnes extensibles centrée / gauche / droite (C, L, R), souvent utilisées par les modèles
\newcolumntype{C}{>{\centering\arraybackslash}X}
\newcolumntype{L}{>{\raggedright\arraybackslash}X}
\newcolumntype{R}{>{\raggedleft\arraybackslash}X}
\usepackage{array}
\usepackage{multicol}
\usepackage{siunitx}
% parse-numbers=false : accepte aussi \SI{2,5 \times 10^{-3}}{...}
\sisetup{locale=FR,output-decimal-marker={,},group-minimum-digits=4,parse-numbers=false}
\DeclareSIUnit\ohm{\ensuremath{\symup{\Omega}}} % U+2126 absent de la police
\DeclareSIUnit\division{div} % réglages d'oscilloscope (ms/div, V/div)
\DeclareSIUnit\tr{tr} % tours (vitesse de rotation, ex. \SI{1500}{\tr\per\minute})
\DeclareSIUnit\ch{ch} % cheval-vapeur (puissance moteur, unité usuelle hors SI)
\DeclareSIUnit\franccfa{FCFA} % franc CFA (coûts, contexte camerounais)
\DeclareSIUnit\dioptre{\ensuremath{\delta}} % vergence d'une lentille (dioptrie)
\usepackage{qrcode}
% \degree : commande fréquemment utilisée par les modèles pour un angle ou une
% température en dehors de \SI{}{} (non fournie par les paquets chargés).
\providecommand{\degree}{^{\circ}}
% \qed (fin de démonstration), utilisé par les modèles sans amsthm
\providecommand{\qed}{\hfill$\square$}
% \barre : double barre des valeurs interdites dans un tableau de variations (tkz-tab)
\providecommand{\barre}{\ensuremath{\|}}
% environnement proof (amsthm non chargé) utilisé par les modèles
\ifdefined\proof\else\newenvironment{proof}[1][Démonstration]{\par\noindent\textit{#1.}\ }{\qed\par}\fi
\usepackage{lastpage}
\usepackage{hyperref}
\hypersetup{hidelinks}

% ============================================================
% Palette « Pop » OnBuch+ — mêmes hex que la classe P (plus_ui.dart)
% ============================================================
\definecolor{poporange}{HTML}{F59321}
\definecolor{popdark}{HTML}{E07A0C}
\definecolor{popcream}{HTML}{FFF6E8}
\definecolor{popink}{HTML}{1C1714}
\definecolor{poppurple}{HTML}{7A5AE0}
\definecolor{popgreen}{HTML}{2E9E6A}
\definecolor{poppink}{HTML}{E0457B}
\definecolor{popblue}{HTML}{2F7DE1}
\definecolor{popgold}{HTML}{C98A12}
\definecolor{popmuted}{HTML}{8A7F76}
\definecolor{popline}{HTML}{EADFD2}
\definecolor{popgray}{HTML}{8A7F76}
\colorlet{popgrey}{popgray}
% Alias tolérants (le modèle nomme parfois les couleurs différemment)
\colorlet{popwhite}{white}
\colorlet{popblack}{popink}
\colorlet{popred}{poppink}
\colorlet{popyellow}{popgold}
% Teintes claires pour fonds de tableaux / zones
\colorlet{popblueL}{popblue!10!white}
\colorlet{poporangeL}{poporange!14!white}
\colorlet{popgreenL}{popgreen!12!white}
\colorlet{poppurpleL}{poppurple!12!white}
\colorlet{poppinkL}{poppink!12!white}
\colorlet{popgoldL}{popgold!14!white}

\pagecolor{popcream}

% ============================================================
% Typographie
% ============================================================
\defaultfontfeatures{Ligatures=TeX}
\setmainfont{PlusJakartaSans-Medium.ttf}[Path=../../../fonts/,BoldFont=PlusJakartaSans-Bold.ttf]
% Maths en sans-serif (Fira Math) avec chiffres et lettres droites de Plus
% Jakarta, pour que les formules se fondent dans le texte.
\usepackage[math-style=ISO,bold-style=ISO]{unicode-math}
\setmathfont{FiraMath-Regular.otf}[Path=../../../fonts/]
\setmathfont{PlusJakartaSans-Medium.ttf}[Path=../../../fonts/,range={up/num,up/{Latin,latin},\mathup}]
\usepackage{icomma} % virgule décimale sans espace en mode math
% Caractères mathématiques Unicode absents des polices de texte (Plus Jakarta,
% Archivo Black) : rendus par leur équivalent mathématique, en texte comme en
% formule — sinon ils s'affichent en carré vide (ex. « Arithmétique dans ℤ »).
\usepackage{newunicodechar}
\newunicodechar{ℝ}{\ensuremath{\mathbb{R}}}
\newunicodechar{ℤ}{\ensuremath{\mathbb{Z}}}
\newunicodechar{ℂ}{\ensuremath{\mathbb{C}}}
\newunicodechar{ℕ}{\ensuremath{\mathbb{N}}}
\newunicodechar{ℚ}{\ensuremath{\mathbb{Q}}}
\newunicodechar{𝔻}{\ensuremath{\mathbb{D}}}
\newunicodechar{α}{\ensuremath{\alpha}}
\newunicodechar{β}{\ensuremath{\beta}}
\newunicodechar{γ}{\ensuremath{\gamma}}
\newunicodechar{δ}{\ensuremath{\delta}}
\newunicodechar{ε}{\ensuremath{\varepsilon}}
\newunicodechar{θ}{\ensuremath{\theta}}
\newunicodechar{λ}{\ensuremath{\lambda}}
\newunicodechar{μ}{\ensuremath{\mu}}
\newunicodechar{σ}{\ensuremath{\sigma}}
\newunicodechar{τ}{\ensuremath{\tau}}
\newunicodechar{φ}{\ensuremath{\varphi}}
\newunicodechar{ψ}{\ensuremath{\psi}}
\newunicodechar{ω}{\ensuremath{\omega}}
\newunicodechar{Σ}{\ensuremath{\Sigma}}
% majuscules produites par \MakeUppercase dans les titres (α-aminés -> Α)
\newunicodechar{Α}{\ensuremath{\alpha}}
\newunicodechar{Β}{\ensuremath{\beta}}
\newunicodechar{Γ}{\ensuremath{\Gamma}}
\newunicodechar{Δ}{\ensuremath{\Delta}}
\newunicodechar{Ε}{\ensuremath{\varepsilon}}
\newunicodechar{Θ}{\ensuremath{\Theta}}
\newunicodechar{Λ}{\ensuremath{\Lambda}}
\newunicodechar{Μ}{\ensuremath{\mu}}
\newunicodechar{Τ}{\ensuremath{\tau}}
\newunicodechar{Φ}{\ensuremath{\Phi}}
\newunicodechar{Ψ}{\ensuremath{\Psi}}
\newunicodechar{Ω}{\ensuremath{\Omega}}
\newunicodechar{↦}{\ensuremath{\mapsto}}
\newunicodechar{∈}{\ensuremath{\in}}
\newunicodechar{∉}{\ensuremath{\notin}}
\newunicodechar{∧}{\ensuremath{\wedge}}
\newunicodechar{⇔}{\ensuremath{\Leftrightarrow}}
\newunicodechar{⟺}{\ensuremath{\Longleftrightarrow}}
\newunicodechar{⇒}{\ensuremath{\Rightarrow}}
\newunicodechar{⟹}{\ensuremath{\Longrightarrow}}
\newunicodechar{∘}{\ensuremath{\circ}}
\newunicodechar{∪}{\ensuremath{\cup}}
\newunicodechar{∩}{\ensuremath{\cap}}
\newunicodechar{≡}{\ensuremath{\equiv}}
\newunicodechar{⊂}{\ensuremath{\subset}}
\newunicodechar{⊥}{\ensuremath{\perp}}
\newunicodechar{∃}{\ensuremath{\exists}}
\newunicodechar{∀}{\ensuremath{\forall}}
\newunicodechar{∛}{\ensuremath{\sqrt[3]{\,}}}
\newunicodechar{∜}{\ensuremath{\sqrt[4]{\,}}}
\newunicodechar{⃗}{\ensuremath{{}^{\to}}}
\newunicodechar{ⁿ}{\ifmmode^{n}\else\textsuperscript{\ensuremath{n}}\fi}
\newunicodechar{ˣ}{\ifmmode^{x}\else\textsuperscript{\ensuremath{x}}\fi}
\newunicodechar{ᵇ}{\ifmmode^{b}\else\textsuperscript{\ensuremath{b}}\fi}
\newunicodechar{ʳ}{\ifmmode^{r}\else\textsuperscript{\ensuremath{r}}\fi}
\newunicodechar{ᵘ}{\ifmmode^{u}\else\textsuperscript{\ensuremath{u}}\fi}
\newunicodechar{ʸ}{\ifmmode^{y}\else\textsuperscript{\ensuremath{y}}\fi}
\newunicodechar{ᵃ}{\ifmmode^{a}\else\textsuperscript{\ensuremath{a}}\fi}
\newunicodechar{ᵉ}{\ifmmode^{e}\else\textsuperscript{\ensuremath{e}}\fi}
\newunicodechar{ᵏ}{\ifmmode^{k}\else\textsuperscript{\ensuremath{k}}\fi}
\newunicodechar{ᵖ}{\ifmmode^{p}\else\textsuperscript{\ensuremath{p}}\fi}
\newunicodechar{ᴺ}{\ifmmode^{N}\else\textsuperscript{\ensuremath{N}}\fi}
\newunicodechar{ᶜ}{\ifmmode^{c}\else\textsuperscript{\ensuremath{c}}\fi}
\newunicodechar{ʰ}{\ifmmode^{h}\else\textsuperscript{\ensuremath{h}}\fi}
\newunicodechar{ᵠ}{\ifmmode^{q}\else\textsuperscript{\ensuremath{q}}\fi}
\newunicodechar{ₙ}{\ifmmode_{n}\else\textsubscript{\ensuremath{n}}\fi}
\newunicodechar{ₐ}{\ifmmode_{a}\else\textsubscript{\ensuremath{a}}\fi}
\newunicodechar{ᵢ}{\ifmmode_{i}\else\textsubscript{\ensuremath{i}}\fi}
\newunicodechar{ᵣ}{\ifmmode_{r}\else\textsubscript{\ensuremath{r}}\fi}
\newunicodechar{ₖ}{\ifmmode_{k}\else\textsubscript{\ensuremath{k}}\fi}
\newunicodechar{ᵦ}{\ifmmode_{\beta}\else\textsubscript{\ensuremath{\beta}}\fi}
\newunicodechar{ₓ}{\ifmmode_{x}\else\textsubscript{\ensuremath{x}}\fi}
\newfontfamily\popdisplay{ArchivoBlack-Regular.ttf}[Path=../../../fonts/]
\newfontfamily\poplogo{SpaceGrotesk-Bold.ttf}[Path=../../../fonts/]
\newfontfamily\popsemi{PlusJakartaSans-SemiBold.ttf}[Path=../../../fonts/]
\newfontfamily\popxbold{PlusJakartaSans-ExtraBold.ttf}[Path=../../../fonts/]
\newfontfamily\popxboldls{PlusJakartaSans-ExtraBold.ttf}[Path=../../../fonts/,LetterSpace=3.0]

\setlist[enumerate]{leftmargin=1.7em,itemsep=5pt,topsep=4pt}
\setlist[itemize]{leftmargin=1.7em,itemsep=4pt,topsep=4pt,label={\color{poporange}\textbullet}}
\setlength{\parindent}{0pt}
\setlength{\parskip}{5pt}
\renewcommand{\arraystretch}{1.25}

% pgfplots : style Pop par défaut
\pgfplotsset{
  every axis/.append style={axis line style={popink,line width=1pt},tick style={popink},
    grid style={popline,line width=0.5pt},label style={font=\small},tick label style={font=\footnotesize}},
  popcurve/.style={poporange,line width=1.8pt,no markers,smooth},
}
% Même style disponible dans un \draw TikZ simple (hors pgfplots)
\tikzset{popcurve/.style={poporange,line width=1.8pt}}
\tikzset{popgrid/.style={popline,very thin}}
\colorlet{popcurve}{poporange} % le nom du style est parfois utilisé comme couleur

% ============================================================
% Pill / squiggle
% ============================================================
\newcommand{\poppill}[3]{%
  \tikz[baseline=-0.55ex]\node[draw=popink,line width=1.2pt,rounded corners=2.6pt,%
    fill=#2,inner xsep=6.5pt,inner ysep=3pt]{%
    \popxboldls\fontsize{7}{7}\selectfont\color{#3}\MakeUppercase{#1}};%
}
\newcommand{\popsquiggle}[1]{%
  \tikz[baseline=-0.4ex]{
    \draw[#1,line width=2.6pt,line cap=round]
      plot [smooth] coordinates {(0,0.09) (0.34,0.01) (0.68,0.21) (1.02,0.01) (1.36,0.10)};
  }%
}
% Mot-clé surligné (marqueur orange sous le texte)
\newcommand{\cle}[1]{{\popxbold\color{popdark}#1}}

% ============================================================
% En-tête / pied
% ============================================================
\pagestyle{fancy}
\fancyhf{}
\renewcommand{\headrulewidth}{0pt}
\renewcommand{\footrulewidth}{0pt}
\lhead{\raisebox{-0.15ex}{\poplogo\fontsize{13}{13}\selectfont\color{popink}OnBuch\color{poporange}+}}
\rhead{\poppill{\DOCMATIERE\ \textbullet\ \DOCNIVEAU\ \textbullet\ Leçon \DOCLECON}{white}{popink}}
\lfoot{\popxbold\fontsize{7.6}{7.6}\selectfont\color{popmuted}\MakeUppercase{Cours signé OnBuch+}\ \textbullet\ {\popsemi\DOCTITRE}}
\rfoot{\tikz[baseline=-0.6ex]\node[rounded corners=8.5pt,fill=popink,minimum height=17pt,minimum width=17pt,inner xsep=5pt,inner sep=0]{\popxbold\fontsize{8}{8}\selectfont\color{popcream}\thepage\,/\,\pageref*{LastPage}};}

% ============================================================
% Titres
% ============================================================
\newcommand{\chaptitre}[2]{%
  \begin{tcolorbox}[enhanced,colback=poporange,colframe=popink,boxrule=2.6pt,arc=11pt,
      drop shadow=popink,left=15pt,right=15pt,top=12pt,bottom=11pt,before skip=3pt,after skip=18pt]
    \poppill{#2}{popink}{popcream}\par\vspace{9pt}
    {\raggedright\hyphenpenalty=10000\exhyphenpenalty=10000\popdisplay\fontsize{22}{24}\selectfont\color{popcream}\MakeUppercase{#1}\par}
  \end{tcolorbox}
}
% Section numérotée : pastille orange avec le numéro + titre Archivo
\newcounter{popsec}
\newcommand{\coursec}[1]{%
  \stepcounter{popsec}\setcounter{popsub}{0}%
  \par\vspace{16pt}\noindent
  \begin{minipage}{\linewidth}
  \tikz[baseline=-0.7ex]\node[draw=popink,line width=1.6pt,fill=poporange,rounded corners=4pt,minimum size=22pt,inner sep=0]{\popdisplay\fontsize{12}{12}\selectfont\color{popcream}\thepopsec};%
  \hspace{8pt}{\popdisplay\fontsize{15}{17}\selectfont\color{popink}\MakeUppercase{#1}}\par
  \vspace{4pt}\noindent{\color{poporange}\rule{36pt}{3.6pt}}
  \end{minipage}\par\nopagebreak\vspace{8pt}
}
\newcounter{popsub}
\newcommand{\courssub}[1]{%
  \stepcounter{popsub}%
  \par\vspace{9pt}\noindent{\popxbold\fontsize{11.5}{13}\selectfont\color{popink}{\color{poporange}\thepopsec.\thepopsub}\hspace{6pt}#1}\par\nopagebreak\vspace{4pt}
}

% ============================================================
% Boîtes pédagogiques — même recette : fond blanc, bordure épaisse colorée,
% ombre dure encre, badge plein. Titre optionnel [#1].
% ============================================================
\tcbset{popbase/.style={breakable,enhanced,colback=white,boxrule=2.3pt,arc=9pt,
  shadow={3pt}{-3pt}{0pt}{popink},left=14pt,right=14pt,top=11pt,bottom=11pt,before skip=11pt,after skip=15pt,
  fontupper=\popsemi\fontsize{10.6}{14.5}\selectfont\color{popink},
  fontlower=\popsemi\fontsize{10.6}{14.5}\selectfont\color{popink}}}
% Coche dessinée (le glyphe ✓ n'existe pas dans les polices du document)
\newcommand{\popcheck}[1]{\tikz[baseline=-0.5ex]\draw[#1,line width=1.6pt,line cap=round,line join=round](-0.13,0.0)--(-0.04,-0.09)--(0.13,0.1);}
\providecommand{\coche}{\popcheck{popgreen}}
\providecommand{\resultat}[1]{\par\smallskip\noindent\fcolorbox{popgreen}{popgreenL}{\parbox{\dimexpr\linewidth-2\fboxsep-2\fboxrule\relax}{\textbf{Résultat.} #1}}\par\smallskip}
\newcommand{\popboxhead}[3]{\poppill{#1}{#2}{white}\ifx\relax#3\relax\else\hspace{6pt}{\popxbold\fontsize{10.6}{12}\selectfont\color{popink}#3}\fi\par\vspace{7pt}}

\newtcolorbox{aretenir}[1][]{popbase,colframe=popgreen,before upper={\popboxhead{À retenir}{popgreen}{#1}}}
\newtcolorbox{exemplebox}[1][]{popbase,colframe=poporange,before upper={\popboxhead{Exemple}{poporange}{#1}}}
\newtcolorbox{definition}[1][]{popbase,colframe=popblue,before upper={\popboxhead{Définition}{popblue}{#1}}}
\newtcolorbox{propriete}[1][]{popbase,colframe=poppurple,before upper={\popboxhead{Propriété}{poppurple}{#1}}}
\newtcolorbox{methode}[1][]{popbase,colframe=popgold,before upper={\popboxhead{Méthode}{popgold}{#1}}}
\newtcolorbox{attention}[1][]{popbase,colframe=poppink,before upper={\popboxhead{Attention}{poppink}{#1}}}
\newtcolorbox{experience}[1][]{popbase,colframe=popblue,colback=popblueL,before upper={\popboxhead{Expérience}{popblue}{#1}}}
% « Le savais-tu ? » : carte violette pleine (contexte, culture, Cameroun)
\newtcolorbox{savaistu}[1][]{popbase,colback=poppurple,colframe=popink,
  fontupper=\popsemi\fontsize{10.4}{14.2}\selectfont\color{white},
  before upper={\poppill{Le savais-tu\,?}{white}{poppurple}\ifx\relax#1\relax\else\hspace{6pt}{\popxbold\color{white}#1}\fi\par\vspace{7pt}}}
% Exercice résolu : énoncé en haut, \tcblower puis la solution sur fond crème
\newcounter{popexo}\newcounter{popexr}
\newtcolorbox{exoresolu}[1][]{popbase,colframe=poporange,colbacklower=popcream,
  segmentation style={popink,line width=1.2pt,dashed},
  before upper={\stepcounter{popexr}\popboxhead{Exercice résolu \thepopexr}{poporange}{#1}},
  before lower={\poppill{Solution}{popink}{popcream}\par\vspace{6pt}}}
% Exercice (non résolu en place) — la correction est dans la partie Corrigés
\newtcolorbox{exercice}[1][]{popbase,colframe=popink,
  before upper={\stepcounter{popexo}\popboxhead{Exercice \thepopexo}{popink}{#1}}}
% Corrigé
\newtcolorbox{corrige}[1][]{popbase,colframe=popgreen,colback=popgreenL,
  before upper={\popboxhead{Corrigé}{popgreen}{#1}}}
% Objectifs / prérequis / bilan
\newtcolorbox{objectifs}{popbase,colframe=popink,colback=poporangeL,
  before upper={\popboxhead{Objectifs de la leçon}{popink}{}}}
\newtcolorbox{prerequis}{popbase,colframe=popink,colback=popblueL,
  before upper={\popboxhead{Prérequis}{popblue}{}}}
\newtcolorbox{bilan}{popbase,colframe=popink,colback=white,boxrule=2.8pt,
  before upper={\popboxhead{Fiche bilan}{poporange}{L'essentiel à connaître pour l'examen}}}
% Figure encadrée avec légende
\newtcolorbox{popfigure}[1][]{enhanced,breakable=false,colback=white,colframe=popline,boxrule=1.4pt,arc=8pt,
  left=10pt,right=10pt,top=10pt,bottom=8pt,before skip=10pt,after skip=12pt,halign=center,
  lower separated=false,halign lower=center,
  fontlower=\popsemi\fontsize{9}{11.5}\selectfont\color{popmuted},
  #1}
\newcommand{\legende}[1]{\par\vspace{6pt}{\popsemi\fontsize{9}{11.5}\selectfont\color{popmuted}#1}}

% ============================================================
% Signature OnBuch+
% ============================================================
\newcommand{\poplogoinline}[1]{{\poplogo\fontsize{#1}{#1}\selectfont\color{popink}OnBuch\color{poporange}+}}
\newcommand{\signatureonbuch}{%
  \begin{tcolorbox}[enhanced,colback=popink,colframe=popink,boxrule=2.2pt,arc=10pt,
      drop shadow=poporange,left=14pt,right=14pt,top=10pt,bottom=10pt,before skip=0pt,after skip=0pt]
    \tikz[baseline=-0.7ex]\node[circle,fill=popgreen,draw=popcream,line width=1.3pt,minimum size=22pt,inner sep=0]{\popcheck{white}};%
    \hspace{9pt}\begin{minipage}[c]{0.62\linewidth}
      {\popxbold\fontsize{12}{13}\selectfont\color{popcream}Signé\ \textbullet\ {\poplogo OnBuch\color{poporange}+}}\par\vspace{2pt}
      {\raggedright\popsemi\fontsize{8}{10.5}\selectfont\color{popline}\MakeUppercase{Équipe pédagogique OnBuch+}\\\MakeUppercase{Conforme au programme officiel MINESEC}\par}
    \end{minipage}\hfill
    {\poplogo\fontsize{22}{22}\selectfont\color{popcream}OnBuch\color{poporange}+}
  \end{tcolorbox}%
}

% ============================================================
% Page de garde
% #1 titre · #2 matière · #3 niveau · #4 module · #5 n° leçon · #6 durée ·
% #7 description / objectifs (texte ou itemize)
% ============================================================
\newcommand{\pagegarde}[7]{%
  \thispagestyle{empty}
  \newgeometry{a4paper,margin=1.7cm,top=1.6cm,bottom=1.4cm}%
  \begin{tikzpicture}[remember picture,overlay]
    \fill[poporange!18] ([xshift=-3.2cm,yshift=-3.6cm]current page.north east) circle (4.6cm);
    \fill[poppurple!14] ([xshift=2.4cm,yshift=7.6cm]current page.south west) circle (3.8cm);
  \end{tikzpicture}%
  {\poplogo\fontsize{30}{30}\selectfont\color{popink}OnBuch\color{poporange}+}\hfill
  \raisebox{0.6ex}{\poppill{Cours signé \textbullet\ Premium}{popink}{popcream}}\par
  \vspace{1pt}\popsquiggle{poppurple}\par
  \vspace{8pt}
  \begin{tcolorbox}[enhanced,colback=poporange,colframe=popink,boxrule=2.8pt,arc=13pt,
      drop shadow=popink,left=18pt,right=18pt,top=12pt,bottom=13pt,after skip=10pt]
    \poppill{#2 \textbullet\ #3}{popink}{popcream}\hspace{5pt}\poppill{#4}{white}{popink}\par\vspace{8pt}
    {\popdisplay\fontsize{14}{14}\selectfont\color{popink}LEÇON #5}\par\vspace{4pt}
    {\raggedright\hyphenpenalty=10000\exhyphenpenalty=10000\popdisplay\fontsize{25}{27}\selectfont\color{popcream}\MakeUppercase{#1}\par}
    \vspace{8pt}
    {\popxbold\fontsize{9.5}{9.5}\selectfont\color{popink}\MakeUppercase{Durée conseillée : #6}}
  \end{tcolorbox}
  \begin{tcolorbox}[enhanced,colback=white,colframe=popink,boxrule=2.2pt,arc=10pt,
      drop shadow=popink,left=16pt,right=16pt,top=10pt,bottom=10pt,after skip=8pt]
    \poppill{Dans cette leçon}{poppurple}{white}\par\vspace{5pt}
    \popsemi\fontsize{9.3}{12.6}\selectfont\color{popink}#7
  \end{tcolorbox}
  \noindent\begin{minipage}[t]{0.487\linewidth}
    \begin{tcolorbox}[enhanced,colback=popgreen,colframe=popink,boxrule=2.2pt,arc=10pt,
        drop shadow=popink,left=11pt,right=11pt,top=10pt,bottom=10pt,height=3.1cm,valign=center]
      \tcbox[colback=white,colframe=popink,boxrule=1.3pt,arc=5pt,boxsep=0pt,nobeforeafter,
        left=3pt,right=3pt,top=3pt,bottom=3pt,valign=center]{\qrcode[height=1.9cm]{https://play.google.com/store/apps/details?id=com.luvvix.onbuch&hl=fr}}%
      \hspace{8pt}\begin{minipage}[c]{0.5\linewidth}
        {\popdisplay\fontsize{11}{12.5}\selectfont\color{white}\MakeUppercase{Télécharge OnBuch}}\par\vspace{4pt}
        {\raggedright\popsemi\fontsize{8.4}{11}\selectfont\color{white}Gratuit \textbullet\ Google Play\\Tuteur IA, annales, concours, résultats\par}
      \end{minipage}
    \end{tcolorbox}
  \end{minipage}\hfill
  \begin{minipage}[t]{0.487\linewidth}
    \begin{tcolorbox}[enhanced,colback=poppurple,colframe=popink,boxrule=2.2pt,arc=10pt,
        drop shadow=popink,left=11pt,right=11pt,top=10pt,bottom=10pt,height=3.1cm,valign=center]
      \tikz[baseline=-0.7ex]\node[circle,fill=white,draw=popink,line width=1.3pt,minimum size=1.6cm,inner sep=0]{\popdisplay\fontsize{24}{24}\selectfont\color{poppurple}+};%
      \hspace{8pt}\begin{minipage}[c]{0.6\linewidth}
        {\popdisplay\fontsize{11}{12.5}\selectfont\color{white}\MakeUppercase{Passe à OnBuch+}}\par\vspace{4pt}
        {\raggedright\popsemi\fontsize{8.4}{11}\selectfont\color{white}Tous les cours signés, Tuteur IA illimité, fiches PDF — onglet OnBuch+ de l'app\par}
      \end{minipage}
    \end{tcolorbox}
  \end{minipage}
  \vspace{8pt}
  \popcontenu
  \vfill
  \signatureonbuch
  \restoregeometry
  \newpage
}

% Rangée « Ce cours contient » (page de garde)
\newcommand{\poptile}[3]{% #1 couleur #2 gros texte #3 légende
  \begin{tcolorbox}[enhanced,colback=#1,colframe=popink,boxrule=2pt,arc=9pt,shadow={2.5pt}{-2.5pt}{0pt}{popink},
      left=6pt,right=6pt,top=9pt,bottom=9pt,halign=center,valign=center,height=1.75cm,width=\linewidth,nobeforeafter]
    {\popdisplay\fontsize{12.5}{13}\selectfont\color{white}\MakeUppercase{#2}}\par\vspace{3pt}
    {\centering\popsemi\fontsize{7.8}{9.5}\selectfont\color{white}#3\par}
  \end{tcolorbox}}
\newcommand{\popcontenu}{%
  \par\noindent{\popxboldls\fontsize{7.5}{7.5}\selectfont\color{popmuted}CE COURS SIGNÉ CONTIENT}\par\vspace{7pt}
  \noindent\begin{minipage}[t]{0.235\linewidth}\poptile{popblue}{Cours}{complet, illustré, pas à pas}\end{minipage}\hfill
  \begin{minipage}[t]{0.235\linewidth}\poptile{poporange}{Méthodes}{et exercices résolus}\end{minipage}\hfill
  \begin{minipage}[t]{0.235\linewidth}\poptile{poppink}{Exercices}{type examen + corrigés}\end{minipage}\hfill
  \begin{minipage}[t]{0.235\linewidth}\poptile{popgreen}{Fiche bilan}{l'essentiel à retenir}\end{minipage}\par}

% Sommaire
\newtcolorbox{sommaire}{enhanced,colback=white,colframe=popink,boxrule=2.2pt,arc=10pt,
  drop shadow=popink,left=18pt,right=18pt,top=13pt,bottom=13pt,before skip=6pt,after skip=20pt,
  before upper={\poppill{Sommaire}{poppurple}{white}\par\vspace{9pt}\popsemi\fontsize{10.6}{14.5}\selectfont\color{popink}}}

% Signature de fin de cours — poussée tout en bas de la dernière page
\newcommand{\findecours}{%
  \par\vfill\nopagebreak
  \begin{center}\popsquiggle{poporange}\end{center}\vspace{4pt}
  \signatureonbuch
}
\AtBeginDocument{\def\llbracket{\mathopen{[\mkern-3.5mu[}}\def\rrbracket{\mathclose{]\mkern-3.5mu]}}} % intervalles d'entiers (glyphe ⟦ absent de la police)
% Glyphes absents de Fira Math : redéfinis à partir de glyphes disponibles
\AtBeginDocument{%
  \def\setminus{\mathbin{\backslash}}\let\smallsetminus\setminus
  \def\vdots{\mathord{\vbox{\baselineskip3.2pt\lineskiplimit0pt\kern4pt\hbox{.}\hbox{.}\hbox{.}}}}%
  \def\ddots{\mathinner{\mkern1mu\raise6.5pt\hbox{.}\mkern2mu\raise3.6pt\hbox{.}\mkern2mu\raise0.7pt\hbox{.}\mkern1mu}}%
  \def\triangle{\mathord{\scalebox{0.9}{$\symup{\Delta}$}}}%
  \def\nabla{\mathord{\raisebox{\depth}{\scalebox{1}[-1]{$\symup{\Delta}$}}}}%
  \def\checkmark{\popcheck{popdark}}%
  \def\star{\mathbin{\ast}}\def\bigstar{\mathord{\ast}}%
  \def\diamond{\mathbin{\scalebox{0.6}{$\Diamond$}}}%
  \def\bigcup{\mathop{\mathchoice{\raisebox{-0.3em}{\scalebox{1.7}{$\cup$}}}{\scalebox{1.25}{$\cup$}}{\cup}{\cup}}\displaylimits}%
  \def\bigcap{\mathop{\mathchoice{\raisebox{-0.3em}{\scalebox{1.7}{$\cap$}}}{\scalebox{1.25}{$\cap$}}{\cap}{\cap}}\displaylimits}%
  \def\lesseqqgtr{\mathrel{\lessgtr}}\def\bigoplus{\mathop{\oplus}\displaylimits}%
}
\providecommand{\ssi}{si et seulement si} % abréviation fréquente
\AtBeginDocument{\providecommand{\wideparen}{\overparen}} % arc de cercle

%% FIGFIX-BEGIN v5 — correctifs globaux des figures (docs/onbuch-generation/tools/figfix)
\usepackage{adjustbox}\usepackage{etoolbox}\tcbuselibrary{raster}
% toute figure TikZ/pgfplots plus large que la ligne est réduite à la largeur disponible
\BeforeBeginEnvironment{tikzpicture}{\begin{adjustbox}{max width=\linewidth}}
\AfterEndEnvironment{tikzpicture}{\end{adjustbox}}
% légendes pgfplots lisibles par-dessus les courbes
\pgfplotsset{every axis/.append style={legend style={fill=white,fill opacity=0.92,text opacity=1,draw=black!15,inner sep=3pt}}}
% étiquettes d'arêtes (syntaxe edge["..."]) avec fond blanc
\tikzset{every edge quotes/.append style={fill=white,inner sep=1.2pt}}
%% FIGFIX-END
\begin{document}

\pagegarde{Énergie cinétique et théorème de l'énergie cinétique}{Physique}{1ère C}{Module 2 : Mouvements et interactions mécaniques}{N04}{4 h}{Cette leçon définit l'énergie cinétique d'un solide en translation et en rotation, puis établit le théorème de l'énergie cinétique. Elle fournit une méthode énergétique simple et puissante pour résoudre des problèmes de mécanique sans passer par l'étude détaillée du mouvement.
\vspace{4pt}
\begin{multicols}{2}\small
\begin{itemize}
\item À la fin de cette leçon, je sais définir l'énergie cinétique et citer son unité dans le Système international
\item À la fin de cette leçon, je sais exprimer et calculer l'énergie cinétique d'un solide en translation
\item À la fin de cette leçon, je sais exprimer et calculer l'énergie cinétique d'un solide en rotation autour d'un axe fixe
\item À la fin de cette leçon, je sais exprimer l'énergie cinétique d'un solide en mouvement combiné translation-rotation (roulement)
\item À la fin de cette leçon, je sais énoncer le théorème de l'énergie cinétique et préciser son contexte d'utilisation
\item À la fin de cette leçon, je sais appliquer le théorème de l'énergie cinétique dans des cas simples (plan horizontal, plan incliné, chute)
\end{itemize}\end{multicols}}

\begin{sommaire}
\begin{multicols}{2}
\begin{enumerate}
\item Notion d'énergie et énergie cinétique de translation
\item Énergie cinétique d'un solide en rotation
\item Solide en mouvement combiné : translation et rotation
\item Théorème de l'énergie cinétique
\item Applications directes du théorème
\item Méthode de résolution et synthèse
\item Activité d'intégration
\item Méthodes et exercices résolus
\item Exercices
\item Corrigés des exercices
\item Fiche bilan
\end{enumerate}
\end{multicols}
\end{sommaire}

\begin{objectifs}
\begin{itemize}
\item À la fin de cette leçon, je sais définir l'énergie cinétique et citer son unité dans le Système international
\item À la fin de cette leçon, je sais exprimer et calculer l'énergie cinétique d'un solide en translation
\item À la fin de cette leçon, je sais exprimer et calculer l'énergie cinétique d'un solide en rotation autour d'un axe fixe
\item À la fin de cette leçon, je sais exprimer l'énergie cinétique d'un solide en mouvement combiné translation-rotation (roulement)
\item À la fin de cette leçon, je sais énoncer le théorème de l'énergie cinétique et préciser son contexte d'utilisation
\item À la fin de cette leçon, je sais appliquer le théorème de l'énergie cinétique dans des cas simples (plan horizontal, plan incliné, chute)
\item À la fin de cette leçon, je sais résoudre un problème de mécanique en choisissant judicieusement la méthode énergétique
\item À la fin de cette leçon, je sais exploiter le théorème de l'énergie cinétique pour estimer des grandeurs liées à la mobilité (distance de freinage, vitesse, coût énergétique)
\end{itemize}
\end{objectifs}

\begin{prerequis}
\begin{itemize}
\item Travail d'une force constante : W\_AB(F) = F·AB·cos(α), travail du poids, travail moteur et résistant
\item Notion de force, bilan des forces et deuxième loi de Newton (relation fondamentale de la dynamique)
\item Grandeurs cinématiques : vitesse, accélération, mouvement rectiligne uniformément varié
\item Moment d'inertie J\_Δ d'un solide par rapport à un axe et vitesse angulaire ω (leçon sur la rotation)
\item Unités du Système international et conversions (km/h vers m/s)
\end{itemize}
\end{prerequis}

\newpage

\coursec{Notion d'énergie et énergie cinétique de translation}

Sur l'axe lourd Yaoundé--Douala, un chauffeur de bus de 70 places aperçoit soudain un camion arrêté en travers de la chaussée. Il écrase la pédale de frein. Les pneus crissent, les passagers sont projetés vers l'avant, et le bus s'immobilise à quelques mètres de l'obstacle. Un constat s'impose à tout usager de la route : plus le bus est chargé et rapide, plus la distance de freinage est grande. Mais pourquoi, précisément, \emph{doubler} la vitesse \emph{quadruple}-t-il la distance d'arrêt ? Comment évaluer l'énergie qu'un véhicule doit dissiper pour s'immobiliser, ou celle qu'un pousseur de brouette doit fournir au marché Mokolo pour lancer sa charge ?

Cette leçon introduit une grandeur fondamentale, l'\cle{énergie cinétique}, puis un théorème puissant qui relie le travail des forces à la variation de vitesse : le \cle{théorème de l'énergie cinétique}. C'est l'outil indispensable pour analyser le coût énergétique de la mobilité, du piéton au camion-citerne.

\courssub{L'intuition de l'énergie : un corps en mouvement peut produire un effet}

\textbf{Observations de la vie quotidienne.}\par

Réfléchis à quelques situations que tu connais bien :
\begin{itemize}
\item une mototaxi lancée à pleine vitesse qui heurte une barrière la \emph{déforme} ;
\item le vent qui fait tourner les pales d'une éolienne les \emph{déplace} ;
\item un clou frappé par un marteau en mouvement \emph{s'enfonce} dans le bois ;
\item les freins d'un camion qui descend la côte de Mbanga \emph{chauffent} au point de fumer.
\end{itemize}

Dans chaque cas, un corps \emph{en mouvement} produit un effet : il déforme, déplace, enfonce ou chauffe un autre objet. Or, en physique, produire un tel effet, c'est fournir un \cle{travail} mécanique. Un corps capable de fournir un travail possède de l'\cle{énergie}.

\begin{definition}[Énergie (approche intuitive)]
L'\cle{énergie} d'un système est la grandeur qui mesure sa capacité à produire un travail, c'est-à-dire à déplacer, déformer ou modifier l'état d'autres systèmes. L'énergie se mesure en \cle{joules} (symbole J).
\end{definition}

Un corps \emph{immobile} posé sur la route ne déforme rien. C'est bien le \emph{mouvement} qui confère au marteau, à la moto ou au bus cette capacité d'action. L'énergie qu'un corps possède du fait de son mouvement porte un nom : l'\cle{énergie cinétique} (du grec \emph{kinêsis}, mouvement).

\textbf{De quels paramètres dépend-elle ?}\par

L'expérience quotidienne suggère deux facteurs :
\begin{itemize}
\item la \cle{vitesse} : à masse égale, une balle de football tirée à \SI{30}{m/s} fait bien plus mal qu'à \SI{5}{m/s} ;
\item la \cle{masse} : à vitesse égale, le choc d'un camion de \SI{30}{t} est incomparablement plus destructeur que celui d'une moto de \SI{200}{kg}.
\end{itemize}
Reste à trouver la \emph{loi mathématique exacte} qui relie l'énergie cinétique à la masse et à la vitesse. C'est l'objet du paragraphe suivant.

\courssub{Établissement de l'expression de l'énergie cinétique de translation}

\textbf{Mise en place du modèle.}\par

Considérons un solide de masse $m$, assimilé à son centre d'inertie $G$, en \cle{translation} rectiligne sur une route horizontale. Un moteur (ou un pousseur) exerce sur lui une force constante $\vec{F}$, parallèle au déplacement. Le solide part du repos ($v_0 = 0$) et on néglige les frottements. On cherche le travail $W$ que la force doit fournir pour amener le solide à la vitesse $v$ : ce travail est précisément l'énergie que le solide aura emmagasinée sous forme de mouvement.

\textbf{Application de la deuxième loi de Newton.}\par

La seule force horizontale est $\vec{F}$ (le poids $\vec{P}$ et la réaction $\vec{R}$ de la route, verticaux, se compensent). La deuxième loi de Newton projetée sur l'axe du mouvement donne :
\[ F = m\,a \quad\Longrightarrow\quad a = \dfrac{F}{m} = \text{constante}. \]
Le mouvement est donc rectiligne uniformément accéléré. Entre le point de départ (vitesse nulle) et le point où la vitesse vaut $v$, la relation indépendante du temps s'écrit :
\[ v^2 - v_0^2 = 2\,a\,d \quad\Longrightarrow\quad v^2 = 2\,a\,d, \]
où $d$ est la distance parcourue. Le travail de la force constante $\vec{F}$ sur cette distance est :
\[ W = F\,d = (m\,a)\,d = m\,(a\,d) = m\,\dfrac{v^2}{2}. \]
Le travail fourni au solide pour le faire passer du repos à la vitesse $v$ vaut donc exactement $\dfrac{1}{2}mv^2$. Par définition de l'énergie, c'est l'énergie que le solide possède désormais du fait de son mouvement.

\begin{definition}[Énergie cinétique d'un solide en translation]
L'\cle{énergie cinétique} d'un solide de masse $m$ en mouvement de translation à la vitesse $v$ est :
\[ \boxed{E_c = \dfrac{1}{2}\,m\,v^2} \]
avec $m$ en kilogrammes (kg), $v$ en mètres par seconde (m/s) et $E_c$ en joules (J). Dans un mouvement de translation, tous les points du solide ont la même vitesse, celle du centre d'inertie $G$.
\end{definition}

\begin{attention}[Convertis toujours la vitesse en m/s]
L'erreur la plus fréquente au Probatoire est d'injecter une vitesse en km/h directement dans la formule. La vitesse \textbf{doit} être exprimée en m/s. Règle de conversion :
\[ v\ (\text{m/s}) = \dfrac{v\ (\text{km/h})}{3,6}. \]
Exemple : \SI{72}{km/h} $= \dfrac{72}{3,6} =$ \SI{20}{m/s}. Une vitesse laissée en km/h fausse le résultat d'un facteur $3,6^2 \approx 13$ !
\end{attention}

\courssub{Unité d'énergie : le joule et les unités usuelles}

\begin{propriete}[Unité SI de l'énergie]
L'unité SI de l'énergie est le \cle{joule} (J). L'analyse dimensionnelle de $E_c = \frac{1}{2}mv^2$ donne :
\[ 1\ \text{J} = 1\ \text{kg}\cdot\left(\text{m/s}\right)^2 = 1\ \text{kg}\cdot\text{m}^2\cdot\text{s}^{-2}. \]
Un joule est l'énergie cinétique d'une masse de \SI{2}{kg} se déplaçant à \SI{1}{m/s} (ou le travail d'une force de \SI{1}{N} sur \SI{1}{m}).
\end{propriete}

Le joule est une unité petite à l'échelle des transports. On utilise couramment :
\begin{center}
\begin{tabularx}{\linewidth}{|l|X|l|}
\hline
\rowcolor{popblueL} \textbf{Unité} & \textbf{Équivalence} & \textbf{Domaine d'usage} \\
\hline
kilojoule (kJ) & $1\ \text{kJ} = 10^3\ \text{J}$ & énergies mécaniques courantes \\
\hline
mégajoule (MJ) & $1\ \text{MJ} = 10^6\ \text{J}$ & véhicules, carburants \\
\hline
watt-heure (Wh) & $1\ \text{Wh} = 3600\ \text{J}$ & batteries, petits appareils \\
\hline
kilowatt-heure (kWh) & $1\ \text{kWh} = 3,6\times 10^6\ \text{J}$ & factures ENEO, électroménager \\
\hline
\end{tabularx}
\end{center}

Retiens la conversion phare : $1\ \text{kWh} = 3,6\times 10^6\ \text{J}$. C'est l'énergie que facture ENEO ; nous verrons qu'un bus lancé à \SI{72}{km/h} transporte une énergie cinétique comparable à une fraction de kilowatt-heure.

\courssub{Propriétés essentielles de l'énergie cinétique}

\textbf{Une grandeur toujours positive.}\par

La masse $m$ est positive et $v^2$ est positif quel que soit le sens du mouvement : $E_c \geq 0$ toujours. L'énergie cinétique est une grandeur \emph{scalaire} : elle ne dépend pas de la \emph{direction} du vecteur vitesse. Un bus roulant vers Douala ou vers Yaoundé à la même vitesse possède la même énergie cinétique. Elle ne s'annule que si le solide est au repos.

\textbf{Une dépendance quadratique en vitesse.}\par

C'est le point capital de cette section. Comme $E_c$ est proportionnelle à $v^2$ :
\begin{itemize}
\item doubler la vitesse ($v \to 2v$) \emph{quadruple} l'énergie : $E_c \to 4E_c$ ;
\item tripler la vitesse ($v \to 3v$) la \emph{multiplie par 9} ;
\item diviser la vitesse par deux la divise par 4.
\end{itemize}

\begin{aretenir}[Le message sécurité routière]
L'énergie cinétique croît comme le \emph{carré} de la vitesse. C'est pourquoi les accidents à \SI{100}{km/h} sont infiniment plus graves qu'à \SI{50}{km/h} : l'énergie à dissiper lors du choc est \textbf{quatre fois} plus grande, et non deux fois.
\end{aretenir}

La courbe ci-dessous montre $E_c$ en fonction de $v$ pour deux véhicules de l'axe Yaoundé--Douala : une moto de \SI{200}{kg} et un bus de \SI{8000}{kg}.

\begin{popfigure}
\begin{center}
\begin{tikzpicture}
\begin{axis}[
width=0.85\linewidth, height=6.5cm,
xlabel={$v$ (m/s)}, ylabel={$E_c$ (J)},
xmin=0, xmax=30, ymin=0, ymax=3.6e6,
xtick={0,5,10,15,20,25,30},
ytick={0,1e6,2e6,3e6},
yticklabels={$0$,$1\times10^6$,$2\times10^6$,$3\times10^6$},
legend style={at={(0.03,0.97)},anchor=north west,font=\small},
grid=both, grid style={popline!40},
]
\addplot[popcurve,domain=0:30,samples=100,color=popblue,very thick]{0.5*200*x^2};
\addlegendentry{moto : $m = 200$ kg}
\addplot[popcurve,domain=0:30,samples=100,color=poporange,very thick]{0.5*8000*x^2};
\addlegendentry{bus : $m = 8000$ kg}
\end{axis}
\end{tikzpicture}
\end{center}
\legende{Énergie cinétique en fonction de la vitesse : croissance parabolique (en $v^2$). À masse égale, la courbe monte de plus en plus vite ; à vitesse égale, le bus de 8 t possède 40 fois l'énergie de la moto de 200 kg.}
\end{popfigure}

\textbf{Lien avec la distance de freinage.}\par

Pour s'arrêter, un véhicule doit \emph{dissiper} toute son énergie cinétique, essentiellement par le travail résistant de la force de freinage $\vec{f}$ (frottements des plaquettes et des pneus). Si cette force est supposée constante, son travail sur la distance d'arrêt $d$ vaut $W = -f\,d$, et l'arrêt complet exige :
\[ f\,d = \dfrac{1}{2}mv^2 \quad\Longrightarrow\quad d = \dfrac{m\,v^2}{2f}. \]
La distance de freinage est donc, elle aussi, proportionnelle à $v^2$ : \textbf{doubler la vitesse quadruple la distance d'arrêt}. C'est la réponse quantitative à la question posée en introduction. (Nous démontrerons ce résultat proprement avec le théorème de l'énergie cinétique dans la section 4.)

\begin{popfigure}
\begin{center}
\begin{tikzpicture}[scale=0.9, >=Stealth]
% Route 1
\draw[fill=popcream,draw=popline] (0,0) rectangle (13,1.2);
\draw[dashed,popmuted] (0.3,0.6) -- (12.7,0.6);
% Voiture à 50 km/h
\draw[fill=popblue,draw=popdark] (0.6,1.2) rectangle (1.8,1.9);
\draw[fill=popblueL,draw=popdark] (0.85,1.9) rectangle (1.55,2.3);
\draw[fill=popdark] (0.9,1.2) circle (0.18);
\draw[fill=popdark] (1.5,1.2) circle (0.18);
\node[font=\small,popdark] at (1.2,2.6) {$v = 50$ km/h};
% Distance d1
\draw[<->,popblue,thick] (1.9,0.35) -- (5.4,0.35) node[midway,below,font=\small,popblue]{$d_1$};
\draw[popblue,thick] (5.4,0.15) -- (5.4,1.2);
\node[font=\small,popblue,anchor=west] at (5.5,0.9) {arrêt};
% Route 2
\draw[fill=popcream,draw=popline] (0,3.2) rectangle (13,4.4);
\draw[dashed,popmuted] (0.3,3.8) -- (12.7,3.8);
% Voiture à 100 km/h
\draw[fill=poporange,draw=popdark] (0.6,4.4) rectangle (1.8,5.1);
\draw[fill=poporangeL,draw=popdark] (0.85,5.1) rectangle (1.55,5.5);
\draw[fill=popdark] (0.9,4.4) circle (0.18);
\draw[fill=popdark] (1.5,4.4) circle (0.18);
\node[font=\small,popdark] at (1.2,5.8) {$v = 100$ km/h};
\draw[<->,poporange,thick] (1.9,2.9) -- (12.4,2.9) node[midway,below,font=\small,poporange]{$d_2 \approx 4\,d_1$};
\draw[poporange,thick] (12.4,2.7) -- (12.4,3.2);
\node[font=\small,poporange,anchor=east] at (12.3,2.5) {arrêt};
\end{tikzpicture}
\end{center}
\legende{À force de freinage égale, la distance d'arrêt est proportionnelle à $v^2$ : en passant de 50 à 100 km/h, elle est multipliée par 4 (environ 15 m contre 60 m sur route sèche, hors distance de réaction).}
\end{popfigure}

\courssub{Exemples chiffrés : du piéton au camion-citerne}

\begin{exemplebox}[Énergie cinétique d'un bus sur l'axe Yaoundé--Douala]
Un bus de 70 places, de masse totale $m = \SI{8}{t} = \SI{8000}{kg}$, roule à $v = \SI{72}{km/h}$. Calculons son énergie cinétique.
\begin{align*}
v &= \dfrac{72}{3,6} = \SI{20}{m/s} \\
E_c &= \dfrac{1}{2}mv^2 = \dfrac{1}{2}\times 8000\times 20^2 = 4000\times 400 = \SI{1,6e6}{J}.
\end{align*}
Soit $E_c = \SI{1,6}{MJ}$, l'équivalent d'environ \SI{0,44}{kWh} : c'est l'énergie qu'il faudra entièrement dissiper dans les freins pour immobiliser le bus.
\end{exemplebox}

Comparons maintenant plusieurs usagers de la route, tous à des vitesses réalistes :

\begin{center}
\begin{tabularx}{\linewidth}{|l|X|X|X|}
\hline
\rowcolor{popblueL} \textbf{Mobile} & \textbf{Masse} & \textbf{Vitesse} & \textbf{$E_c = \frac{1}{2}mv^2$} \\
\hline
Piéton (quartier Mokolo) & \SI{60}{kg} & \SI{5}{km/h} = \SI{1,4}{m/s} & $E_c \approx \SI{58}{J}$ \\
\hline
Mototaxi + passager & \SI{200}{kg} & \SI{50}{km/h} = \SI{13,9}{m/s} & $E_c \approx \SI{1,9e4}{J}$ \\
\hline
Bus 70 places & \SI{8000}{kg} & \SI{72}{km/h} = \SI{20}{m/s} & $E_c = \SI{1,6e6}{J}$ \\
\hline
Camion-citerne & \SI{30000}{kg} & \SI{90}{km/h} = \SI{25}{m/s} & $E_c \approx \SI{9,4e6}{J}$ \\
\hline
\end{tabularx}
\end{center}

Le contraste est saisissant : le camion-citerne emporte plus de \emph{cent soixante mille fois} l'énergie cinétique du piéton. C'est toute cette énergie qui se transforme en déformations et en chaleur lors d'un choc.

\begin{exoresolu}[La mototaxi et le marchand de brouettes]
\textbf{Énoncé.} (a) Une mototaxi de masse totale $m = \SI{200}{kg}$ roule à $v_1 = \SI{50}{km/h}$. Calcule son énergie cinétique. (b) Quelle devient cette énergie si la vitesse double, $v_2 = \SI{100}{km/h}$ ? Conclus. (c) Un pousseur de brouette au marché Mokolo lance sa charge ($m' = \SI{120}{kg}$) à $v' = \SI{2,0}{m/s}$. Quel travail minimal a-t-il dû fournir ?
\tcblower
\textbf{Solution.} (a) Conversion : $v_1 = \dfrac{50}{3,6} \approx \SI{13,9}{m/s}$. Alors :
\[ E_{c1} = \dfrac{1}{2}\times 200\times (13,9)^2 \approx \SI{1,9e4}{J}. \]
(b) $v_2 = \dfrac{100}{3,6} \approx \SI{27,8}{m/s}$, soit $E_{c2} = \dfrac{1}{2}\times 200\times (27,8)^2 \approx \SI{7,7e4}{J}$. On vérifie que $E_{c2} = 4\,E_{c1}$ : doubler la vitesse quadruple l'énergie cinétique. (c) Le travail minimal fourni est l'énergie cinétique acquise (frottements négligés) :
\[ W = \dfrac{1}{2}m'v'^2 = \dfrac{1}{2}\times 120\times (2,0)^2 = \SI{240}{J}. \]
\end{exoresolu}

\begin{savaistu}[Une chute de 13 étages]
À \SI{100}{km/h}, l'énergie cinétique d'une voiture équivaut à celle qu'elle acquerrait en chutant du sommet d'un immeuble d'environ 13 étages (soit une hauteur voisine de \SI{39}{m}, car $v^2 = 2gh$ donne $h = \dfrac{v^2}{2g} \approx \dfrac{(27,8)^2}{2\times 9,8} \approx \SI{39}{m}$). Percuter un mur à \SI{100}{km/h}, c'est donc subir la même violence qu'une chute du dernier étage d'une tour du centre-ville de Yaoundé. Voilà pourquoi le port de la ceinture et le respect des limitations de vitesse sauvent des vies : la physique ne négocie pas.
\end{savaistu}

\begin{attention}[Bilan des pièges de cette section]
\begin{itemize}
\item Vitesse \textbf{toujours} en m/s (diviser les km/h par 3,6) avant tout calcul.
\item Ne pas oublier le \textbf{carré} sur la vitesse : $E_c = \frac{1}{2}mv^2$, pas $\frac{1}{2}mv$.
\item La masse s'exprime en \textbf{kg} : \SI{8}{t} = \SI{8000}{kg}.
\item $E_c$ est toujours positive ; elle ne dépend pas du sens du mouvement.
\item Ne pas confondre joule (énergie) et watt (puissance) : $1\ \text{W} = 1\ \text{J/s}$.
\end{itemize}
\end{attention}

Nous savons désormais calculer l'énergie cinétique d'un solide en translation. Mais les roues du bus, elles, \emph{tournent} : la section suivante montrera qu'un solide en rotation possède lui aussi une énergie cinétique, de forme analogue, avant que nous n'énoncions le théorème qui relie ces énergies au travail des forces.

\coursec{Énergie cinétique d'un solide en rotation}

Dans la section précédente, nous avons vu qu'un solide en \emph{translation} possède une énergie cinétique $E_c = \dfrac{1}{2}mv^2$, où $v$ est la vitesse commune de tous ses points. Mais observe un instant le tambour d'une machine à laver en plein essorage, ou la roue d'une mototaxi qui démarre sur le boulevard : ici, le solide ne se déplace pas globalement, il \emph{tourne} autour d'un axe fixe. Ses points n'ont plus tous la même vitesse : un point proche de l'axe se déplace lentement, un point au bord file très vite. Comment, dès lors, exprimer l'énergie cinétique d'un tel solide ? C'est toute l'ambition de cette section.

\courssub{Observation : chaque point d'un solide en rotation a sa propre vitesse}

Considérons un solide indéformable en rotation autour d'un axe fixe $\Delta$, à la vitesse angulaire $\omega$ (en $\mathrm{rad\cdot s^{-1}}$). Tous les points du solide décrivent des cercles centrés sur l'axe, avec la \emph{même} vitesse angulaire $\omega$, mais des \emph{rayons différents}.

Pour un point $M_i$ de masse $m_i$ situé à la distance $r_i$ de l'axe $\Delta$, la trajectoire est un cercle de rayon $r_i$, parcouru à la vitesse linéaire :
\[ v_i = r_i\,\omega \]

Cette relation est fondamentale : la vitesse linéaire d'un point \textbf{croît proportionnellement à sa distance à l'axe}. Un point situé deux fois plus loin de l'axe va deux fois plus vite. C'est pourquoi, sur une roue de moto, la bande de roulement s'use et chauffe bien plus que la zone proche du moyeu.

\begin{popfigure}
\begin{tikzpicture}[scale=1.05]
  % Disque
  \fill[popblueL] (0,0) circle (2.6);
  \draw[popblue, thick] (0,0) circle (2.6);
  % Axe
  \fill[popdark] (0,0) circle (0.07);
  \node[below right, popdark] at (0.05,-0.05) {$\Delta$};
  % Flèche de rotation
  \draw[->, poporange, very thick] (140:3.0) arc (140:60:3.0);
  \node[poporange] at (100:3.35) {$\omega$};
  % Point 1 proche de l'axe
  \fill[popgreen] (0.8,0) circle (0.09);
  \draw[->, popgreen, very thick] (0.8,0) -- (0.8,0.75);
  \node[popgreen, right] at (0.85,0.55) {$\vec{v}_1$};
  \node[popgreen, below] at (0.8,-0.12) {$M_1$};
  \draw[popmuted, dashed] (0,0) -- (0.8,0) node[midway, below, popdark] {\footnotesize $r_1$};
  % Point 2 intermédiaire
  \fill[poppurple] (30:1.7) circle (0.09);
  \draw[->, poppurple, very thick] (30:1.7) -- ++(120:1.5);
  \node[poppurple, above left] at ($(30:1.7)+(120:1.5)$) {$\vec{v}_2$};
  \node[poppurple, below right] at (30:1.62) {$M_2$};
  \draw[popmuted, dashed] (0,0) -- (30:1.7) node[midway, below left, popdark] {\footnotesize $r_2$};
  % Point 3 au bord
  \fill[popink] (-2.5,0) circle (0.09);
  \draw[->, popink, very thick] (-2.5,0) -- (-2.5,-2.2);
  \node[popink, below left] at (-2.5,-2.2) {$\vec{v}_3$};
  \node[popink, above] at (-2.5,0.12) {$M_3$};
  \draw[popmuted, dashed] (0,0) -- (-2.5,0) node[midway, above, popdark] {\footnotesize $r_3$};
\end{tikzpicture}
\legende{Disque en rotation autour de l'axe fixe $\Delta$. Les points $M_1$, $M_2$, $M_3$ ont la même vitesse angulaire $\omega$, mais des vitesses linéaires différentes : $v_i = r_i\omega$, d'autant plus grandes que le point est éloigné de l'axe. Les vecteurs vitesse sont tangents aux trajectoires circulaires.}
\end{popfigure}

\courssub{De la somme des énergies cinétiques à l'expression générale}

\textbf{Intuition.} Un solide en rotation n'est rien d'autre qu'un ensemble de points matériels en mouvement. L'énergie cinétique totale du solide est donc la \emph{somme} des énergies cinétiques de tous ses points. Toute la physique de cette section tient dans cette idée simple.

Découpons mentalement le solide en petits éléments de matière, assimilés à des points matériels $M_i$ de masse $m_i$, situés à la distance $r_i$ de l'axe $\Delta$. Chacun possède l'énergie cinétique $\dfrac{1}{2}m_i v_i^2$. L'énergie cinétique totale du solide s'écrit :
\[ E_c = \sum_i \frac{1}{2}m_i v_i^2 \]

Or tous les points tournent avec la même vitesse angulaire $\omega$, et $v_i = r_i\omega$. En reportant :
\begin{align*}
E_c &= \sum_i \frac{1}{2}m_i (r_i\omega)^2 \\
    &= \sum_i \frac{1}{2}m_i r_i^2 \omega^2 \\
    &= \frac{1}{2}\left(\sum_i m_i r_i^2\right)\omega^2
\end{align*}

La factorisation par $\omega^2$ est l'étape décisive : la quantité entre parenthèses ne dépend que du solide (sa masse et la répartition de cette masse autour de l'axe), et \emph{pas} de la vitesse de rotation. On lui donne un nom : le \cle{moment d'inertie} du solide par rapport à l'axe $\Delta$.

\begin{definition}[Moment d'inertie]
Le \cle{moment d'inertie} d'un solide par rapport à un axe $\Delta$ est la grandeur :
\[ J_\Delta = \sum_i m_i r_i^2 \]
où $m_i$ est la masse de chaque élément du solide et $r_i$ sa distance à l'axe $\Delta$. Il s'exprime en $\mathrm{kg\cdot m^2}$.
\end{definition}

\begin{definition}[Énergie cinétique d'un solide en rotation]
L'\cle{énergie cinétique} d'un solide de moment d'inertie $J_\Delta$ tournant à la vitesse angulaire $\omega$ autour d'un axe fixe $\Delta$ est :
\[ \boxed{E_c = \frac{1}{2}J_\Delta\,\omega^2} \]
avec $E_c$ en joules ($\mathrm{J}$), $J_\Delta$ en $\mathrm{kg\cdot m^2}$ et $\omega$ en $\mathrm{rad\cdot s^{-1}}$.
\end{definition}

Vérifions l'homogénéité de cette expression par analyse dimensionnelle :
\[ [J_\Delta] = \mathrm{kg\cdot m^2},\quad [\omega^2] = \mathrm{s^{-2}} \quad\Rightarrow\quad [J_\Delta\,\omega^2] = \mathrm{kg\cdot m^2\cdot s^{-2}} = \mathrm{J} \]
C'est bien une énergie.

\begin{propriete}[Signification physique du moment d'inertie]
Le moment d'inertie $J_\Delta$ dépend :
\begin{itemize}
  \item de la \textbf{masse totale} du solide : à géométrie identique, un solide deux fois plus massif a un moment d'inertie deux fois plus grand ;
  \item de la \textbf{répartition de la masse autour de l'axe} : plus la masse est éloignée de l'axe (grands $r_i$), plus $J_\Delta$ est grand, car les distances interviennent au \emph{carré}.
\end{itemize}
Le moment d'inertie mesure la \emph{résistance} du solide à être mis en rotation : il joue en rotation le rôle que joue la masse en translation. Son unité SI est le $\mathrm{kg\cdot m^2}$.
\end{propriete}

\begin{savaistu}[Pourquoi les patineurs tournent-ils plus vite bras repliés ?]
Quand un patineur replie les bras, il rapproche de la masse de son axe de rotation : son moment d'inertie $J_\Delta$ diminue. Comme la quantité $J_\Delta\,\omega$ (le moment cinétique, que tu étudieras plus tard) se conserve, sa vitesse angulaire $\omega$ augmente spectaculairement. Même principe pour le plongeur qui se groupe en boule pour effectuer ses rotations avant d'entrer dans l'eau.
\end{savaistu}

\courssub{Moments d'inertie des solides usuels}

Le calcul de $J_\Delta = \sum_i m_i r_i^2$ pour un solide continu relève de l'intégration, qui dépasse le programme de Première. Au Probatoire, les moments d'inertie des solides usuels sont \emph{donnés} dans l'énoncé ou rappelés ci-dessous ; tu dois savoir les \emph{utiliser} correctement.

\begin{aretenir}[Moments d'inertie à connaître]
Pour un solide de masse $m$, par rapport à son axe de symétrie de révolution :
\begin{itemize}
  \item \textbf{Cerceau} (ou anneau fin) de rayon $R$ : toute la masse est à la distance $R$ de l'axe :
  \[ J_\Delta = mR^2 \]
  \item \textbf{Cylindre plein ou disque plein} de rayon $R$ : la masse est répartie de l'axe jusqu'au rayon $R$ :
  \[ J_\Delta = \frac{1}{2}mR^2 \]
  \item \textbf{Tige} mince de longueur $L$, par rapport à un axe perpendiculaire passant par son milieu :
  \[ J_\Delta = \frac{1}{12}mL^2 \]
\end{itemize}
\end{aretenir}

Remarque la cohérence avec la propriété précédente : à masse et rayon égaux, le cerceau ($J = mR^2$, masse entièrement au bord) a un moment d'inertie \emph{double} de celui du disque plein ($J = \tfrac{1}{2}mR^2$, masse répartie jusqu'au centre). C'est exactement ce que prédit la formule $J_\Delta = \sum_i m_i r_i^2$.

\courssub{Analogie formelle entre translation et rotation}

L'expression $E_c = \dfrac{1}{2}J_\Delta\omega^2$ a exactement la même structure que $E_c = \dfrac{1}{2}mv^2$. Ce n'est pas un hasard : toute la mécanique de la rotation est construite en parallèle de celle de la translation, selon le dictionnaire suivant.

\begin{popfigure}
\begin{center}
\begin{tabular}{|c|c|c|}
    \hline
    \rowcolor{popblueL} \textbf{Grandeur} & \textbf{Translation} & \textbf{Rotation} \\
    \hline
    Position & $x$ ($\mathrm{m}$) & $\theta$ ($\mathrm{rad}$) \\
    \hline
    Vitesse & $v$ ($\mathrm{m\cdot s^{-1}}$) & $\omega$ ($\mathrm{rad\cdot s^{-1}}$) \\
    \hline
    Inertie & masse $m$ ($\mathrm{kg}$) & $J_\Delta$ ($\mathrm{kg\cdot m^2}$) \\
    \hline
    Énergie cinétique & $\dfrac{1}{2}mv^2$ & $\dfrac{1}{2}J_\Delta\omega^2$ \\
    \hline
\end{tabular}
\end{center}
\legende{Dictionnaire translation -- rotation. À chaque grandeur de la translation correspond une grandeur de la rotation : la masse $m$ est remplacée par le moment d'inertie $J_\Delta$, et la vitesse linéaire $v$ par la vitesse angulaire $\omega$.}
\end{popfigure}

Cette analogie est un formidable outil de mémorisation : si tu connais la formule en translation, tu connais celle de la rotation, moyennant les correspondances $m \leftrightarrow J_\Delta$ et $v \leftrightarrow \omega$.

\courssub{Méthode de calcul et exemples concrets}

\begin{methode}[Calculer l'énergie cinétique d'un solide en rotation]
\begin{enumerate}
  \item \textbf{Identifier} la nature du solide (disque, cerceau, tige...) et choisir l'expression de $J_\Delta$ correspondante.
  \item \textbf{Convertir} toutes les données en unités SI : masses en $\mathrm{kg}$, rayons et longueurs en $\mathrm{m}$.
  \item \textbf{Convertir la vitesse angulaire en $\mathrm{rad\cdot s^{-1}}$}. Si la rotation est donnée en tours par minute ($N$ en $\mathrm{tr/min}$) :
  \[ \omega = \frac{2\pi N}{60} \]
  (un tour vaut $2\pi$ radians, une minute vaut $60$ secondes).
  \item \textbf{Appliquer} $E_c = \dfrac{1}{2}J_\Delta\omega^2$ et exprimer le résultat en joules avec un nombre raisonnable de chiffres significatifs.
\end{enumerate}
\end{methode}

\begin{attention}[Le piège du tr/min]
L'erreur la plus fréquente au Probatoire consiste à injecter directement la vitesse de rotation en $\mathrm{tr/min}$ dans la formule $E_c = \dfrac{1}{2}J_\Delta\omega^2$. C'est \textbf{faux} : la formule exige $\omega$ en $\mathrm{rad\cdot s^{-1}}$. Une rotation de $1500~\mathrm{tr/min}$ n'est \emph{pas} $\omega = 1500$, mais $\omega = \dfrac{2\pi \times 1500}{60} \approx 157~\mathrm{rad\cdot s^{-1}}$. Oublier la conversion fait perdre un facteur $\left(\dfrac{2\pi}{60}\right)^2 \approx 0{,}011$, soit une énergie environ 91 fois trop grande !
\end{attention}

\begin{exemplebox}[Le volant d'inertie d'un groupe électrogène]
Dans de nombreux quartiers de Yaoundé ou de Douala, les coupures d'électricité rendent les groupes électrogènes indispensables. Ces appareils comportent un \textbf{volant d'inertie} : un disque lourd qui, une fois lancé, stocke de l'énergie cinétique de rotation et régularise la rotation du moteur.

Prenons un volant assimilé à un \emph{disque plein} de masse $m = 20~\mathrm{kg}$ et de rayon $R = 30~\mathrm{cm} = 0{,}30~\mathrm{m}$, tournant à $N = 1500~\mathrm{tr/min}$.

\textbf{Étape 1 -- Moment d'inertie} (disque plein) :
\[ J_\Delta = \frac{1}{2}mR^2 = \frac{1}{2}\times 20 \times (0{,}30)^2 = 0{,}90~\mathrm{kg\cdot m^2} \]

\textbf{Étape 2 -- Conversion de la vitesse angulaire} :
\[ \omega = \frac{2\pi N}{60} = \frac{2\pi \times 1500}{60} = 50\pi \approx 157~\mathrm{rad\cdot s^{-1}} \]

\textbf{Étape 3 -- Énergie cinétique} :
\[ E_c = \frac{1}{2}J_\Delta\omega^2 = \frac{1}{2}\times 0{,}90\times (157)^2 \approx 1{,}1\times 10^{4}~\mathrm{J} \]

Le volant stocke environ $11~\mathrm{kJ}$ d'énergie cinétique : de quoi soulever une masse d'une tonne de plus d'un mètre ! C'est cette réserve qui permet au groupe électrogène de tourner régulièrement entre deux explosions du moteur.
\end{exemplebox}

\begin{exoresolu}[Roue de moto et tambour de machine à laver]
\textbf{Énoncé.} (a) La roue avant d'une mototaxi est assimilée à un cerceau de masse $m = 8{,}0~\mathrm{kg}$ et de rayon $R = 30~\mathrm{cm}$. Elle tourne à $300~\mathrm{tr/min}$. Calcule son énergie cinétique de rotation. (b) Le tambour d'une machine à laver, assimilé à un cylindre creux (cerceau) de masse $m' = 4{,}0~\mathrm{kg}$ et de rayon $R' = 25~\mathrm{cm}$, essore à $1200~\mathrm{tr/min}$. Calcule son énergie cinétique de rotation et compare.
\tcblower
\textbf{Solution.} (a) \emph{Roue de moto.} Cerceau : $J_\Delta = mR^2 = 8{,}0\times(0{,}30)^2 = 0{,}72~\mathrm{kg\cdot m^2}$.

Conversion : $\omega = \dfrac{2\pi\times 300}{60} = 10\pi \approx 31{,}4~\mathrm{rad\cdot s^{-1}}$.

Énergie cinétique :
\[ E_c = \frac{1}{2}\times 0{,}72\times(31{,}4)^2 \approx 3{,}6\times10^{2}~\mathrm{J} \]

(b) \emph{Tambour.} $J'_\Delta = m'R'^2 = 4{,}0\times(0{,}25)^2 = 0{,}25~\mathrm{kg\cdot m^2}$.

Conversion : $\omega' = \dfrac{2\pi\times 1200}{60} = 40\pi \approx 126~\mathrm{rad\cdot s^{-1}}$.

Énergie cinétique :
\[ E'_c = \frac{1}{2}\times 0{,}25\times(126)^2 \approx 2{,}0\times10^{3}~\mathrm{J} \]

\emph{Comparaison.} Bien que le tambour soit deux fois moins lourd et plus petit que la roue, il emmagasine environ $5{,}5$ fois plus d'énergie, car il tourne 4 fois plus vite et que l'énergie cinétique varie comme le \textbf{carré} de $\omega$ (facteur $4^2 = 16$). C'est pourquoi l'essorage est la phase la plus sollicitante mécaniquement pour une machine à laver.
\end{exoresolu}

\begin{attention}[Autres erreurs fréquentes]
\begin{itemize}
  \item \textbf{Confondre cerceau et disque plein} : l'énoncé précise toujours (directement ou par le moment d'inertie donné) de quel modèle il s'agit. Un cylindre \emph{plein} : $J = \tfrac{1}{2}mR^2$ ; un cylindre \emph{creux} ou une roue à jante : $J = mR^2$.
  \item \textbf{Oublier de convertir le rayon en mètres} : $R = 30~\mathrm{cm}$ doit devenir $0{,}30~\mathrm{m}$ avant d'être élevé au carré.
  \item \textbf{Prendre le diamètre pour le rayon} : lis attentivement l'énoncé !
\end{itemize}
\end{attention}

\begin{aretenir}[L'essentiel de la section]
\begin{itemize}
  \item Un solide en rotation autour d'un axe fixe $\Delta$ possède une énergie cinétique $E_c = \dfrac{1}{2}J_\Delta\omega^2$, avec $\omega$ obligatoirement en $\mathrm{rad\cdot s^{-1}}$.
  \item Le moment d'inertie $J_\Delta = \sum_i m_i r_i^2$ (en $\mathrm{kg\cdot m^2}$) dépend de la masse et de sa répartition autour de l'axe.
  \item À retenir : cerceau $J = mR^2$ ; disque ou cylindre plein $J = \tfrac{1}{2}mR^2$ ; tige $J = \tfrac{1}{12}mL^2$.
  \item Conversion indispensable : $\omega = \dfrac{2\pi N}{60}$ si $N$ est en $\mathrm{tr/min}$.
  \item Analogie : $m \leftrightarrow J_\Delta$ et $v \leftrightarrow \omega$.
\end{itemize}
\end{aretenir}

Nous savons désormais évaluer l'énergie cinétique d'un solide en translation \emph{ou} en rotation. Mais que se passe-t-il lorsque le solide fait les deux à la fois, comme une roue qui roule sur la route ? C'est l'objet de la section suivante.

\coursec{Solide en mouvement combiné : translation et rotation}

Dans la section précédente, nous avons appris à calculer l'énergie cinétique d'un solide en \cle{rotation} autour d'un axe fixe : $E_c = \dfrac{1}{2}J_\Delta\,\omega^2$. Mais ce cas est idéalisé : l'axe de rotation était \emph{immobile}. Or, dans la vie quotidienne, les objets qui tournent se déplacent en même temps. La roue d'une mototaxi qui file sur le boulevard de la Liberté à Douala tourne sur elle-même \emph{et} avance avec le véhicule. Un ballon de football frappé avec effet vole vers les buts tout en tournoyant. Comment exprimer l'énergie cinétique de tels objets ? C'est l'objet de cette section, qui aboutit à l'une des formules les plus utiles de la mécanique : le \cle{théorème de König}.

\courssub{Observation : une roue qui roule fait deux mouvements à la fois}

\begin{experience}[Décomposer le mouvement d'une roue]
Prenons une roue de bicyclette posée au sol et marquons d'un point de peinture blanche son centre $G$ et un point $M$ de sa jante. Faisons rouler la roue sur le sol et photographions la scène en pose longue (ou observons une vidéo image par image).

\textbf{Observations :}
\begin{itemize}
\item le centre $G$ se déplace en ligne droite, parallèlement au sol, à vitesse constante : c'est un mouvement de \cle{translation} ;
\item le point $M$ de la jante décrit une courbe en forme d'arches (une \emph{cycloïde}) : il tourne autour de $G$ tout en étant entraîné vers l'avant ;
\item le point de contact $I$ entre la roue et le sol est, à chaque instant, \emph{immobile} par rapport au sol si la roue ne glisse pas : c'est ce qui permet l'adhérence des pneus sur la route.
\end{itemize}

\textbf{Interprétation :} le mouvement de la roue est la \emph{combinaison} (on dit la \emph{superposition}) de deux mouvements simples : une translation à la vitesse $\vec{v}_G$ du centre d'inertie, et une rotation de vitesse angulaire $\omega$ autour de l'axe passant par $G$.
\end{experience}

\begin{popfigure}
\begin{center}
\begin{tikzpicture}[scale=1.0, >=Stealth]
% Sol
\fill[popcream] (-0.5,-0.15) rectangle (11.5,0);
\draw[popdark, thick] (-0.5,0) -- (11.5,0);
\foreach \x in {0,0.6,...,11.4} \draw[popmuted] (\x,0) -- (\x-0.25,-0.15);
% Roue
\draw[popblue, very thick, fill=popblueL] (3,1.5) circle (1.5);
\foreach \a in {0,45,...,315} \draw[popblue!70] (3,1.5) -- ({3+1.35*cos(\a)},{1.5+1.35*sin(\a)});
\fill[popdark] (3,1.5) circle (0.07) node[above right=1pt] {$G$};
\fill[poporange] ({3+1.5*cos(60)},{1.5+1.5*sin(60)}) circle (0.09) node[right=2pt, popdark] {$M$};
\fill[popdark] (3,0) circle (0.07) node[below=2pt] {$I$};
% Vecteur vitesse de G
\draw[->, poporange, very thick] (3,1.5) -- (5.4,1.5) node[midway, above, popdark] {$\vec{v}_G$};
% Rotation
\draw[->, poppurple, very thick] (4.6,1.5) arc (0:-120:1.6) node[midway, right=3pt, popdark] {$\omega$};
% Flèche translation globale
\draw[->, popgreen, very thick] (0.3,3.6) -- (2.6,3.6) node[midway, above, popdark] {translation de $G$};
% Décomposition à droite
\node[popdark, font=\bfseries] at (8.3,3.9) {Décomposition du mouvement};
% Translation seule
\draw[popblue, thick] (7.2,1.5) circle (0.55);
\fill[popdark] (7.2,1.5) circle (0.05);
\draw[->, poporange, thick] (7.2,1.5) -- (8.5,1.5) node[midway, above, popdark] {$\vec{v}_G$};
\node[popdark] at (7.2,0.6) {translation};
\node[popdark] at (8.9,1.5) {$+$};
% Rotation seule
\draw[popblue, thick] (10.1,1.5) circle (0.55);
\fill[popdark] (10.1,1.5) circle (0.05);
\draw[->, poppurple, thick] (10.85,1.5) arc (0:-110:0.75) node[midway, right=1pt, popdark] {$\omega$};
\node[popdark] at (10.1,0.6) {rotation autour de $G$};
\end{tikzpicture}
\end{center}
\legende{Une roue roulant sur une route : son mouvement est la somme d'une translation du centre d'inertie $G$ à la vitesse $\vec{v}_G$ et d'une rotation de vitesse angulaire $\omega$ autour de l'axe passant par $G$. Le point de contact $I$ est instantanément au repos dans le roulement sans glissement.}
\end{popfigure}

\begin{aretenir}[Mouvement combiné]
Un solide est en \cle{mouvement combiné} (translation et rotation) lorsque son centre d'inertie $G$ se déplace \emph{et} que le solide tourne autour d'un axe passant par $G$. C'est le cas général du mouvement d'un solide : roues de véhicules, ballons, toupies qui se déplacent, planètes (révolution + rotation propre).
\end{aretenir}

\courssub{Théorème de König : l'énergie cinétique totale}

Puisque le mouvement se décompose en deux mouvements simples, il est naturel de penser que l'énergie cinétique totale se décompose elle aussi en deux contributions. C'est exact, et c'est un résultat général et remarquable de la mécanique.

\begin{propriete}[Théorème de König (admis)]
Pour un solide de masse $m$ en mouvement quelconque, l'énergie cinétique totale est la somme de deux termes :
\[ E_c = \underbrace{\dfrac{1}{2}\,m\,v_G^{2}}_{\text{translation de }G} \;+\; \underbrace{\dfrac{1}{2}\,J_G\,\omega^{2}}_{\text{rotation autour de }G} \]
où :
\begin{itemize}
\item $v_G$ est la vitesse du \cle{centre d'inertie} $G$ du solide (en $\SI{}{m/s}$) ;
\item $J_G$ est le moment d'inertie du solide par rapport à l'axe de rotation \emph{passant par $G$} (en $\SI{}{kg.m^2}$) ;
\item $\omega$ est la vitesse angulaire de rotation (en $\SI{}{rad/s}$).
\end{itemize}
Ce théorème est \emph{admis} en Première C ; sa démonstration n'est pas exigible au Probatoire. En revanche, savoir l'\emph{appliquer} est exigible.
\end{propriete}

\begin{attention}[L'axe de rotation passe par G]
Dans le second terme, le moment d'inertie doit être calculé par rapport à l'axe passant par le \emph{centre d'inertie} $G$, et non par rapport à un axe quelconque (par exemple l'axe du point de contact avec le sol). De même, c'est bien la vitesse de $G$ qui intervient dans le premier terme, pas celle d'un point de la périphérie.
\end{attention}

\textbf{Vérification par analyse dimensionnelle.} Les deux termes doivent être homogènes à une énergie ($\SI{}{J} = \SI{}{kg.m^2/s^2}$) :
\begin{align*}
\left[\tfrac{1}{2}mv_G^2\right] &= \SI{}{kg}\times\left(\SI{}{m/s}\right)^2 = \SI{}{kg.m^2/s^2} = \SI{}{J} \quad\checkmark\\
\left[\tfrac{1}{2}J_G\omega^2\right] &= \SI{}{kg.m^2}\times\left(\SI{}{rad/s}\right)^2 = \SI{}{kg.m^2/s^2} = \SI{}{J} \quad\checkmark
\end{align*}
(le radian est sans dimension). La formule est dimensionnellement cohérente.

\courssub{Condition de roulement sans glissement}

Reprenons notre roue de rayon $R$ qui roule sur le sol. Lorsque la roue effectue \emph{un tour complet} (angle $2\pi$ radians), son centre $G$ avance exactement d'une distance égale au périmètre $2\pi R$ : chaque point du pneu a touché le sol une fois et une seule, sans frotter en glissant. En divisant par la durée $T$ d'un tour :
\[ \frac{2\pi R}{T} = R\,\frac{2\pi}{T} \quad\Longrightarrow\quad v_G = R\,\omega \]

\begin{propriete}[Roulement sans glissement]
Pour une roue (ou un cylindre, une sphère) de rayon $R$ roulant \emph{sans glisser} sur une surface, la vitesse du centre d'inertie et la vitesse angulaire sont liées par :
\[ \boxed{\;v_G = R\,\omega\;} \]
Si la roue glisse (freinage bloqué, dérapage sur route mouillée), cette relation n'est \emph{plus} vérifiée.
\end{propriete}

\begin{savaistu}[Pourquoi les pneus ont-ils des sculptures ?]
Quand une mototaxi freine brutalement sur le sable d'une route non bitumée, les roues peuvent se bloquer : elles glissent au lieu de rouler, la relation $v_G = R\omega$ est rompue, et l'adhérence chute brutalement. Les sculptures des pneus et les systèmes de freinage modernes servent précisément à maintenir le \emph{roulement sans glissement}, condition d'un freinage efficace et contrôlé.
\end{savaistu}

\courssub{Application fondamentale : énergie cinétique d'une roue pleine}

Considérons une \cle{roue pleine} (cylindre homogène) de masse $m$ et de rayon $R$, roulant sans glisser à la vitesse $v_G = v$. Son moment d'inertie par rapport à son axe est :
\[ J_G = \dfrac{1}{2}\,m\,R^2 \]

Appliquons le théorème de König en remplaçant $\omega$ par $\dfrac{v}{R}$ (condition de roulement sans glissement) :
\begin{align*}
E_c &= \dfrac{1}{2}mv^2 + \dfrac{1}{2}J_G\,\omega^2 \\
&= \dfrac{1}{2}mv^2 + \dfrac{1}{2}\times\left(\dfrac{1}{2}mR^2\right)\times\left(\dfrac{v}{R}\right)^2 \\
&= \dfrac{1}{2}mv^2 + \dfrac{1}{4}mv^2 \\
&= \boxed{\;\dfrac{3}{4}mv^2\;}
\end{align*}

Remarquons la simplification remarquable : $R$ disparaît ! L'énergie cinétique totale d'une roue pleine qui roule sans glisser ne dépend que de sa masse et de sa vitesse, pas de son rayon.

\begin{methode}[Calculer l'énergie cinétique d'un solide roulant sans glisser]
\begin{enumerate}
\item Identifier la masse $m$, la vitesse $v_G$ du centre d'inertie et la forme du solide (pour connaître $J_G$).
\item Écrire le théorème de König : $E_c = \dfrac{1}{2}mv_G^2 + \dfrac{1}{2}J_G\,\omega^2$.
\item Utiliser la condition de roulement sans glissement pour éliminer $\omega$ : $\omega = \dfrac{v_G}{R}$.
\item Simplifier : $R^2$ se simplifie toujours ; regrouper les termes en $v_G^2$.
\item Calculer numériquement en veillant aux unités SI ($m$ en kg, $v$ en m/s, $E_c$ en J).
\end{enumerate}
\end{methode}

\begin{attention}[Ne jamais oublier le terme de rotation !]
L'erreur la plus fréquente au Probatoire est d'écrire $E_c = \dfrac{1}{2}mv^2$ pour une roue qui roule. C'est \textbf{faux} : cette expression ne tient compte que de la translation. Pour une roue pleine, le terme de rotation $\dfrac{1}{4}mv^2$ représente la \emph{moitié} du terme de translation $\dfrac{1}{2}mv^2$, soit le \emph{tiers} de l'énergie totale. L'oublier, c'est sous-estimer l'énergie de $33\,\%$ !
\end{attention}

\courssub{Comparaison : solide roulant contre solide glissant}

Comparons, à masse $m$ et vitesse $v$ égales :
\begin{itemize}
\item un bloc qui \emph{glisse} sans tourner (translation pure) : $E_c^{\text{gliss}} = \dfrac{1}{2}mv^2$ ;
\item une roue pleine de même masse qui \emph{roule} : $E_c^{\text{roul}} = \dfrac{3}{4}mv^2 = 1{,}5\times\dfrac{1}{2}mv^2$.
\end{itemize}

\begin{aretenir}[Conclusion physique]
À vitesse égale et à masse égale, un solide \emph{roulant} possède \textbf{plus} d'énergie cinétique qu'un solide \emph{glissant}, car il faut aussi fournir de l'énergie pour le faire tourner. Pour la roue pleine, l'énergie totale vaut $1{,}5$ fois celle de la translation seule. C'est pourquoi mettre en mouvement (ou arrêter) un véhicule exige plus d'énergie que le simple calcul $\frac{1}{2}mv^2$ ne le suggère : les roues, les arbres de transmission et le moteur stockent de l'énergie de rotation.
\end{aretenir}

\begin{popfigure}
\begin{center}
\begin{tikzpicture}
\begin{axis}[
    ybar, bar width=0.9cm,
    width=0.85\linewidth, height=6.2cm,
    symbolic x coords={Translation, Rotation, Totale},
    xtick=data,
    ymin=0, ymax=0.85,
    ylabel={$E_c$ (en unités de $mv^2$)},
    ylabel style={font=\small},
    nodes near coords,
    every node near coord/.append style={font=\small\bfseries, color=popdark},
    axis lines=left,
    x tick label style={font=\small},
    y tick label style={font=\small},
    enlarge x limits=0.25,
    grid=major, grid style={popline!40},
]
\addplot[fill=popblue, draw=popdark] coordinates {(Translation,0.5) (Rotation,0.25) (Totale,0.75)};
\end{axis}
\end{tikzpicture}
\end{center}
\legende{Répartition de l'énergie cinétique d'une roue pleine roulant sans glisser : $\frac{1}{2}mv^2$ pour la translation, $\frac{1}{4}mv^2$ pour la rotation, soit $\frac{3}{4}mv^2$ au total. La rotation compte pour un tiers de l'énergie totale.}
\end{popfigure}

\courssub{Exemple chiffré détaillé}

\begin{exoresolu}[Énergie cinétique d'une roue pleine de mototaxi]
Une roue de mototaxi, assimilée à un cylindre plein homogène de masse $m = \SI{15}{kg}$ et de rayon $R = \SI{0,30}{m}$, roule sans glisser à la vitesse $v = \SI{10}{m/s}$ (soit $\SI{36}{km/h}$).
\begin{enumerate}
\item Calculer la vitesse angulaire $\omega$ de la roue.
\item Calculer son moment d'inertie $J_G$.
\item Calculer les énergies cinétiques de translation, de rotation, puis l'énergie cinétique totale.
\item Comparer avec l'énergie d'un bloc de même masse glissant à la même vitesse.
\end{enumerate}
\tcblower
\textbf{Solution détaillée.}

\textbf{1. Vitesse angulaire.} La condition de roulement sans glissement donne :
\[ \omega = \frac{v}{R} = \frac{10}{0{,}30} \approx \SI{33,3}{rad/s} \]
soit environ $5{,}3$ tours par seconde.

\textbf{2. Moment d'inertie} de la roue pleine par rapport à son axe :
\[ J_G = \frac{1}{2}mR^2 = \frac{1}{2}\times 15\times (0{,}30)^2 = \frac{1}{2}\times 15\times 0{,}09 = \SI{0,675}{kg.m^2} \]

\textbf{3. Énergies cinétiques.}
\begin{align*}
E_c^{\text{trans}} &= \frac{1}{2}mv^2 = \frac{1}{2}\times 15\times 10^2 = \SI{750}{J} \\
E_c^{\text{rot}} &= \frac{1}{2}J_G\,\omega^2 = \frac{1}{2}\times 0{,}675\times (33{,}3)^2 \approx \SI{375}{J} \\
E_c^{\text{totale}} &= 750 + 375 = \SI{1125}{J}
\end{align*}
Vérification rapide par la formule établie : $E_c = \dfrac{3}{4}mv^2 = \dfrac{3}{4}\times 15\times 100 = \SI{1125}{J}$. \checkmark

\textbf{4. Comparaison.} Un bloc de $\SI{15}{kg}$ glissant à $\SI{10}{m/s}$ ne posséderait que $\frac{1}{2}mv^2 = \SI{750}{J}$. La roue roulante possède donc $\SI{375}{J}$ de plus, soit $50\,\%$ d'énergie supplémentaire stockée dans la rotation. C'est cette énergie qu'il faut également dissiper au freinage.
\end{exoresolu}

\begin{exercice}[Roue de brouette]
Une roue de brouette, assimilée à un cylindre plein de masse $m = \SI{4,0}{kg}$ et de rayon $R = \SI{0,20}{m}$, roule sans glisser à $v = \SI{2,0}{m/s}$ sur un chantier de Yaoundé. Calculer : (a) sa vitesse angulaire ; (b) son énergie cinétique totale ; (c) la part (en pourcentage) de l'énergie de rotation dans l'énergie totale.
\end{exercice}

\begin{corrige}[Exercice : roue de brouette]
(a) $\omega = \dfrac{v}{R} = \dfrac{2{,}0}{0{,}20} = \SI{10}{rad/s}$.

(b) $E_c = \dfrac{3}{4}mv^2 = \dfrac{3}{4}\times 4{,}0\times (2{,}0)^2 = \SI{12}{J}$.

(c) $E_c^{\text{rot}} = \dfrac{1}{4}mv^2 = \SI{4}{J}$, soit $\dfrac{4}{12} = \dfrac{1}{3} \approx 33\,\%$ de l'énergie totale. Ce pourcentage est indépendant de la masse et de la vitesse : c'est une propriété générale de la roue pleine.
\end{corrige}

\begin{savaistu}[Et la Terre dans tout ça ?]
Le théorème de König s'applique aussi aux astres ! L'énergie cinétique de la Terre est la somme de son énergie de translation autour du Soleil ($\frac{1}{2}Mv_G^2$, gigantesque) et de son énergie de rotation propre sur elle-même ($\frac{1}{2}J_G\omega^2$). Le ralentissement progressif de la rotation terrestre, dû aux marées, montre que cette énergie de rotation se dissipe lentement : les journées s'allongent d'environ $2$ millisecondes par siècle.
\end{savaistu}

\begin{aretenir}[L'essentiel de la section]
\begin{itemize}
\item Un solide en mouvement général combine une \cle{translation} de son centre d'inertie $G$ et une \cle{rotation} autour d'un axe passant par $G$.
\item Théorème de König (admis) : $E_c = \dfrac{1}{2}mv_G^2 + \dfrac{1}{2}J_G\,\omega^2$.
\item Roulement sans glissement : $v_G = R\,\omega$.
\item Roue pleine roulant sans glisser : $E_c = \dfrac{3}{4}mv^2$, dont un tiers en rotation.
\item À masse et vitesse égales, un solide roulant est toujours \emph{plus énergétique} qu'un solide glissant.
\end{itemize}
\end{aretenir}

Nous disposons désormais de tous les outils pour exprimer l'énergie cinétique d'un solide, quel que soit son mouvement. La prochaine étape est de comprendre \emph{comment cette énergie varie} lorsque des forces agissent sur le solide : ce sera le \cle{théorème de l'énergie cinétique}, l'un des résultats les plus puissants de la mécanique de Première C.

\coursec{Théorème de l'énergie cinétique}

Dans la section précédente, nous avons vu comment calculer l'énergie cinétique d'un solide, qu'il soit en translation pure, en rotation pure autour d'un axe fixe, ou en mouvement combiné (translation du centre d'inertie + rotation autour de ce centre). Nous savons maintenant exprimer $E_c$ en fonction des masses, vitesses, moments d'inertie et vitesses angulaires. Mais comment cette énergie cinétique \emph{évolue-t-elle} au cours du temps ? Qu'est-ce qui la fait augmenter, diminuer ou rester constante ? C'est précisément la réponse que nous apporte le \textbf{théorème de l'énergie cinétique}, l'un des outils les plus puissants de la mécanique newtonienne pour résoudre des problèmes sans avoir à déterminer la loi horaire complète du mouvement.

\courssub{Rappel : travail d'une force et travail des forces usuelles}

Avant d'énoncer le théorème, rappelons la définition du travail d'une force, notion centrale qui fait le lien entre la cause (les forces) et l'effet (variation d'énergie cinétique).

\begin{definition}[Travail d'une force constante]
Soit une force $\vec{F}$ constante (en norme, direction et sens) appliquée à un point matériel $M$. Si $M$ se déplace de $A$ à $B$, le travail de $\vec{F}$ le long du déplacement $\overrightarrow{AB}$ est le scalaire :
\[
W_{AB}(\vec{F}) = \vec{F} \cdot \overrightarrow{AB} = F \, AB \, \cos(\vec{F}, \overrightarrow{AB})
\]
L'unité SI est le \cle{joule} (\textbf{J}). $1\,\text{J} = 1\,\text{N}\cdot\text{m} = 1\,\text{kg}\cdot\text{m}^2\cdot\text{s}^{-2}$.
\end{definition}

\begin{propriete}[Travail des forces usuelles en translation]
Pour un solide en translation (tous les points ont le même déplacement $\overrightarrow{AB}$) :
\begin{itemize}
    \item \textbf{Poids} $\vec{P} = m\vec{g}$ : $W_{AB}(\vec{P}) = \vec{P} \cdot \overrightarrow{AB} = mg\,(z_A - z_B) = \pm mgh$.
        \begin{itemize}
            \item Descente ($z_B < z_A$) : $W > 0$ (travail \cle{moteur}).
            \item Montée ($z_B > z_A$) : $W < 0$ (travail \cle{résistant}).
        \end{itemize}
    \item \textbf{Réaction normale} $\vec{R}_N$ (plan horizontal ou incliné sans frottement) : $\vec{R}_N \perp \overrightarrow{AB} \Rightarrow W = 0$.
    \item \textbf{Frottements solides} $\vec{f}$ (opposés au mouvement) : $W_{AB}(\vec{f}) = -f\, AB < 0$ (toujours résistant).
    \item \textbf{Tension d'un fil} $\vec{T}$ : $W = T\, AB \cos\alpha$ (moteur si $\alpha < 90^\circ$, résistant si $\alpha > 90^\circ$).
\end{itemize}
\end{propriete}

\begin{attention}[Pièges classiques sur le travail du poids]
\begin{itemize}
    \item \textbf{Erreur de signe} : On écrit souvent $W = mgh$ sans réfléchir au sens. \textbf{Règle} : $W = mg(z_A - z_B)$. Si l'objet descend, $z_A > z_B \Rightarrow W > 0$. Si l'objet monte, $W < 0$.
    \item \textbf{Trajectoire non verticale} : Le travail du poids ne dépend \textbf{que} de la différence d'altitude $h = z_A - z_B$, pas de la longueur du chemin parcouru (pente, courbe, chute libre). C'est une force \cle{conservative}.
    \item \textbf{Réaction normale} : Elle ne travaille \textbf{jamais} pour un solide glissant sur un plan fixe (déplacement tangent au plan, réaction normale). Elle travaille seulement si le plan lui-même bouge (ex: ascenseur, camion en accélération).
\end{itemize}
\end{attention}

\courssub{Énoncé du théorème de l'énergie cinétique}

\begin{propriete}[Théorème de l'énergie cinétique — Énoncé général]
Dans un \cle{référentiel galiléen}, pour tout système solide (indéformable) passant de la position $A$ à la position $B$, la variation de son énergie cinétique est égale à la somme des travaux de \textbf{toutes les forces extérieures} appliquées au solide :
\[
\boxed{\Delta E_c = E_c(B) - E_c(A) = \sum_{\text{forces extérieures}} W_{AB}(\vec{F}_i)}
\]
Ce théorème est valable :
\begin{itemize}
    \item Pour \textbf{toute trajectoire} (rectiligne, circulaire, curviligne quelconque).
    \item Pour des forces \textbf{variables} (en norme, direction, sens) — le travail se calcule alors par intégration $W = \int \vec{F} \cdot d\vec{l}$.
    \item Pour un solide en \textbf{translation}, en \textbf{rotation} autour d'un axe fixe, ou en \textbf{mouvement combiné}.
\end{itemize}
\end{propriete}

\begin{aretenir}[Interprétation physique]
\begin{itemize}
    \item $\sum W > 0$ (travail \textbf{moteur} net) $\Rightarrow \Delta E_c > 0 \Rightarrow$ la vitesse (ou la vitesse angulaire) \textbf{augmente}.
    \item $\sum W < 0$ (travail \textbf{résistant} net) $\Rightarrow \Delta E_c < 0 \Rightarrow$ la vitesse \textbf{diminue}.
    \item $\sum W = 0 \Rightarrow \Delta E_c = 0 \Rightarrow E_c(B) = E_c(A)$ : l'énergie cinétique est \textbf{conservée} (vitesse constante en norme).
\end{itemize}
Le théorème relie l'état \textbf{initial} (A) et l'état \textbf{final} (B) sans avoir besoin de connaître les détails du mouvement entre les deux (loi horaire, accélération à chaque instant). C'est un \textbf{principe global}.
\end{aretenir}

\courssub{Démonstration guidée (cas simple : mouvement rectiligne, force résultante constante)}

Bien que le théorème soit général, sa démonstration dans le cas simple permet d'en comprendre l'origine newtonienne. Suivons les étapes logiques.

\begin{experience}[Démonstration : du principe fondamental à l'énergie]
\textbf{Matériel conceptuel} : Un solide de masse $m$ (modélisé en point matériel) se déplaçant sur un axe $Ox$ galiléen, soumis à une force résultante $\vec{F}_{\text{rés}} = F\,\vec{u}_x$ \textbf{constante}.
\textbf{Protocole et raisonnement} :
\begin{enumerate}
    \item \textbf{Deuxième loi de Newton} (PFDS) : $\vec{F}_{\text{rés}} = m \vec{a}$. Comme le mouvement est rectiligne sur $Ox$, $F = m a$ (relation scalaire algébrique).
    \item \textbf{Relation cinématique} (M.R.U.A.) : Pour une accélération constante $a$, on a la relation entre vitesses et déplacement : $v_B^2 - v_A^2 = 2 a \, (x_B - x_A) = 2 a \, AB$.
    \item \textbf{Éliminer l'accélération} : On exprime $a = F/m$ et on substitue dans la relation cinématique :
    \[
    v_B^2 - v_A^2 = 2 \frac{F}{m} AB
    \]
    \item \textbf{Multiplier par $m/2$} :
    \[
    \frac{1}{2} m v_B^2 - \frac{1}{2} m v_A^2 = F \cdot AB
    \]
    \item \textbf{Identifier} : Le membre de gauche est $\Delta E_c = E_c(B) - E_c(A)$. Le membre de droite est le travail de la force résultante $W_{AB}(\vec{F}_{\text{rés}}) = F \cdot AB$ (car $\vec{F}_{\text{rés}} \parallel \overrightarrow{AB}$).
    \item \textbf{Conclusion} : $\Delta E_c = W_{AB}(\vec{F}_{\text{rés}}) = \sum W_{AB}(\vec{F}_i)$.
\end{enumerate}
\textbf{Interprétation} : Le théorème n'est qu'une \textbf{reformulation intégrée} de la deuxième loi de Newton. Il « somme » l'effet de la force sur le déplacement.
\end{experience}

\begin{popfigure}
\begin{tikzpicture}[node distance=1.5cm, >=Stealth, font=\small]
    % Nodes
    \node[draw, rounded corners, fill=popblueL, thick] (start) {PFDS : $\vec{F}_{\text{rés}} = m\vec{a}$};
    \node[draw, rounded corners, fill=popgreenL, thick, right=of start] (cin) {Relation cinématique MRUA : $v_B^2 - v_A^2 = 2a\,AB$};
    \node[draw, rounded corners, fill=poporangeL, thick, right=of cin] (sub) {Substitution $a = F/m$};
    \node[draw, rounded corners, fill=poppinkL, thick, right=of sub] (mult) {Multiplication par $m/2$};
    \node[draw, rounded corners, fill=popgoldL, thick, right=of mult] (end) {Théorème : $\Delta E_c = \sum W$};
    
    % Arrows
    \draw[thick, popdark] (start) -- (cin) node[midway, above] {\textit{1}};
    \draw[thick, popdark] (cin) -- (sub) node[midway, above] {\textit{2}};
    \draw[thick, popdark] (sub) -- (mult) node[midway, above] {\textit{3}};
    \draw[thick, popdark] (mult) -- (end) node[midway, above] {\textit{4}};
\end{tikzpicture}
\legende{Organigramme de la démonstration du théorème dans le cas simple (force constante, mouvement rectiligne).}
\end{popfigure}

\begin{propriete}[Généralisation admise]
La démonstration ci-dessus suppose une force résultante constante et un mouvement rectiligne. Au programme de Première, on \textbf{admet} que le théorème reste valable :
\begin{itemize}
    \item Pour des forces \textbf{variables} : le travail se calcule par intégration curviligne $W = \int_A^B \vec{F} \cdot d\vec{l}$, et la relation $\Delta E_c = \sum W$ est obtenue en intégrant $\vec{F} = m\vec{a}$ scalairement le long de la trajectoire ($\vec{F}\cdot d\vec{l} = m\vec{a}\cdot\vec{v}\,dt = d(\frac{1}{2}mv^2)$).
    \item Pour \textbf{toute trajectoire} curviligne.
    \item Pour un \textbf{solide indéformable} en rotation ou mouvement combiné : $\Delta E_c = \sum W_{\text{forces}} + \sum W_{\text{moments}}$. En rotation pure autour d'un axe fixe $\Delta$, $W_{\text{moment}} = \int M_\Delta \, d\theta = M_\Delta \Delta\theta$ (si moment constant).
\end{itemize}
\end{propriete}

\courssub{Cas particulier : solide en rotation autour d'un axe fixe}

Pour un solide tournant autour d'un axe fixe $\Delta$ (ex: roue, poulie, volant de moteur), l'énergie cinétique est $E_c = \frac{1}{2} J_\Delta \omega^2$. Le théorème s'écrit :

\begin{propriete}[Théorème de l'énergie cinétique en rotation]
\[
\Delta E_c = \frac{1}{2} J_\Delta (\omega_B^2 - \omega_A^2) = \sum_{\text{forces extérieures}} W_{AB}(\vec{F}_i) = \sum_{\text{moments extérieurs}} W_{AB}(\mathcal{M}_i)
\]
Le travail d'un moment $\mathcal{M}_\Delta$ constant (par rapport à l'axe $\Delta$) lors d'une rotation d'angle $\Delta\theta = \theta_B - \theta_A$ (en radians) est :
\[
W_{AB}(\mathcal{M}_\Delta) = \mathcal{M}_\Delta \cdot \Delta\theta
\]
\textbf{Convention de signe} : $\mathcal{M}_\Delta$ et $\Delta\theta$ sont des scalaires algébriques (positifs si même sens que l'orientation choisie de l'axe, négatifs sinon). Le travail est moteur si $\mathcal{M}_\Delta \cdot \Delta\theta > 0$.
\end{propriete}

\begin{exemplebox}[Freinage d'un volant d'inertie (contexte industriel camerounais)]
Un volant d'inertie d'une centrale thermique (ex: \textbf{centrale de Kribi} ou \textbf{Dibamba}) modélisé par un disque plein de masse $M = \SI{200}{kg}$ et rayon $R = \SI{0,80}{m}$ tourne à $\omega_A = \SI{150}{rad/s}$ ($\approx \SI{1430}{tr/min}$). On actionne un frein exerçant un couple résistant constant $\mathcal{M}_f = \SI{-500}{N.m}$. Quelle distance angulaire $\Delta\theta$ parcourt-il avant de s'arrêter ($\omega_B = 0$) ?
\begin{itemize}
    \item Moment d'inertie : $J_\Delta = \frac{1}{2} M R^2 = \frac{1}{2} \times 200 \times 0,80^2 = \SI{64}{kg.m^2}$.
    \item Énergie cinétique initiale : $E_c(A) = \frac{1}{2} J_\Delta \omega_A^2 = \frac{1}{2} \times 64 \times 150^2 = \SI{720000}{J} = \SI{720}{kJ}$.
    \item $E_c(B) = 0$.
    \item Travail du couple résistant : $W = \mathcal{M}_f \Delta\theta = -500 \Delta\theta$.
    \item Théorème : $\Delta E_c = W \Rightarrow 0 - 720\,000 = -500 \Delta\theta$.
    \item $\Delta\theta = \frac{720\,000}{500} = \SI{1440}{rad}$.
    \item Nombre de tours : $N = \frac{1440}{2\pi} \approx 229\,\text{tours}$.
\end{itemize}
\end{exemplebox}

\courssub{Application : solide sur plan incliné (bilan des forces et travaux)}

C'est le cas d'école par excellence pour illustrer le théorème : translation sur une pente, forces constantes, travaux faciles à calculer.

\begin{popfigure}
\begin{tikzpicture}[scale=1.2, >=Stealth]
    % Plan incliné
    \coordinate (O) at (0,0);
    \coordinate (B) at (5,0);
    \coordinate (A) at (5,3);
    \draw[thick, popdark] (O) -- (A) -- (B) -- cycle;
    \draw[pattern=north east lines, pattern color=popmuted] (O) -- (A) -- (B) -- cycle;
    
    % Solide (bloc)
    \coordinate (G) at (3.5, 2.1); % Centre du bloc sur la pente
    \draw[fill=popblue!20, draw=popblue, thick] ($(G)+(-0.4,-0.3)$) rectangle ($(G)+(0.4,0.3)$);
    \node at (G) {$m$};
    
    % Forces
    \draw[->, very thick, popred] (G) -- ($(G)+(0,-1.2)$) node[below right] {$\vec{P}=m\vec{g}$};
    \draw[->, very thick, popgreen] (G) -- ($(G)+(0.6,1.0)$) node[above right] {$\vec{R}_N$};
    \draw[->, very thick, poporange] (G) -- ($(G)+(-0.9,-0.5)$) node[below left] {$\vec{f}$};
    
    % Décomposition du poids
    \draw[->, thick, popred, dashed] (G) -- ($(G)+(0,-0.8)$) node[midway, right] {$P_\perp = mg\cos\alpha$};
    \draw[->, thick, popred, dashed] (G) -- ($(G)+(-0.65,-0.35)$) node[midway, below] {$P_\parallel = mg\sin\alpha$};
    
    % Angle alpha
    \draw (0.5,0) arc (0:31:0.5) node[midway, right, xshift=2pt] {$\alpha$};
    
    % Points A et B
    \node[popred, fill=white, circle, inner sep=1pt] at (A) {\textbf{A}};
    \node[popred, fill=white, circle, inner sep=1pt] at (B) {\textbf{B}};
    \draw[<->, thick, popdark] (A) -- (B) node[midway, above, sloped] {$AB = L$};
    
    % Coordonnées
    \draw[->] (O) -- (6,0) node[right] {$x$ (horizontal)};
    \draw[->] (O) -- (0,4) node[above] {$z$ (vertical)};
\end{tikzpicture}
\legende{Bilan des forces sur un solide glissant d'A vers B sur un plan incliné d'angle $\alpha$. Travaux : $W(P) = mgL\sin\alpha > 0$, $W(R_N) = 0$, $W(f) = -fL < 0$.}
\end{popfigure}

\begin{exoresolu}[Chariot sur une côte — Application directe]
\textbf{Énoncé} : Un chariot de masse $m = \SI{50}{kg}$ est lâché sans vitesse initiale ($v_A = 0$) du point A, au sommet d'une pente rectiligne de longueur $L = \SI{10}{m}$ et de dénivelé $h = \SI{5}{m}$. Le coefficient de frottement solide est $\mu = 0,10$. Déterminer sa vitesse $v_B$ au point B (bas de la pente). On prend $g = \SI{9,8}{N/kg}$.

\textbf{Solution détaillée} :
\begin{enumerate}
    \item \textbf{Système / Référentiel} : Système = chariot (solide indéformable en translation). Référentiel = Terre (galiléen au voisinage de Yaoundé).
    \item \textbf{Bilan des forces} (voir schéma ci-dessus) : $\vec{P}$, $\vec{R}_N$, $\vec{f}$.
    \item \textbf{Calcul des travaux entre A et B} :
    \begin{itemize}
        \item Poids : $W_{AB}(\vec{P}) = mg(z_A - z_B) = mgh = 50 \times 9,8 \times 5 = \SI{2450}{J}$. (Moteur, $>0$).
        \item Réaction normale : $\vec{R}_N \perp \overrightarrow{AB} \Rightarrow W_{AB}(\vec{R}_N) = 0$.
        \item Frottements : $f = \mu R_N = \mu mg\cos\alpha$. Mais on n'a pas besoin de $\alpha$ ! $W_{AB}(\vec{f}) = -f L = -\mu mg\cos\alpha \cdot L$.
            Or $\cos\alpha = \frac{\sqrt{L^2 - h^2}}{L} = \frac{\sqrt{100-25}}{10} = \frac{\sqrt{75}}{10} \approx 0,866$.
            $W_f = -0,10 \times 50 \times 9,8 \times 0,866 \times 10 \approx \SI{-424}{J}$. (Résistant, $<0$).
    \end{itemize}
    \item \textbf{Somme des travaux} : $\sum W = 2450 - 424 = \SI{2026}{J}$.
    \item \textbf{Théorème de l'énergie cinétique} :
    \[
    \Delta E_c = E_c(B) - E_c(A) = \frac{1}{2} m v_B^2 - 0 = \sum W = \SI{2026}{J}
    \]
    \item \textbf{Résolution} :
    \[
    v_B^2 = \frac{2 \times 2026}{50} = 81,04 \quad\Rightarrow\quad v_B = \sqrt{81,04} \approx \SI{9,0}{m/s}
    \]
    \textbf{Réponse} : Le chariot atteint une vitesse d'environ $\mathbf{\SI{9,0}{m/s}}$ (soit $\approx \SI{32,4}{km/h}$) au bas de la pente.
\end{enumerate}
\end{exoresolu}

\begin{attention}[Erreurs fréquentes à ne pas commettre]
\begin{itemize}
    \item \textbf{Oublier des forces} : Ne pas lister $\vec{R}_N$ (même si son travail est nul, il faut le mentionner et justifier $W=0$). Oublier les frottements s'ils existent.
    \item \textbf{Signe du travail du poids} : Écrire $W(P) = -mgh$ pour une descente. \textbf{Règle} : $W(P) = mg(z_A - z_B)$. Descente $\Rightarrow z_A > z_B \Rightarrow W > 0$.
    \item \textbf{Confondre travail et moment} : En rotation, on utilise le travail des \textbf{moments} ($M\Delta\theta$), pas le travail des forces directement (sauf si on calcule le travail de chaque force sur le déplacement de son point d'application).
    \item \textbf{Référentiel non galiléen} : Le théorème \textbf{ne s'applique pas} dans un référentiel non galiléen (ex: cabine d'ascenseur en accélération, camion qui freine) sans ajouter les travaux des forces d'inertie (forces fictives). Au Probatoire, on travaille \textbf{toujours} dans un référentiel galiléen (lié à la Terre).
    \item \textbf{Unités} : Vérifier que tous les travaux sont en joules (J). $1\,\text{J} = 1\,\text{N}\cdot\text{m}$. Attention aux cm, km, g, tonnes : \textbf{convertir en SI} avant de calculer.
\end{itemize}
\end{attention}

\courssub{Méthode : Appliquer le théorème de l'énergie cinétique en 5 étapes}

\begin{methode}[Résolution type « Théorème de l'énergie cinétique »]
\begin{enumerate}
    \item \textbf{Définir le système} (solide indéformable) et choisir le \textbf{référentiel galiléen} (généralement terrestre). Préciser si c'est translation, rotation ou combiné.
    \item \textbf{Faire le bilan des forces extérieures} (schéma obligatoire : poids, contact, frottements, tensions, actions mécaniques...). Identifier les forces \textbf{moteures} et \textbf{résistantes}.
    \item \textbf{Calculer le travail de chaque force} entre l'état initial (A) et l'état final (B).
        \begin{itemize}
            \item Forces constantes : $W = \vec{F}\cdot\overrightarrow{AB} = F\,AB\cos\theta$.
            \item Poids : $W(P) = mg(z_A - z_B)$.
            \item Forces variables / trajectoire courbe : Intégration (rare en 1ère C, mais savoir que $W = \int \vec{F}\cdot d\vec{l}$).
            \item Rotation : $W(\mathcal{M}) = \mathcal{M}\Delta\theta$ (moment constant).
        \end{itemize}
    \item \textbf{Écrire le théorème} : $\Delta E_c = E_c(B) - E_c(A) = \sum W_{AB}$.
        \begin{itemize}
            \item Translation : $\frac{1}{2}m(v_B^2 - v_A^2)$.
            \item Rotation : $\frac{1}{2}J_\Delta(\omega_B^2 - \omega_A^2)$.
            \item Combiné : $\frac{1}{2}m(v_{G,B}^2 - v_{G,A}^2) + \frac{1}{2}J_G(\omega_B^2 - \omega_A^2)$.
        \end{itemize}
    \item \textbf{Résoudre l'équation} pour l'inconnue demandée (vitesse, distance, coefficient de frottement, couple...). Vérifier l'homogénéité des unités et l'ordre de grandeur.
\end{enumerate}
\end{methode}

\courssub{Contexte d'utilisation : quand choisir le théorème ?}

Le théorème de l'énergie cinétique est \textbf{privilégié} par rapport à la dynamique newtonienne (PFDS + cinématique) dans les situations suivantes :

\begin{itemize}
    \item \textbf{Forces non constantes} (ressort, gravitation universelle, frottements fluides $f = -kv$) : La PFDS donne une équation différentielle difficile ; le théorème donne directement la relation $v(x)$ via l'intégration du travail.
    \item \textbf{Problème de « vitesse en fonction de la position »} : On ne cherche pas $x(t)$ ou $v(t)$, mais $v_B$ connaissant $v_A$ et la géométrie (hauteur, longueur, angle).
    \item \textbf{Forces de liaison « travail nul »} : Réaction normale, tension de fil inextensible (si poulie fixe), réaction d'un axe fixe. Elles disparaissent naturellement du bilan énergétique, simplifiant le problème.
    \item \textbf{Systèmes complexes} : Plusieurs solides liés (poulie, engrenages) : on applique le théorème à l'\textbf{ensemble} du système (somme des $E_c$ = somme des travaux des forces \textbf{extérieures} seulement ; les actions intérieures s'annulent si pas de frottement interne).
\end{itemize}

\begin{savaistu}[Énergie et mobilité au Cameroun]
Le théorème de l'énergie cinétique est au cœur du calcul de la \textbf{consommation énergétique des transports}.
\begin{itemize}
    \item \textbf{Mototaxis (benskin)} : Pour monter une côte (ex: montée vers \textbf{Nkolbisson} ou \textbf{Ekounou} à Yaoundé), le moteur fournit un travail moteur $W_{\text{moteur}} > 0$ pour augmenter $E_c$ et vaincre le travail résistant du poids ($W_P < 0$) et des frottements ($W_f < 0$). $\Delta E_c = W_{\text{moteur}} + W_P + W_f$. La consommation d'essence est proportionnelle à l'énergie fournie par le moteur (rendement $\approx 25-30\%$).
    \item \textbf{Freinage régénératif (bus électriques, future mobilité urbaine)} : Au freinage, $\Delta E_c < 0$. Au lieu de dissiper cette énergie en chaleur (freins classiques, $W_{\text{frein}} < 0$), un moteur-générateur récupère $|\Delta E_c|$ pour recharger la batterie (travail moteur négatif sur le volant, mais énergie stockée).
    \item \textbf{Barrages hydroélectriques (Lom Pangar, Nachtigal, Edéa)} : L'eau chute de hauteur $h$. Son énergie potentielle se convertit en énergie cinétique ($mg h = \frac{1}{2}mv^2$ en l'absence de pertes). La turbine récupère cette énergie cinétique (travail résistant sur l'eau, travail moteur sur l'alternateur).
\end{itemize}
Comprendre le théorème, c'est comprendre comment on \textbf{paye} (en carburant, en kWh) chaque mètre parcouru et chaque mètre/seconde gagné.
\end{savaistu}

\courssub{Synthèse : Énergie cinétique et Travail — Le duo inséparable}

\begin{aretenir}[Tableau récapitulatif pour le Probatoire]
\begin{center}
\begin{tabularx}{\linewidth}{|l|X|X|}\hline
\textbf{Situation} & \textbf{Énergie cinétique $E_c$} & \textbf{Travail $\sum W$ (forces extérieures)} \\ \hline
Translation (masse $m$, vitesse $\vec{v}_G$) & $E_c = \frac{1}{2} m v_G^2$ & $\sum \vec{F}_i \cdot \overrightarrow{AB}$ (ou $\int \vec{F}\cdot d\vec{l}$) \\ \hline
Rotation axe fixe $\Delta$ (moment $J_\Delta$, $\omega$) & $E_c = \frac{1}{2} J_\Delta \omega^2$ & $\sum \mathcal{M}_{\Delta,i} \Delta\theta$ (ou $\int \mathcal{M}_\Delta d\theta$) \\ \hline
Combiné (translation $G$ + rotation autour de $G$) & $E_c = \frac{1}{2} m v_G^2 + \frac{1}{2} J_G \omega^2$ & $\sum W_{\text{forces}} + \sum W_{\text{moments}/G}$ \\ \hline
\textbf{Théorème (toujours valide en référentiel galiléen)} & \multicolumn{2}{c|}{\textbf{$\Delta E_c = E_c(B) - E_c(A) = \sum W_{AB}$}} \\ \hline
\end{tabularx}
\end{center}
\end{aretenir}

\begin{exercice}[Entraînement : Ascenseur et contrepoids]
Un ascenseur de masse $m_c = \SI{800}{kg}$ (cabine + charge) est relié par un câble inextensible passant sur une poulie massive (moment d'inertie $J = \SI{50}{kg.m^2}$, rayon $R = \SI{0,40}{m}$) à un contrepoids de masse $m_p = \SI{600}{kg}$. Le système est lâché au repos. On néglige les frottements d'axe et l'élasticité du câble.
\begin{enumerate}
    \item Faire le bilan des forces extérieures sur le système \{cabine + câble + poulie + contrepoids\}.
    \item Exprimer les travaux de ces forces lorsque la cabine descend de $h = \SI{10}{m}$.
    \item En déduire la vitesse $v$ de la cabine après cette descente de $\SI{10}{m}$.
    \item Calculer la vitesse angulaire $\omega$ de la poulie à cet instant.
\end{enumerate}
\textit{Indication : Le câble ne glisse pas sur la poulie $\Rightarrow v = \omega R$. Les tensions sont des forces intérieures au système global (travail nul total car câble inextensible). Seuls les poids travaillent.}
\end{exercice}

\begin{corrige}[Exercice : Ascenseur et contrepoids]
\begin{enumerate}
    \item \textbf{Système} : \{Cabine + Contrepoids + Poulie + Câble\}. \textbf{Référentiel} : Terrestre (galiléen). Forces extérieures : $\vec{P}_c$, $\vec{P}_p$, Réactions d'axe sur poulie (travail nul), Poids de la poulie (travail nul car centre fixe).
    \item \textbf{Travaux} (descente cabine $h=\SI{10}{m}$ $\Rightarrow$ montée contrepoids $h=\SI{10}{m}$, rotation poulie $\Delta\theta = h/R$) :
    \begin{align*}
        W(P_c) &= +m_c g h = 800 \times 9,8 \times 10 = \SI{78400}{J} \\
        W(P_p) &= -m_p g h = -600 \times 9,8 \times 10 = \SI{-58800}{J} \\
        W(\text{axe}) &= 0 \quad (\text{point d'application fixe}) \\
        \sum W &= \SI{19600}{J}
    \end{align*}
    \item \textbf{Énergies cinétiques finales} (initialement nulles) :
    \begin{align*}
        E_c(\text{cabine}) &= \frac{1}{2} m_c v^2 \\
        E_c(\text{contrepoids}) &= \frac{1}{2} m_p v^2 \quad (v \text{ même norme}) \\
        E_c(\text{poulie}) &= \frac{1}{2} J \omega^2 = \frac{1}{2} J (v/R)^2 \\
        \Delta E_c &= \frac{1}{2} \left( m_c + m_p + \frac{J}{R^2} \right) v^2
    \end{align*}
    $\frac{J}{R^2} = \frac{50}{0,16} = \SI{312,5}{kg}$. Masse équivalente totale $= 800 + 600 + 312,5 = \SI{1712,5}{kg}$.
    \item \textbf{Théorème} : $\frac{1}{2} \times 1712,5 \times v^2 = 19600$.
    $v^2 = \frac{39200}{1712,5} \approx 22,89 \Rightarrow v \approx \SI{4,78}{m/s}$.
    $\omega = v/R = 4,78 / 0,40 \approx \SI{11,95}{rad/s}$.
\end{enumerate}
\end{corrige}

\coursec{Applications directes du théorème}

Dans la section précédente, nous avons établi le théorème de l'énergie cinétique : la variation d'énergie cinétique d'un solide entre deux instants est égale à la somme des travaux des forces appliquées au solide pendant ce déplacement. Ce théorème est l'un des outils les plus puissants de la mécanique de Première, et c'est précisément lui qui permet de répondre à la question de notre famille de situations : \emph{quel est le coût énergétique de la mobilité ?} Combien de mètres faut-il à un bus lancé à \SI{90}{km/h} pour s'arrêter ? Quelle énergie faut-il dépenser pour lancer une mototaxi de \SI{0}{} à \SI{60}{km/h} ? Nous allons maintenant appliquer le théorème à six situations concrètes et classiques, en dégageant à chaque fois la démarche et les formules à retenir.

\begin{methode}[Quand penser au théorème de l'énergie cinétique ?]
Le théorème de l'énergie cinétique est l'outil à privilégier dès que le problème fait intervenir des \cle{vitesses} et des \cle{distances} (ou des hauteurs, des angles) \textbf{sans que le temps n'intervienne}. Démarche type :
\begin{enumerate}
\item Définir le \cle{système} étudié et l'état initial (A) et l'état final (B) ;
\item Faire le \cle{bilan des forces} appliquées au solide et représenter les vecteurs sur un schéma ;
\item Calculer le \cle{travail} de chaque force entre A et B (attention aux forces perpendiculaires au déplacement, dont le travail est nul) ;
\item Écrire le théorème : $E_{c}(B) - E_{c}(A) = \sum W_{AB}(\vec{F})$ ;
\item Isoler la grandeur cherchée (vitesse, distance, hauteur, angle\dots) et faire l'application numérique en unités SI.
\end{enumerate}
\end{methode}

\courssub{Application 1 : distance de freinage d'un véhicule sur route horizontale}

\textbf{Situation.} Un bus de masse $m$ roule à la vitesse $v$ sur une route horizontale (par exemple l'axe lourd Yaoundé--Douala). Le conducteur freine : les freins exercent sur le véhicule une force de freinage $\vec{f}$ supposée constante, opposée au mouvement. Quelle distance $d$ le bus parcourt-il avant l'arrêt ?

\textbf{Bilan des forces.} Le bus est soumis à trois forces : son poids $\vec{P}$, la réaction normale de la route $\vec{R}_{N}$ et la force de freinage $\vec{f}$. Sur route horizontale, $\vec{P}$ et $\vec{R}_{N}$ sont perpendiculaires au déplacement : leurs travaux sont nuls. Seule $\vec{f}$ travaille, avec un travail résistant :
\[ W(\vec{f}) = -f \times d. \]

\textbf{Application du théorème.} Entre l'instant où le freinage commence (vitesse $v$) et l'arrêt (vitesse nulle) :
\[ E_{c}(B) - E_{c}(A) = W(\vec{f}) \quad\Longrightarrow\quad 0 - \frac{1}{2}mv^{2} = -f\,d. \]
D'où la formule fondamentale de la distance de freinage :
\begin{propriete}[Distance de freinage]
Sur route horizontale, avec une force de freinage constante $f$, la distance d'arrêt d'un véhicule de masse $m$ lancé à la vitesse $v$ est :
\[ d = \frac{m\,v^{2}}{2f}. \]
Vérification dimensionnelle : $\dfrac{\mathrm{kg}\cdot\mathrm{m}^{2}\cdot\mathrm{s}^{-2}}{\mathrm{kg}\cdot\mathrm{m}\cdot\mathrm{s}^{-2}} = \mathrm{m}$. La formule est homogène.
\end{propriete}

\begin{attention}[La distance de freinage varie comme le carré de la vitesse]
C'est la conséquence la plus importante pour la sécurité routière : $d$ est proportionnelle à $v^{2}$, pas à $v$. \textbf{Doubler la vitesse multiplie la distance de freinage par quatre.} C'est pourquoi les limitations de vitesse à l'entrée des villes camerounaises ne sont pas un détail : un véhicule qui passe de \SI{90}{km/h} à \SI{45}{km/h} divise sa distance de freinage par quatre. Ne jamais écrire que $d$ est proportionnelle à $v$ : c'est l'erreur classique du Probatoire.
\end{attention}

\begin{exemplebox}[Freinage d'un bus de 8 tonnes]
Un bus de masse $m = \SI{8000}{kg}$ roule à $v = \SI{72}{km/h} = \SI{20}{m/s}$. La force de freinage maximale vaut $f = \SI{40}{kN} = \SI{40000}{N}$.
\[ d = \frac{m v^{2}}{2f} = \frac{8000 \times 20^{2}}{2 \times 40000} = \frac{8000 \times 400}{80000} = \SI{40}{m}. \]
À $v' = \SI{36}{km/h} = \SI{10}{m/s}$ (vitesse deux fois plus faible) :
\[ d' = \frac{8000 \times 10^{2}}{2 \times 40000} = \SI{10}{m}. \]
La vitesse a été divisée par 2, la distance de freinage par 4 : on passe de \SI{40}{m} à \SI{10}{m}. En ville, c'est souvent la différence entre un arrêt sans dommage et un accident.
\end{exemplebox}

\begin{popfigure}
\begin{center}
\begin{tikzpicture}
\begin{axis}[
width=0.85\linewidth, height=6.5cm,
xlabel={Vitesse $v$ (km/h)},
ylabel={Distance de freinage $d$ (m)},
xmin=20, xmax=130, ymin=0, ymax=140,
grid=both, grid style={popline!40},
xtick={20,40,60,80,100,120},
label style={font=\small},
tick label style={font=\small},
]
\addplot[popcurve,domain=20:130,samples=100]{0.007716*x^2};
\addplot[mark=*,only marks,poporange] coordinates {(72,40) (36,10)};
\node[font=\scriptsize,popdark,anchor=west] at (axis cs:75,40) {72 km/h : $d = 40$ m};
\node[font=\scriptsize,popdark,anchor=west] at (axis cs:40,10) {36 km/h : $d = 10$ m};
\end{axis}
\end{tikzpicture}
\end{center}
\legende{Distance de freinage en fonction de la vitesse pour un bus de 8 t avec $f = \SI{40}{kN}$ (route sèche). La courbe est une parabole : $d$ croît comme $v^{2}$.}
\end{popfigure}

\begin{savaistu}[Les zones de freinage d'urgence en gravillons]
Sur les longues descentes du Cameroun, comme sur la colline de Nkolbisson à la sortie de Yaoundé ou sur certaines côtes de l'Ouest, on aménage des \cle{zones de freinage d'urgence} : des voies remplies de gravillons épais où un camion dont les freins ont lâché peut s'engager. Le principe est exactement celui de notre formule : les gravillons exercent une force de frottement $f$ considérable, et toute l'énergie cinétique $\frac{1}{2}mv^{2}$ du camion est convertie en travail des forces de frottement (et en échauffement). Plus la vitesse est élevée, plus le camion s'enfonce loin dans le lit de gravillons : c'est $d = mv^{2}/(2f)$ qui décide de la longueur de la zone à construire !
\end{savaistu}

\courssub{Application 2 : vitesse en bas d'un plan incliné avec frottements}

\textbf{Situation.} Un solide de masse $m$ est lâché sans vitesse initiale du sommet A d'un plan incliné d'angle $\alpha$ sur l'horizontale, de longueur $L$. Il est freiné par une force de frottement $\vec{f}$ constante, parallèle au plan. Quelle est sa vitesse $v$ au point B, en bas du plan ?

\textbf{Bilan des forces.} Le solide est soumis à son poids $\vec{P} = m\vec{g}$, à la réaction normale $\vec{R}_{N}$ (perpendiculaire au plan, travail nul) et aux frottements $\vec{f}$ (opposés au mouvement, travail $-fL$). Le travail du poids entre A et B ne dépend que du dénivelé $h = L\sin\alpha$ :
\[ W_{AB}(\vec{P}) = m\,g\,h = m\,g\,L\sin\alpha. \]

\textbf{Application du théorème.} Le solide part sans vitesse :
\[ \frac{1}{2}mv^{2} - 0 = m\,g\,L\sin\alpha - f\,L. \]
En divisant par $m$ et en multipliant par 2 :
\begin{propriete}[Vitesse en bas d'un plan incliné avec frottements]
\[ v^{2} = 2\,g\,L\left(\sin\alpha - \frac{f}{m\,g}\right). \]
Sans frottements ($f = 0$), on retrouve $v^{2} = 2gL\sin\alpha = 2gh$ : la vitesse en bas ne dépend que du dénivelé $h$, pas de la longueur du plan ni de la masse.
\end{propriete}

\begin{attention}[Condition de mouvement]
La formule n'a de sens que si $\sin\alpha > \dfrac{f}{mg}$, c'est-à-dire si la composante du poids le long du plan ($mg\sin\alpha$) est supérieure aux frottements. Sinon, le solide ne démarre pas (ou s'arrête) : le radicande serait négatif, ce qui est physiquement impossible. Si le calcul aboutit à $v^{2} < 0$, conclure : \og le solide ne parvient pas en bas \fg.
\end{attention}

\begin{popfigure}
\begin{center}
\begin{tikzpicture}[scale=1.05, >=Stealth]
\coordinate (A) at (0,0);
\coordinate (B) at (6,0);
\coordinate (C) at (6,3);
\draw[thick, popdark] (A) -- (B) -- (C) -- cycle;
\draw[pattern=north east lines, pattern color=popmuted] (A) -- (B) -- ++(0,-0.15) -- ++(-6,0) -- cycle;
\draw[popdark] (4.6,0) arc (0:26.565:1.4) node[midway, right, font=\small] {$\alpha$};
\draw[fill=popblueL, draw=popblue, thick] (4.35,2.175) rectangle ++(0.55,0.4);
\node[font=\small, popdark] at (4.62,3.0) {solide de masse $m$};
\draw[->, very thick, popink] (4.62,2.37) -- ++(0,-1.4) node[below, font=\small] {$\vec{P}$};
\draw[->, very thick, poporange] (4.62,2.57) -- ++(-0.9,0.45) node[above left, font=\small] {$\vec{R}_N$};
\draw[->, very thick, poppurple] (4.35,2.37) -- ++(-1.2,-0.6) node[below left, font=\small] {$\vec{f}$};
\draw[->, very thick, popgreen] (4.9,2.17) -- ++(1.2,0.6) node[above right, font=\small] {$\vec{v}$};
\draw[<->, popmuted] (0.15,0.12) -- (5.85,2.97) node[midway, above left, font=\small] {$L$};
\draw[<->, popmuted] (6.25,0) -- (6.25,3) node[midway, right, font=\small] {$h = L\sin\alpha$};
\node[font=\small, popdark] at (0,-0.35) {B (bas)};
\node[font=\small, popdark] at (6,3.3) {A (haut)};
\end{tikzpicture}
\end{center}
\legende{Solide glissant sur un plan incliné : cotation de la longueur $L$, de l'angle $\alpha$ et du dénivelé $h = L\sin\alpha$ ; bilan des forces appliquées au solide.}
\end{popfigure}

\begin{exoresolu}[Descente d'une caisse sur un plan incliné]
\textbf{Énoncé.} Une caisse de masse $m = \SI{50}{kg}$ est lâchée sans vitesse du sommet d'un plan incliné de longueur $L = \SI{10}{m}$, incliné de $\alpha = 30^{\circ}$. Les frottements équivalent à une force constante $f = \SI{100}{N}$. On prend $g = \SI{10}{N/kg}$. Calculer la vitesse de la caisse en bas du plan.
\tcblower
\textbf{Solution.} Vérifions d'abord que la caisse glisse : $mg\sin\alpha = 50 \times 10 \times 0{,}5 = \SI{250}{N} > f = \SI{100}{N}$. La condition est remplie.
Théorème de l'énergie cinétique entre le sommet A (vitesse nulle) et le bas B :
\[ \frac{1}{2}mv^{2} = mgL\sin\alpha - fL. \]
Application numérique :
\begin{align*}
\frac{1}{2} \times 50 \times v^{2} &= 50 \times 10 \times 10 \times 0{,}5 - 100 \times 10 \\
25\,v^{2} &= 2500 - 1000 = 1500 \\
v^{2} &= 60 \quad\Longrightarrow\quad v = \sqrt{60} \approx \SI{7,7}{m/s}.
\end{align*}
Sans frottements, on aurait trouvé $v = \sqrt{2 \times 10 \times 10 \times 0{,}5} = \SI{10}{m/s}$ : les frottements ont bien fait perdre de la vitesse, donc de l'énergie cinétique.
\end{exoresolu}

\courssub{Application 3 : hauteur maximale d'un projectile lancé verticalement}

\textbf{Situation.} Un enfant lance une balle verticalement vers le haut avec la vitesse initiale $v_{0}$. On néglige la résistance de l'air. Jusqu'à quelle hauteur $h$ la balle monte-t-elle ?

\textbf{Application du théorème.} La seule force qui travaille est le poids (travail résistant pendant la montée : $W(\vec{P}) = -mgh$). Au point le plus haut, la vitesse s'annule :
\[ 0 - \frac{1}{2}mv_{0}^{2} = -m\,g\,h \quad\Longrightarrow\quad \boxed{h = \frac{v_{0}^{2}}{2g}}. \]

\begin{exemplebox}[Lancer d'une balle]
Avec $v_{0} = \SI{14}{m/s}$ et $g = \SI{10}{N/kg}$ :
\[ h = \frac{14^{2}}{2 \times 10} = \frac{196}{20} = \SI{9,8}{m}. \]
La balle monte à près de \SI{10}{m}, soit la hauteur d'un immeuble de trois étages. Remarque : la masse de la balle n'intervient pas, car elle se simplifie dans l'équation.
\end{exemplebox}

\courssub{Application 4 : travail nécessaire pour accélérer un véhicule}

\textbf{Situation.} Quelle énergie faut-il fournir à une mototaxi de masse $m$ (avec son conducteur) pour passer de l'arrêt à la vitesse $v$ sur route horizontale, en négligeant les frottements ?

\textbf{Application du théorème.} Le poids et la réaction de la route ne travaillent pas. Le travail $W$ de la force motrice se convertit intégralement en énergie cinétique :
\[ W = \frac{1}{2}mv^{2} - 0 = \frac{1}{2}mv^{2}. \]

\begin{exemplebox}[Coût énergétique d'une mototaxi]
Pour $m = \SI{250}{kg}$ (moto + conducteur + passager) et $v = \SI{60}{km/h} \approx \SI{16,7}{m/s}$ :
\[ W = \frac{1}{2} \times 250 \times 16{,}7^{2} \approx \SI{3,5e4}{J} = \SI{35}{kJ}. \]
À chaque arrêt (feu, ralentisseur, embouteillage), cette énergie est perdue dans les freins et il faut la dépenser à nouveau en carburant. C'est pourquoi la conduite en ville, avec ses arrêts répétés, consomme beaucoup plus que la conduite à vitesse stable : le \cle{coût énergétique de la mobilité} est directement lié à l'énergie cinétique que l'on crée puis que l'on détruit.
\end{exemplebox}

\courssub{Application 5 : arrêt d'un volant en rotation freiné par un couple constant}

\textbf{Situation.} Un volant de moment d'inertie $J_{\Delta}$ tourne à la vitesse angulaire $\omega_{0}$. On coupe le moteur et un couple de freinage résistant de moment $\mathcal{M}$ constant s'applique. De quel angle $\theta$ le volant tourne-t-il avant de s'arrêter ?

\textbf{Application du théorème.} Pour un solide en rotation, l'énergie cinétique est $E_{c} = \frac{1}{2}J_{\Delta}\omega^{2}$ et le travail d'un couple constant sur un angle $\theta$ est $W = \mathcal{M}\,\theta$ (ici résistant, donc négatif) :
\[ 0 - \frac{1}{2}J_{\Delta}\omega_{0}^{2} = -\mathcal{M}\,\theta \quad\Longrightarrow\quad \theta = \frac{J_{\Delta}\,\omega_{0}^{2}}{2\mathcal{M}}. \]

\begin{exemplebox}[Volant d'une machine de menuiserie]
Un volant de $J_{\Delta} = \SI{2}{kg.m^{2}}$ tourne à $\omega_{0} = \SI{20}{rad/s}$. Le couple de freinage vaut $\mathcal{M} = \SI{10}{N.m}$ :
\[ \theta = \frac{2 \times 20^{2}}{2 \times 10} = \SI{40}{rad} \approx 6{,}4 \text{ tours}. \]
Le volant effectue encore plus de six tours avant l'arrêt : d'où les protections obligatoires sur les machines tournantes des ateliers.
\end{exemplebox}

\courssub{Application 6 : solide en roulement sans glisser (translation + rotation)}

\textbf{Situation.} Un cylindre plein de masse $m$, de rayon $R$ et de moment d'inertie $J_{\Delta} = \frac{1}{2}mR^{2}$ (axe de symétrie) roule sans glisser sur un plan incliné d'angle $\alpha$, de longueur $L$, depuis le repos. Déterminer sa vitesse $v$ au bas du plan.

\textbf{Bilan des forces et travaux.} Forces : poids $\vec{P}$, réaction normale $\vec{R}_N$, frottement de roulement $\vec{f}$ (statique). Le point de contact est instantanément au repos : le travail du frottement est \textbf{nul}. Seuls le poids travaille.
\[ W_{AB}(\vec{P}) = mgL\sin\alpha = mgh. \]

\textbf{Énergie cinétique finale.} Mouvement combiné : translation du centre d'inertie ($v$) et rotation autour de lui ($\omega$). Condition de non-glissement : $v = \omega R$.
\[ E_{c}(B) = \frac{1}{2}mv^{2} + \frac{1}{2}J_{\Delta}\omega^{2} = \frac{1}{2}mv^{2} + \frac{1}{2}\left(\frac{1}{2}mR^{2}\right)\left(\frac{v}{R}\right)^{2} = \frac{1}{2}mv^{2} + \frac{1}{4}mv^{2} = \frac{3}{4}mv^{2}. \]

\textbf{Théorème de l'énergie cinétique.}
\[ \frac{3}{4}mv^{2} = mgL\sin\alpha \quad\Longrightarrow\quad \boxed{v^{2} = \frac{4}{3}gL\sin\alpha = \frac{4}{3}gh}. \]

\begin{propriete}[Vitesse d'un cylindre roulant sans glisser]
Pour un cylindre plein ($J = \frac{1}{2}mR^2$) partant du repos sur un plan incliné de hauteur $h$ :
\[ v = \sqrt{\frac{4}{3}gh}. \]
Il est plus lent qu'un bloc qui glisserait sans frottement ($v = \sqrt{2gh}$) car une partie de l'énergie potentielle est convertie en énergie cinétique de rotation.
\end{propriete}

\begin{attention}[Frottement de roulement et travail]
Le frottement de roulement (statique) ne travaille \textbf{pas} car son point d'application ne se déplace pas (vitesse nulle instantanée). Il est indispensable pour \emph{transmettre} l'énergie de translation en rotation, mais ne \emph{dissipe} pas d'énergie. Ne pas confondre avec le frottement de glissement (cinétique) qui travaille toujours négativement.
\end{attention}

\begin{attention}[Ce que le théorème ne dit pas]
Le théorème de l'énergie cinétique donne des relations entre des \cle{grandeurs globales} : vitesses initiale et finale, distance, hauteur, angle. Il ne donne \textbf{ni le temps} du parcours, \textbf{ni la loi horaire} $x(t)$ du mouvement. Pour obtenir une durée ou une équation horaire, il faut revenir à la relation fondamentale de la dynamique. Autre piège : ne pas oublier que seules les forces dont le point d'application se déplace et qui ne sont pas perpendiculaires au déplacement travaillent ; sur un plan incliné, la réaction normale ne travaille jamais.
\end{attention}

\begin{exercice}[Pour s'entraîner]
Un camion de masse $m = \SI{12}{t}$ descend une côte de longueur $L = \SI{200}{m}$ inclinée de $\alpha = 6^{\circ}$ ($\sin 6^{\circ} \approx 0{,}105$), moteur coupé, en partant à $v_{0} = \SI{36}{km/h}$. Les frottements équivalent à $f = \SI{6000}{N}$. Calculer sa vitesse en bas de la côte ($g = \SI{10}{N/kg}$). Le conducteur freine ensuite sur l'horizontale avec $f' = \SI{30000}{N}$ : quelle distance lui faut-il pour s'arrêter ?
\end{exercice}

\begin{corrige}[Exercice : camion en descente]
\textbf{1. Vitesse en bas de la côte.}
Système : camion. États : A (haut de la côte, $v_A = \SI{10}{m/s}$), B (bas de la côte, $v_B = ?$).
Forces travaillant : Poids ($W_P = mgL\sin\alpha$), Frottements ($W_f = -fL$).
Théorème :
\[ \frac{1}{2}m v_B^2 - \frac{1}{2}m v_A^2 = mgL\sin\alpha - fL. \]
Numérique ($m = 12000\,\mathrm{kg}$, $g=10$, $L=200$, $\sin\alpha=0.105$, $f=6000$, $v_A=10$) :
\begin{align*}
\frac{1}{2}\times 12000 (v_B^2 - 100) &= 12000\times 10 \times 200 \times 0.105 - 6000 \times 200 \\
6000 (v_B^2 - 100) &= 25\,200\,000 - 1\,200\,000 = 24\,000\,000 \\
v_B^2 - 100 &= 4000 \quad\Rightarrow\quad v_B^2 = 4100 \\
v_B &\approx \SI{64,0}{m/s} \quad (\text{soit } \approx \SI{230}{km/h}).
\end{align*}
\textbf{2. Distance de freinage sur l'horizontale.}
État initial : $v_B \approx 64\,\mathrm{m/s}$, État final : arrêt ($v=0$). Force de freinage $f' = 30000\,\mathrm{N}$.
\[ 0 - \frac{1}{2}m v_B^2 = -f' d \quad\Rightarrow\quad d = \frac{m v_B^2}{2f'} = \frac{12000 \times 4100}{2 \times 30000} = \frac{49\,200\,000}{60\,000} = \SI{820}{m}. \]
Le camion met \SI{820}{m} pour s'arrêter : l'importance de l'entretien des freins en zone montagneuse est vitale.
\end{corrige}

\coursec{Méthode de résolution et synthèse}

Dans la section précédente, nous avons appliqué le théorème de l'énergie cinétique à des situations concrètes : distance de freinage d'une mototaxi, vitesse acquise sur une pente, travail d'une force de traction. Tu as sans doute remarqué que, malgré la diversité des énoncés, la démarche était toujours la même : choisir un système, faire le bilan des forces, écrire les travaux, appliquer $\Delta E_c = \sum W$. Cette régularité n'est pas un hasard : c'est la force de la méthode énergétique. Cette dernière section te donne la \cle{fiche méthode} complète à recopier et à utiliser systématiquement au Probatoire, les critères pour choisir entre la relation fondamentale de la dynamique (RFD) et le théorème de l'énergie cinétique (TEC), un tableau de synthèse de toutes les expressions du chapitre, et un complément sur la puissance, très apprécié des correcteurs.

\courssub{Fiche méthode : résoudre un problème par l'énergie cinétique}

L'intuition d'abord : le théorème de l'énergie cinétique est une \emph{comptabilité}. Le solide « gagne » de l'énergie cinétique quand les forces qui le poussent dans le sens du mouvement travaillent positivement (travail moteur), et il en « perd » quand des forces s'opposent au mouvement (travail résistant). Pour faire cette comptabilité correctement, il faut une procédure rigoureuse, toujours la même.

\begin{methode}[Fiche de résolution type en 5 étapes]
\begin{enumerate}
\item \textbf{Choisir le système et le référentiel.} Préciser le solide étudié (par exemple « la moto et son conducteur, assimilés à un solide de masse $m$ ») et le référentiel galiléen (le plus souvent le référentiel terrestre). Identifier l'état initial A (vitesse $v_A$) et l'état final B (vitesse $v_B$).
\item \textbf{Faire le bilan exhaustif des forces extérieures} appliquées au système et les représenter sur un schéma : poids $\vec{P}$, réaction du support $\vec{R}$ (normale $\vec{R}_N$ et frottements $\vec{f}$), force motrice $\vec{F}$, tension d'un câble, etc. Ne rien oublier, ne rien inventer.
\item \textbf{Exprimer le travail de chaque force} entre A et B :
      \begin{itemize}
      \item Poids : $W_{AB}(\vec{P}) = -mg(z_B - z_A) = -mgh$, où $h = z_B - z_A$ est la dénivellation algébrique (positive en montée, négative en descente).
      \item Frottements constants $\vec{f}$ opposés au déplacement : $W_{AB}(\vec{f}) = -f \times AB$.
      \item Force motrice $\vec{F}$ constante formant un angle $\alpha$ avec le déplacement : $W_{AB}(\vec{F}) = F \times AB \times \cos\alpha$.
      \item Réaction normale $\vec{R}_N$ : $W_{AB}(\vec{R}_N) = 0$ car $\vec{R}_N \perp$ déplacement.
      \end{itemize}
\item \textbf{Appliquer le théorème de l'énergie cinétique} :
\[ \frac{1}{2}mv_B^2 - \frac{1}{2}mv_A^2 = \sum W_{AB}(\vec{F}_{\text{ext}}) \]
\item \textbf{Extraire l'inconnue, calculer et vérifier} : convertir toutes les données en unités SI (vitesses en m/s, masses en kg, distances en m), effectuer l'application numérique, vérifier l'homogénéité du résultat (joules, m/s, watts) et sa plausibilité physique.
\end{enumerate}
\end{methode}

\begin{attention}[Homogénéité : le réflexe qui sauve des points]
Avant d'encadrer un résultat, vérifie toujours son unité : une énergie s'exprime en \textbf{joules} ($1~\text{J} = 1~\text{kg}\cdot\text{m}^2\cdot\text{s}^{-2}$), une vitesse en \textbf{m/s}, une puissance en \textbf{watts}. Si tu trouves une vitesse en joules ou une \emph{variation d'énergie} négative alors que le solide accélère sous l'effet d'un moteur, il y a une erreur de signe ou de conversion. L'analyse dimensionnelle rapide de $\tfrac{1}{2}mv^2$ donne bien des $\text{kg}\cdot\text{m}^2\cdot\text{s}^{-2}$ : c'est le contrôle minimal exigible.
\end{attention}

\courssub{Méthode énergétique ou relation fondamentale de la dynamique ?}

Deux outils s'offrent à toi en mécanique : la RFD, $\sum \vec{F}_{\text{ext}} = m\vec{a}$, qui est \emph{vectorielle} et donne accès à l'accélération, puis aux équations horaires $x(t)$ et $v(t)$ ; et le TEC, qui est \emph{scalaire} et relie directement les vitesses aux travaux, sans passer par le temps. Le choix dépend de la question posée et des données disponibles.

\begin{itemize}
\item Si l'énoncé demande une \cle{accélération}, une \cle{durée}, une \cle{loi horaire}, ou si le mouvement n'est pas rectiligne uniformément varié, la RFD est l'outil adapté.
\item Si l'énoncé relie une \cle{vitesse} à une \cle{distance} ou à une \cle{altitude} (distance de freinage, vitesse en bas d'une pente, force de traction sur un parcours), le TEC est beaucoup plus rapide : une seule équation scalaire suffit.
\item Si les forces ne sont pas constantes ou si la trajectoire est courbe, le TEC garde toute sa puissance là où la RFD devient lourde.
\end{itemize}

\begin{popfigure}
\begin{center}
\begin{tikzpicture}[
  font=\small,
  startstop/.style={rectangle, rounded corners, draw=popdark, fill=popblueL, text width=4.6cm, align=center, minimum height=1cm},
  decision/.style={diamond, draw=popdark, fill=popgoldL, text width=3.4cm, align=center, aspect=2.2, inner sep=1pt},
  process/.style={rectangle, draw=popdark, fill=popgreenL, text width=4.4cm, align=center, minimum height=1cm},
  fleche/.style={-{Stealth[length=2.5mm]}, thick, popdark}]
\node[startstop] (depart) {Lire l'énoncé : identifier la question posée et les données};
\node[decision, below=0.7cm of depart] (q1) {Cherche-t-on $a$, $t$ ou une loi horaire ?};
\node[process, below left=0.9cm and 0.2cm of q1] (rfd) {\textbf{RFD} : $\sum \vec{F}_{\text{ext}} = m\vec{a}$, puis intégration pour $v(t)$ et $x(t)$};
\node[decision, below right=0.9cm and 0.2cm of q1] (q2) {Relie-t-on vitesse et distance (ou altitude) ?};
\node[process, below=0.7cm of q2] (tec) {\textbf{TEC} : $\tfrac{1}{2}mv_B^2 - \tfrac{1}{2}mv_A^2 = \sum W_{AB}$};
\node[process, below=0.7cm of rfd] (verif1) {Bilan des forces, schéma, calcul, vérification des unités};
\draw[fleche] (depart) -- (q1);
\draw[fleche] (q1.west) -| node[pos=0.25, above left]{oui} (rfd.north);
\draw[fleche] (q1.east) -| node[pos=0.25, above right]{non} (q2.north);
\draw[fleche] (q2) -- node[right]{oui} (tec);
\draw[fleche] (q2.west) -| node[pos=0.3, above]{non : combiner les deux} (verif1.east);
\draw[fleche] (rfd) -- (verif1);
\draw[fleche] (tec.south) |- ($(verif1.east)+(0,-0.35)$) -- (verif1.east);
\end{tikzpicture}
\end{center}
\legende{Organigramme de décision : quel outil de mécanique utiliser selon les données du problème.}
\end{popfigure}

\courssub{Bilan des expressions de l'énergie cinétique}

Rassemblons les trois expressions rencontrées dans cette leçon. Pour un solide de masse $m$ dont le centre d'inertie G a la vitesse $v_G$ :

\begin{aretenir}[Les trois expressions de l'énergie cinétique]
\begin{itemize}
\item \textbf{Translation} (tous les points ont la même vitesse) : $E_c = \dfrac{1}{2}mv_G^2$.
\item \textbf{Rotation} autour d'un axe fixe ($\Delta$) à la vitesse angulaire $\omega$ : $E_c = \dfrac{1}{2}J_{\Delta}\omega^2$, où $J_{\Delta}$ est le moment d'inertie du solide par rapport à l'axe (en $\text{kg}\cdot\text{m}^2$).
\item \textbf{Mouvement combiné} (translation de G + rotation autour d'un axe passant par G) : $E_c = \dfrac{1}{2}mv_G^2 + \dfrac{1}{2}J_G\omega^2$.
\end{itemize}
Dans tous les cas, $E_c$ est une grandeur scalaire, toujours positive, exprimée en joules.
\end{aretenir}

\begin{exemplebox}[Roue de mototaxi en mouvement combiné]
Une roue de moto de masse $m = 12~\text{kg}$, assimilée à un cerceau de rayon $r = 0{,}30~\text{m}$ (donc $J_G = mr^2$), roule sans glisser à $v = 54~\text{km/h} = 15~\text{m/s}$. Alors $\omega = v/r = 50~\text{rad/s}$ et :
\[ E_c = \frac{1}{2}mv^2 + \frac{1}{2}mr^2\omega^2 = \frac{1}{2}\times 12 \times 15^2 + \frac{1}{2}\times 12 \times 0{,}30^2 \times 50^2 = 1350 + 1350 = 2700~\text{J}. \]
La rotation contribue ici autant que la translation : l'oublier reviendrait à sous-estimer l'énergie de moitié !
\end{exemplebox}

\courssub{Tableau de synthèse : énergies cinétiques et travaux usuels}

Le tableau suivant rassemble tout ce qu'il faut avoir en tête le jour de l'épreuve.

\begin{popfigure}
\begin{center}
\begin{tabular}{|l|c|c|}
\hline
\rowcolor{popblueL} \textbf{Grandeur} & \textbf{Expression} & \textbf{Remarque} \\
\hline
$E_c$ (translation) & $\dfrac{1}{2}mv^2$ & $v$ en m/s \\
\hline
$E_c$ (rotation) & $\dfrac{1}{2}J_{\Delta}\omega^2$ & $\omega$ en rad/s \\
\hline
$E_c$ (combiné) & $\dfrac{1}{2}mv_G^2 + \dfrac{1}{2}J_G\omega^2$ & roulement : $v_G = r\omega$ \\
\hline
$W(\vec{P})$ poids & $-mg\Delta z = -mgh$ & $h = z_B - z_A$ (algébrique) \\
\hline
$W(\vec{f})$ frottements & $-f \times AB$ & toujours résistant \\
\hline
$W(\vec{F})$ force motrice & $F \times AB \times \cos\alpha$ & moteur si $\alpha < 90^{\circ}$ \\
\hline
$W(\vec{R}_N)$ réaction normale & $0$ & $\vec{R}_N \perp$ déplacement \\
\hline
Théorème & $\dfrac{1}{2}mv_B^2 - \dfrac{1}{2}mv_A^2 = \sum W_{AB}$ & référentiel galiléen \\
\hline
\end{tabular}
\end{center}
\legende{Tableau récapitulatif des expressions de l'énergie cinétique et des travaux usuels.}
\end{popfigure}

\begin{attention}[Les quatre erreurs classiques]
\begin{enumerate}
\item \textbf{Confusion km/h et m/s} : $90~\text{km/h} = 25~\text{m/s}$ (diviser par $3{,}6$). Une vitesse laissée en km/h dans $\tfrac{1}{2}mv^2$ donne une énergie environ 13 fois trop grande ($(3{,}6)^2 \approx 13$) !
\item \textbf{Signe des travaux} : le poids résiste en montée ($W = -mgh$) et aide en descente ($W = +mgh$) ; les frottements fournissent toujours un travail négatif.
\item \textbf{Oubli d'une force} : fais systématiquement le schéma avec tous les vecteurs forces avant d'écrire le TEC. La réaction normale, elle, travaille nul : ne l'écris pas dans le bilan des travaux.
\item \textbf{Mauvais état initial/final} : identifie clairement A et B ; échanger $v_A$ et $v_B$ inverse le signe de $\Delta E_c$ et fausse toute la conclusion.
\end{enumerate}
\end{attention}

\courssub{Complément Probatoire : puissance moyenne et énergie cinétique}

Le théorème de l'énergie cinétique ne fait pas intervenir le temps. Pourtant, dans la vie réelle, la durée compte : un moteur qui fournit un travail donné en 5~s est plus « puissant » qu'un moteur qui le fournit en 50~s. D'où la notion de \cle{puissance}.

\begin{propriete}[Puissance moyenne — complément utile au Probatoire]
La \cle{puissance moyenne} développée par un ensemble de forces entre les instants $t_A$ et $t_B$ est le quotient du travail total fourni par la durée :
\[ P_m = \frac{\sum W_{AB}}{\Delta t} \qquad \text{avec} \quad \Delta t = t_B - t_A. \]
Elle s'exprime en \textbf{watts} ($1~\text{W} = 1~\text{J/s}$). Lorsque seules les forces motrices travaillent (frottements négligés, mouvement horizontal), le théorème de l'énergie cinétique donne $\sum W = \Delta E_c$, donc :
\[ P_m = \frac{\Delta E_c}{\Delta t}. \]
Cette relation n'est pas exigible en tant que telle, mais son utilisation correcte est valorisée au Probatoire.
\end{propriete}

Vérification dimensionnelle : $\dfrac{[\text{J}]}{[\text{s}]} = \dfrac{\text{kg}\cdot\text{m}^2\cdot\text{s}^{-2}}{\text{s}} = \text{kg}\cdot\text{m}^2\cdot\text{s}^{-3} = \text{W}$. La relation est homogène.

\begin{exoresolu}[Puissance minimale du moteur d'une moto]
Énoncé. Une moto et son conducteur, de masse totale $m = 250~\text{kg}$, partent du repos sur une route horizontale et atteignent $v = 90~\text{km/h}$ en $\Delta t = 10~\text{s}$. En négligeant les frottements, calculer la puissance moyenne minimale que doit développer le moteur.
\tcblower
Solution détaillée. Conversion : $v = 90/3{,}6 = 25~\text{m/s}$. Système : \{moto + conducteur\}, référentiel terrestre galiléen. Sur route horizontale, le poids et la réaction normale sont perpendiculaires au déplacement : ils ne travaillent pas. Seule la force motrice travaille. Le TEC entre l'état initial (repos) et l'état final donne :
\[ W_{\text{moteur}} = \Delta E_c = \frac{1}{2}mv^2 - 0 = \frac{1}{2}\times 250 \times 25^2 = \num{78125}~\text{J} \approx 7{,}8\times 10^4~\text{J}. \]
La puissance moyenne minimale est donc :
\[ P_m = \frac{\Delta E_c}{\Delta t} = \frac{\num{78125}}{10} \approx 7{,}8\times 10^3~\text{W} = 7{,}8~\text{kW}. \]
Interprétation : il s'agit d'un minimum, car en réalité les frottements de l'air et de la route consomment une partie du travail du moteur. Un moteur réel devra donc développer davantage que $7{,}8~\text{kW}$ pour réaliser cette accélération en 10~s.
\end{exoresolu}

\courssub{Exercice de synthèse}

\begin{exercice}[Montée d'une côte à Ngaoundéré]
Une mototaxi de masse totale $m = 200~\text{kg}$ aborde une côte rectiligne de longueur $L = 100~\text{m}$ avec la vitesse $v_A = 54~\text{km/h}$. La côte s'élève d'une dénivellation $h = 8{,}0~\text{m}$. Les frottements équivalent à une force constante $f = 150~\text{N}$. Le moteur exerce une force motrice constante $F = 400~\text{N}$ parallèle à la route. On prend $g = 10~\text{N/kg}$.
\begin{enumerate}
\item Faire le bilan des forces et calculer le travail de chacune entre le bas A et le haut B de la côte.
\item Déterminer la vitesse $v_B$ au sommet.
\item Calculer la puissance moyenne développée par le moteur pendant la montée.
\end{enumerate}
\end{exercice}

\begin{corrige}[Exercice : montée d'une côte à Ngaoundéré]
1. Conversion : $v_A = 54/3{,}6 = 15~\text{m/s}$. Forces : poids $\vec{P}$, réaction normale $\vec{R}_N$ (travail nul), frottements $\vec{f}$, force motrice $\vec{F}$.
\begin{align*}
W(\vec{P}) &= -mgh = -200\times 10 \times 8{,}0 = -\num{16000}~\text{J},\\
W(\vec{f}) &= -fL = -150 \times 100 = -\num{15000}~\text{J},\\
W(\vec{F}) &= FL = 400 \times 100 = +\num{40000}~\text{J}.
\end{align*}
Somme des travaux : $\sum W = \num{40000} - \num{16000} - \num{15000} = \num{9000}~\text{J}$.

2. Théorème de l'énergie cinétique : $\tfrac{1}{2}mv_B^2 = \tfrac{1}{2}mv_A^2 + \sum W = \tfrac{1}{2}\times 200 \times 15^2 + 9000 = \num{22500} + 9000 = \num{31500}~\text{J}$. D'où $v_B = \sqrt{2\times 31500/200} = \sqrt{315} \approx 17{,}7~\text{m/s} \approx 64~\text{km/h}$. La moto accélère dans la côte : la force motrice surcompense le poids et les frottements, ce que confirme le signe positif de $\sum W$.

3. Durée de la montée (mouvement rectiligne uniformément varié car forces constantes) : $\Delta t = \dfrac{2L}{v_A + v_B} = \dfrac{200}{15 + 17{,}7} \approx 6{,}1~\text{s}$. Puissance moyenne du moteur : $P_m = \dfrac{W(\vec{F})}{\Delta t} = \dfrac{\num{40000}}{6{,}1} \approx 6{,}5~\text{kW}$.
\end{corrige}

\begin{aretenir}[L'essentiel de la leçon]
\begin{itemize}
\item L'énergie cinétique mesure la « quantité de mouvement énergétique » d'un solide : $\tfrac{1}{2}mv^2$ en translation, $\tfrac{1}{2}J_{\Delta}\omega^2$ en rotation, somme des deux en mouvement combiné.
\item Le théorème de l'énergie cinétique relie la variation de $E_c$ à la somme des travaux des forces extérieures : $\Delta E_c = \sum W_{AB}$, dans un référentiel galiléen.
\item La méthode en 5 étapes (système, bilan des forces, travaux, TEC, calcul vérifié) résout tous les problèmes simples du programme.
\item Complément : la puissance moyenne $P_m = W/\Delta t$ relie l'énergie au temps ; elle s'exprime en watts.
\end{itemize}
\end{aretenir}

Tu disposes maintenant de tous les outils du chapitre. Entraîne-toi sur des situations variées — freinage, côtes, treuils, lancements — en rédigeant toujours la fiche méthode complète : au Probatoire, c'est la rigueur de la démarche autant que le résultat numérique qui est récompensée.

\newpage

\coursec{Activité d'intégration}

\courssub{Objectifs de l'activité}

À l'issue de cette activité d'intégration, tu seras capable de :

\begin{itemize}
\item calculer l'\cle{énergie cinétique de translation} d'un véhicule lourd à partir de sa masse et de sa vitesse ;
\item calculer l'\cle{énergie cinétique de rotation} de roues en \cle{roulement sans glissement} et l'énergie cinétique totale d'un solide en mouvement combiné ;
\item appliquer le \cle{théorème de l'énergie cinétique} pour déterminer une distance de freinage et analyser l'influence de la vitesse sur la sécurité routière ;
\item utiliser le théorème de l'énergie cinétique dans le cas d'une montée à vitesse constante (travail du poids) ;
\item convertir des joules en kilowattheures et estimer un \cle{coût énergétique} réel en carburant, en FCFA ;
\item argumenter un conseil de sécurité routière à l'aide de résultats chiffrés.
\end{itemize}

\courssub{Situation}

La liaison Yaoundé--Bafoussam, longue d'environ 300 km, est l'un des axes routiers les plus fréquentés du Cameroun. Après l'ascension redoutée de la côte de Santchou et les longues descentes de l'Ouest, les accidents impliquant des cars de transport en commun y sont malheureusement fréquents. La société de transport «\;Voyages de l'Ouest\;» souhaite sensibiliser ses chauffeurs à l'énergie mise en jeu par leurs véhicules et fait appel à ta classe de Première C pour produire une affiche pédagogique.

On étudie un car de masse totale $M = \SI{12}{t} = \num{12000}{kg}$ (carrosserie, passagers et bagages compris), roulant sur route horizontale à la vitesse $v = \SI{90}{km/h}$. Le car possède $4$ roues pleines, assimilées à des cylindres homogènes de masse $m = \SI{80}{kg}$ chacune et de rayon $R = \SI{0,5}{m}$. Les roues \emph{roulent sans glisser} sur la chaussée.

\textbf{Données :}
\begin{itemize}
\item intensité de la pesanteur : $g = \SI{9,8}{N/kg}$ ;
\item moment d'inertie d'un cylindre plein autour de son axe : $J_{\Delta} = \dfrac{1}{2}mR^{2}$ ;
\item force de freinage maximale (freins + adhérence) : $F = \SI{60}{kN} = \num{60000}{N}$, supposée constante ;
\item côte gravie à vitesse constante : longueur $L = \SI{8}{km}$, pente de $5\,\%$ (dénivelé de $5$ m pour $100$ m parcourus) ;
\item pouvoir énergétique du gasoil : environ $\SI{10}{kWh}$ par litre, rendement du moteur diesel : $30\,\%$ ;
\item prix du litre de gasoil à la station : $800$ FCFA ;
\item rappel : $\SI{1}{kWh} = \num{3,6e6}{J} = \SI{3,6}{MJ}$.
\end{itemize}

\textbf{Problème posé :} quelle énergie le car emmagasine-t-il à $90$ km/h ? Quelle distance faut-il pour l'arrêter ? Quel est le coût énergétique, en carburant et en francs CFA, de la montée d'une grande côte ? Quel conseil de sécurité peut-on en tirer ?

\courssub{Consignes}

Travail en groupes de 4 élèves. Chaque groupe rédigera ses réponses soigneusement, avec les unités SI, puis préparera une affiche de restitution.

\begin{enumerate}
\item Convertir la vitesse $v = \SI{90}{km/h}$ en m/s. Calculer l'énergie cinétique de translation $E_{c,\text{trans}}$ du car.
\item Calculer la vitesse angulaire $\omega$ des roues en roulement sans glissement, puis le moment d'inertie $J_{\Delta}$ d'une roue. En déduire l'énergie cinétique de rotation des quatre roues, puis l'énergie cinétique totale $E_c$ du car. La rotation des roues est-elle négligeable ?
\item À l'aide du théorème de l'énergie cinétique, déterminer la distance de freinage $d$ du car lancé à $90$ km/h sur route horizontale, soumis à la force de freinage $F = \SI{60}{kN}$. Recommencer le calcul pour $v' = \SI{45}{km/h}$. Comparer et commenter : que se passe-t-il quand la vitesse est divisée par $2$ ?
\item Le car gravit maintenant la côte de $\SI{8}{km}$ à $5\,\%$ \emph{à vitesse constante}. Calculer le dénivelé $h$, puis, en appliquant le théorème de l'énergie cinétique, le travail supplémentaire $W_{\text{moteur}}$ que doit fournir le moteur contre le poids.
\item Estimer en joules puis en kWh l'énergie mécanique totale dépensée (mise en vitesse à $90$ km/h + montée de la côte). Sachant qu'un litre de gasoil ne fournit utilement que $30\,\%$ de ses $\SI{10}{kWh}$, calculer le volume de carburant nécessaire et son coût en FCFA.
\item Rédiger sur l'affiche un conseil de sécurité routière (une phrase) destiné aux chauffeurs, justifié par au moins deux de tes résultats chiffrés.
\end{enumerate}

\courssub{Ressources à mobiliser}

\begin{aretenir}[Boîte à outils de la leçon]
\begin{itemize}
\item Énergie cinétique de translation : $E_{c,\text{trans}} = \dfrac{1}{2}Mv^{2}$ (en joules si $M$ en kg et $v$ en m/s).
\item Énergie cinétique de rotation : $E_{c,\text{rot}} = \dfrac{1}{2}J_{\Delta}\omega^{2}$.
\item Roulement sans glissement : $v = R\omega$.
\item Solide en mouvement combiné : $E_c = \dfrac{1}{2}Mv^{2} + \dfrac{1}{2}J_{\Delta}\omega^{2}$ (somme sur toutes les parties en rotation).
\item Théorème de l'énergie cinétique : entre deux instants, $\Delta E_c = E_{c,\text{finale}} - E_{c,\text{initiale}} = \sum W(\vec{F}_{\text{ext}})$.
\item Travail d'une force constante sur un trajet rectiligne : $W = F\,d\cos\alpha$ ; travail du poids en montée : $W(\vec{P}) = -Mgh$.
\item Conversion : $\SI{1}{km/h} = \dfrac{1}{3,6}\ \text{m/s}$ ; $\SI{1}{kWh} = \SI{3,6}{MJ}$.
\end{itemize}
\end{aretenir}

\courssub{Démarche et corrigé commenté}

\begin{exoresolu}[Le coût énergétique du transport Yaoundé--Bafoussam]
\textbf{Énoncé.} Un car de masse $M = \num{12000}$ kg, muni de 4 roues pleines ($m = 80$ kg, $R = 0,5$ m), roule à $90$ km/h. 1) Énergie cinétique de translation. 2) Énergie cinétique de rotation des roues et énergie totale. 3) Distance de freinage sous $F = 60$ kN à $90$ km/h puis à $45$ km/h. 4) Travail du moteur pour gravir à vitesse constante une côte de $8$ km à $5\,\%$. 5) Énergie totale en kWh, consommation de gasoil (10 kWh/L, rendement $30\,\%$) et coût à $800$ FCFA/L.

\tcblower

\textbf{1) Énergie cinétique de translation.}

Conversion de la vitesse : $v = \dfrac{90}{3,6} = \SI{25}{m/s}$.
\[
E_{c,\text{trans}} = \frac{1}{2}Mv^{2} = \frac{1}{2}\times \num{12000}\times 25^{2} = \num{6000}\times 625 = \num{3,75e6}{J}.
\]
\[
\boxed{E_{c,\text{trans}} = \SI{3,75}{MJ}}
\]
C'est l'énergie qu'il a fallu fournir pour lancer le car à $90$ km/h, et celle qu'il faudra dissiper pour l'arrêter.

\textbf{2) Énergie cinétique de rotation des roues.}

Roulement sans glissement : $\omega = \dfrac{v}{R} = \dfrac{25}{0,5} = \SI{50}{rad/s}$.

Moment d'inertie d'une roue (cylindre plein) :
\[
J_{\Delta} = \frac{1}{2}mR^{2} = \frac{1}{2}\times 80\times (0,5)^{2} = \SI{10}{kg.m^{2}}.
\]

Énergie de rotation des quatre roues :
\[
E_{c,\text{rot}} = 4\times\frac{1}{2}J_{\Delta}\omega^{2} = 4\times\frac{1}{2}\times 10\times 50^{2} = \num{5,0e4}{J} = \SI{50}{kJ}.
\]

Énergie cinétique totale :
\[
\boxed{E_c = E_{c,\text{trans}} + E_{c,\text{rot}} = \num{3,75e6} + \num{0,05e6} = \SI{3,80}{MJ}}
\]

\emph{Commentaire :} la rotation des roues ne représente qu'environ $1,3\,\%$ du total. Elle est faible ici, mais non nulle : un calcul rigoureux l'inclut.

\textbf{3) Distance de freinage.}

Sur route horizontale, le poids et la réaction normale de la route, perpendiculaires au déplacement, ne travaillent pas. Seule la force de freinage $\vec{F}$, opposée au mouvement, travaille : $W = -F\,d$.

Théorème de l'énergie cinétique entre le début du freinage ($E_c = \SI{3,80}{MJ}$) et l'arrêt ($E_c = 0$) :
\[
0 - E_c = -F\,d \quad\Longrightarrow\quad d = \frac{E_c}{F} = \frac{\num{3,80e6}}{\num{60000}} \approx \SI{63,3}{m}.
\]
\[
\boxed{d \approx \SI{63}{m} \text{ à } 90\ \text{km/h}}
\]

À $v' = \SI{45}{km/h} = \SI{12,5}{m/s}$ (vitesse divisée par 2) :
\[
E'_{c,\text{trans}} = \frac{1}{2}\times\num{12000}\times (12,5)^{2} = \num{937500}{J}, \qquad
E'_{c,\text{rot}} = 4\times\frac{1}{2}\times 10\times 25^{2} = \num{12500}{J},
\]
soit $E'_c = \SI{950}{kJ}$, d'où :
\[
\boxed{d' = \frac{\num{950000}}{\num{60000}} \approx \SI{15,8}{m}}
\]

\emph{Commentaire :} en divisant la vitesse par $2$, l'énergie cinétique et la distance de freinage sont divisées par $4$ (car $E_c \propto v^{2}$). C'est l'argument physique majeur des limitations de vitesse.

\textbf{4) Montée de la côte à vitesse constante.}

Dénivelé : $h = L\times 5\,\% = \num{8000}\times 0,05 = \SI{400}{m}$.

À vitesse constante, $\Delta E_c = 0$, donc $\sum W = 0$ : le moteur doit compenser exactement le travail résistant du poids (on néglige les frottements) :
\[
W_{\text{moteur}} = -W(\vec{P}) = Mgh = \num{12000}\times 9,8\times 400.
\]
\[
\boxed{W_{\text{moteur}} \approx \SI{47,0}{MJ}}
\]

\emph{Commentaire :} gravir la côte coûte plus de \textbf{douze fois} l'énergie nécessaire pour lancer le car à $90$ km/h. C'est pourquoi la consommation explose dans les reliefs de l'Ouest camerounais.

\textbf{5) Coût énergétique et carburant.}

Énergie mécanique totale : $E_{\text{tot}} = 3,8 + 47,0 = \SI{50,8}{MJ}$, soit :
\[
E_{\text{tot}} = \frac{50,8}{3,6} \approx \SI{14,1}{kWh}.
\]

Un litre de gasoil ne fournit utilement que $10\times 0,30 = \SI{3}{kWh}$. Volume nécessaire :
\[
V = \frac{14,1}{3} \approx \SI{4,7}{L} \quad\Longrightarrow\quad \text{Coût} \approx 4,7\times 800 \approx \num{3760}{FCFA}.
\]
\[
\boxed{V \approx \SI{4,7}{L}, \quad \text{coût} \approx \num{3800}{FCFA}}
\]

\textbf{6) Conseil de sécurité (exemple pour l'affiche).}

«\;Chauffeur, à $90$ km/h ton car de 12 tonnes emmagasine près de $4$ millions de joules : il te faudra plus de $60$ m pour t'arrêter. En roulant à $45$ km/h, cette distance tombe à moins de $16$ m. Diviser la vitesse par deux, c'est diviser le danger par quatre. Ralentis avant les descentes de l'Ouest !\;»
\end{exoresolu}

\begin{attention}[Pièges fréquents rencontrés pendant l'activité]
\begin{itemize}
\item \textbf{Oublier la conversion km/h $\to$ m/s} : avec $v = 90$ au lieu de $25$, l'énergie est fausse d'un facteur $13$ ! Toujours convertir avant de calculer.
\item \textbf{Appliquer $E_c = \frac{1}{2}Mv^{2}$ aux roues} : une roue qui roule possède les deux formes d'énergie ; sa seule rotation fait intervenir $J_{\Delta}$ et $\omega$, pas $v$.
\item \textbf{Confondre pente et angle} : une pente de $5\,\%$ signifie $h = 0,05\,L$ ; inutile de passer par le sinus de l'angle pour une pente faible.
\item \textbf{Dans la montée à vitesse constante}, écrire $\Delta E_c = W_{\text{moteur}}$ en oubliant le poids : c'est la \emph{somme} des travaux qui est nulle.
\item \textbf{Oublier le rendement} : les $10$ kWh d'un litre de gasoil ne sont pas intégralement convertis en énergie mécanique.
\end{itemize}
\end{attention}

\courssub{Bilan de l'activité}

Cette activité a permis de mobiliser \emph{tous} les savoir-faire de la leçon dans une situation authentique de mobilité :

\begin{itemize}
\item l'énergie cinétique de translation donne l'ordre de grandeur de l'énergie «\;embarquée\;» par un véhicule : quelques mégajoules pour un car de 12 tonnes, soit l'équivalent de l'énergie électrique consommée par une famille en plusieurs jours ;
\item le calcul de l'énergie de rotation des roues montre qu'un solide réel est en \emph{mouvement combiné} : $E_c = \frac{1}{2}Mv^{2} + \sum\frac{1}{2}J_{\Delta}\omega^{2}$, même si la part rotationnelle est ici modeste ;
\item le théorème de l'énergie cinétique s'est révélé être un outil puissant et direct : il relie une \emph{variation d'énergie} à des \emph{travaux de forces}, sans avoir besoin de résoudre l'équation du mouvement ;
\item la dépendance en $v^{2}$ de l'énergie cinétique explique quantitativement pourquoi la distance de freinage est multipliée par $4$ quand la vitesse double : un résultat au cœur de la prévention routière ;
\item enfin, la conversion en kWh et en FCFA a donné un sens concret au joule : l'énergie n'est pas une abstraction, elle s'achète à la station-service, et le relief camerounais (côtes de l'Ouest, montées de l'Adamaoua) pèse lourd dans le budget carburant des transporteurs.
\end{itemize}

Tu as ainsi franchi le pas essentiel du physicien : passer d'une situation réelle à un modèle (solide en translation + roues en rotation), appliquer des lois, calculer, puis \emph{revenir à la réalité} pour formuler un conseil argumenté. C'est exactement la démarche attendue au Probatoire comme dans la vie professionnelle.

\newpage

\coursec{Méthodes et exercices résolus}

\coursec{Énergie cinétique et théorème de l'énergie cinétique}

\courssub{Notion d'énergie et énergie cinétique de translation}

\begin{definition}[Énergie cinétique de translation]
L'énergie cinétique de translation d'un solide de masse $m$ en translation de vitesse $\vec{v}_G$ (vitesse de son centre d'inertie $G$) dans un référentiel galiléen est la grandeur scalaire définie par :
\[ E_c = \frac{1}{2}\,m\,v_G^2. \]
Son unité SI est le \cle{joule} (symbole \cle{J}) : $1\,\mathrm{J} = 1\,\mathrm{kg\cdot m^2\cdot s^{-2}}$.
L'énergie cinétique est toujours \cle{positive ou nulle} ; elle ne dépend que du module de la vitesse, pas de sa direction.
\end{definition}

\begin{methode}[Calculer l'énergie cinétique de translation]
\begin{enumerate}
\item \textbf{Vérifier la nature du mouvement} : le solide est en translation si tous ses points ont, à chaque instant, la même vitesse que le centre d'inertie $G$.
\item \textbf{Relever les données} : masse $m$ du solide (en kilogrammes) et vitesse $v_G$ du centre d'inertie (en $\mathrm{m\cdot s^{-1}}$ ; convertir si donnée en $\mathrm{km\cdot h^{-1}}$).
\item \textbf{Convertir la vitesse} si nécessaire : $v\ (\mathrm{m\cdot s^{-1}}) = \dfrac{v\ (\mathrm{km\cdot h^{-1}})}{3,6}$.
\item \textbf{Appliquer la formule} : $E_c = \dfrac{1}{2}mv_G^2$.
\item \textbf{Exprimer le résultat en joules (J)}, avec un nombre raisonnable de chiffres significatifs, et vérifier l'homogénéité : $\mathrm{kg\cdot m^2\cdot s^{-2}} = \mathrm{J}$.
\end{enumerate}
\begin{aretenir}[Réflexes]
Si la vitesse double, l'énergie cinétique est multipliée par \cle{quatre} (variation quadratique).
\end{aretenir}
\end{methode}

\begin{exoresolu}[Le mototaxi de Bafoussam]
Un mototaxi (bend-skin) de masse totale $m = \SI{180}{kg}$ (moto + conducteur + passager) circule à $v = \SI{54}{km\cdot h^{-1}}$ sur l'avenue principale de Bafoussam.
\begin{enumerate}
\item Calculer son énergie cinétique.
\item Que devient cette énergie cinétique si la vitesse est réduite à $\SI{27}{km\cdot h^{-1}}$ ? Conclure.
\end{enumerate}
\tcblower
\textbf{1. Énergie cinétique à $\SI{54}{km\cdot h^{-1}}$.}

La moto est en translation : tous ses points ont la même vitesse que son centre d'inertie.

Conversion de la vitesse en unité SI :
\[ v = \dfrac{54}{3,6} = \SI{15}{m\cdot s^{-1}}. \]

Application de la formule :
\[ E_c = \frac{1}{2}mv^2 = \frac{1}{2}\times 180 \times 15^2 = 90 \times 225 = \SI{20250}{J} \approx \SI{2,0e4}{J}. \]

Vérification de l'homogénéité : $\mathrm{kg}\times(\mathrm{m\cdot s^{-1}})^2 = \mathrm{kg\cdot m^2\cdot s^{-2}} = \mathrm{J}$. Le résultat est cohérent.

\textbf{2. Vitesse réduite à $\SI{27}{km\cdot h^{-1}}$.}

Nouvelle vitesse : $v' = \dfrac{27}{3,6} = \SI{7,5}{m\cdot s^{-1}}$, c'est-à-dire $v' = \dfrac{v}{2}$.

\[ E_c' = \frac{1}{2}mv'^2 = \frac{1}{2}\times 180 \times (7,5)^2 = 90 \times 56,25 = \SI{5062,5}{J} \approx \SI{5,1e3}{J}. \]

On constate que $E_c' = \dfrac{E_c}{4}$.

\textbf{Conclusion :} l'énergie cinétique du mototaxi vaut environ $\SI{2,0e4}{J}$ à $\SI{54}{km\cdot h^{-1}}$ ; lorsque la vitesse est divisée par deux, l'énergie cinétique est divisée par \textbf{quatre}, car elle varie comme le \emph{carré} de la vitesse. C'est pourquoi les distances de freinage augmentent très vite avec la vitesse.
\end{exoresolu}

\courssub{Énergie cinétique d'un solide en rotation}

\begin{definition}[Moment d'inertie et énergie cinétique de rotation]
Le \cle{moment d'inertie} $J_{\Delta}$ d'un solide par rapport à un axe fixe $(\Delta)$ caractérise sa répartition de masse par rapport à cet axe. Son unité SI est $\mathrm{kg\cdot m^2}$.
L'énergie cinétique de rotation d'un solide tournant autour d'un axe fixe $(\Delta)$ à la vitesse angulaire $\omega$ est :
\[ E_c = \frac{1}{2}\,J_{\Delta}\,\omega^2. \]
\end{definition}

\begin{propriete}[Moments d'inertie usuels (axe de symétrie)]
\begin{center}
\begin{tabularx}{\linewidth}{|l|X|X|}\hline
\textbf{Solide} & \textbf{Axe} & \textbf{Moment d'inertie} \\ \hline
Cylindre plein (disque) & Axe de révolution $(\Delta)$ & $J_{\Delta} = \dfrac{1}{2}MR^2$ \\ \hline
Cerceau (anneau mince) & Axe de révolution $(\Delta)$ & $J_{\Delta} = MR^2$ \\ \hline
Tige homogène & Axe $\perp$ à la tige passant par $G$ & $J_{\Delta} = \dfrac{1}{12}ML^2$ \\ \hline
Sphère pleine & Tout diamètre & $J_{\Delta} = \dfrac{2}{5}MR^2$ \\ \hline
\end{tabularx}
\end{center}
\end{propriete}

\begin{methode}[Calculer l'énergie cinétique de rotation]
\begin{enumerate}
\item \textbf{Identifier l'axe de rotation} et vérifier que le solide tourne autour d'un axe fixe $(\Delta)$.
\item \textbf{Déterminer le moment d'inertie} $J_{\Delta}$ (en $\mathrm{kg\cdot m^2}$) : donné par l'énoncé ou calculé via la formule adaptée à la géométrie.
\item \textbf{Convertir la vitesse angulaire} en $\mathrm{rad\cdot s^{-1}}$ si elle est donnée en tours par minute : $\omega\ (\mathrm{rad\cdot s^{-1}}) = N\ (\mathrm{tr/min}) \times \dfrac{2\pi}{60}$.
\item \textbf{Appliquer la formule} : $E_c = \dfrac{1}{2}J_{\Delta}\,\omega^2$.
\item \textbf{Exprimer le résultat en joules} et vérifier l'homogénéité : $\mathrm{kg\cdot m^2}\times\mathrm{rad^2\cdot s^{-2}} = \mathrm{J}$ (le radian est sans dimension).
\end{enumerate}
\begin{attention}[Réflexe de conversion]
L'erreur la plus fréquente est d'oublier de convertir les tours par minute en $\mathrm{rad\cdot s^{-1}}$. Un tour correspond à $2\pi$ radians, et une minute à $\SI{60}{s}$ : ne jamais écrire $\omega = N$ directement.
\end{attention}
\end{methode}

\begin{exoresolu}[Le volant d'une machine à décortiquer]
Un volant de machine agricole, assimilé à un cylindre plein homogène de masse $M = \SI{40}{kg}$ et de rayon $R = \SI{0,30}{m}$, tourne autour de son axe de révolution $(\Delta)$ à la fréquence de $N = \SI{300}{tr/min}$.
\begin{enumerate}
\item Calculer son moment d'inertie $J_{\Delta}$.
\item Calculer son énergie cinétique de rotation.
\end{enumerate}
On donne : pour un cylindre plein, $J_{\Delta} = \dfrac{1}{2}MR^2$.
\tcblower
\textbf{1. Moment d'inertie.}

\[ J_{\Delta} = \frac{1}{2}MR^2 = \frac{1}{2}\times 40 \times (0,30)^2 = 20 \times 0,09 = \SI{1,8}{kg\cdot m^2}. \]

\textbf{2. Énergie cinétique de rotation.}

Conversion de la vitesse angulaire :
\[ \omega = N\times\frac{2\pi}{60} = 300\times\frac{2\pi}{60} = 10\pi \approx \SI{31,4}{rad\cdot s^{-1}}. \]

Application de la formule :
\[ E_c = \frac{1}{2}J_{\Delta}\,\omega^2 = \frac{1}{2}\times 1,8 \times (10\pi)^2 = 0,9 \times 100\pi^2 = 90\pi^2. \]

Numériquement, avec $\pi^2 \approx 9,87$ :
\[ E_c \approx 90 \times 9,87 \approx \SI{8,9e2}{J}. \]

Vérification de l'homogénéité : $\mathrm{kg\cdot m^2 \cdot s^{-2}} = \mathrm{J}$.

\textbf{Conclusion :} le moment d'inertie du volant vaut $J_{\Delta} = \SI{1,8}{kg\cdot m^2}$ et son énergie cinétique de rotation vaut environ $\SI{8,9e2}{J}$. Cette énergie stockée permet au volant de régulariser la rotation de la machine entre deux impulsions du moteur.
\end{exoresolu}

\courssub{Solide en mouvement combiné : translation et rotation}

\begin{propriete}[Théorème de König pour l'énergie cinétique]
L'énergie cinétique d'un solide en mouvement quelconque dans un référentiel galiléen est la somme de l'énergie cinétique de translation de son centre d'inertie $G$ (affecté de toute la masse $m$) et de l'énergie cinétique de rotation autour d'un axe passant par $G$ :
\[ E_c = \underbrace{\frac{1}{2}mv_G^2}_{\text{translation de }G} + \underbrace{\frac{1}{2}J_G\,\omega^2}_{\text{rotation autour de }G}. \]
Si le solide roule \emph{sans glisser} sur un support immobile, la condition de roulement lie les deux mouvements : $v_G = R\,\omega$ (où $R$ est le rayon de roulement).
\end{propriete}

\begin{methode}[Décomposer puis additionner les deux énergies cinétiques]
\begin{enumerate}
\item \textbf{Décomposer le mouvement} : translation de $G$ (vitesse $\vec{v}_G$) et rotation autour de $G$ (vitesse angulaire $\omega$).
\item \textbf{Calculer séparément} les deux termes : $E_c = \frac{1}{2}mv_G^2 + \frac{1}{2}J_G\,\omega^2$.
\item \textbf{Exploiter la condition de roulement sans glissement} si le solide roule : $v_G = R\,\omega$, ce qui permet de relier les deux termes.
\item \textbf{Additionner} les deux contributions et exprimer le résultat en joules.
\end{enumerate}
\end{methode}

\begin{exoresolu}[La roue de la moto]
Une roue de moto, assimilée à un cerceau de masse $m = \SI{12}{kg}$ et de rayon $R = \SI{0,30}{m}$, roule \emph{sans glisser} sur une route à la vitesse $v = \SI{72}{km\cdot h^{-1}}$. On donne : pour un cerceau, $J_G = mR^2$.
\begin{enumerate}
\item Calculer la vitesse angulaire de rotation de la roue.
\item Calculer son énergie cinétique totale.
\end{enumerate}
\tcblower
\textbf{1. Vitesse angulaire.}

Conversion de la vitesse : $v = \dfrac{72}{3,6} = \SI{20}{m\cdot s^{-1}}$.

Le roulement a lieu sans glissement, donc la condition de roulement s'applique :
\[ v_G = R\,\omega \quad\Longrightarrow\quad \omega = \frac{v_G}{R} = \frac{20}{0,30} \approx \SI{66,7}{rad\cdot s^{-1}}. \]

\textbf{2. Énergie cinétique totale.}

Le mouvement de la roue se décompose en translation de $G$ et rotation autour de $G$ :

Énergie cinétique de translation :
\[ E_{c,\mathrm{trans}} = \frac{1}{2}mv_G^2 = \frac{1}{2}\times 12 \times 20^2 = 6 \times 400 = \SI{2400}{J}. \]

Moment d'inertie du cerceau :
\[ J_G = mR^2 = 12 \times (0,30)^2 = 12 \times 0,09 = \SI{1,08}{kg\cdot m^2}. \]

Énergie cinétique de rotation :
\[ E_{c,\mathrm{rot}} = \frac{1}{2}J_G\,\omega^2 = \frac{1}{2}\times 1,08 \times \left(\frac{20}{0,30}\right)^2 = 0,54 \times \frac{400}{0,09} = \SI{2400}{J}. \]

Énergie cinétique totale :
\[ E_c = E_{c,\mathrm{trans}} + E_{c,\mathrm{rot}} = 2400 + 2400 = \SI{4800}{J} = \SI{4,8e3}{J}. \]

\textbf{Conclusion :} la roue tourne à $\omega \approx \SI{66,7}{rad\cdot s^{-1}}$ et possède une énergie cinétique totale de $\SI{4,8e3}{J}$. Remarque : pour un cerceau qui roule sans glisser, l'énergie de rotation est \emph{égale} à l'énergie de translation ; négliger la rotation sous-estimerait l'énergie d'un facteur deux.
\end{exoresolu}

\courssub{Théorème de l'énergie cinétique}

\begin{propriete}[Théorème de l'énergie cinétique]
Dans un \cle{référentiel galiléen}, la variation de l'énergie cinétique d'un système (solide indéformable) entre deux instants $t_A$ et $t_B$ est égale à la somme des travaux des forces extérieures appliquées au système au cours du déplacement $AB$ :
\[ \Delta E_c = E_c(B) - E_c(A) = \sum W_{AB}(\vec{F}_{\mathrm{ext}}). \]
\emph{Condition d'application :} le référentiel doit être galiléen (référentiel terrestre supposé galiléen au lycée).
\end{propriete}

\begin{methode}[Appliquer le théorème de l'énergie cinétique dans un cas simple]
\begin{enumerate}
\item \textbf{Définir le système} (le solide étudié) et le \textbf{référentiel} (en général le référentiel terrestre, supposé galiléen).
\item \textbf{Préciser l'état initial et l'état final} : positions $A$ et $B$, vitesses $v_A$ et $v_B$.
\item \textbf{Faire le bilan des forces extérieures} appliquées au solide pendant le déplacement (schéma à l'appui).
\item \textbf{Exprimer le travail de chaque force} entre $A$ et $B$ :
      \begin{itemize}
      \item Force constante : $W_{AB}(\vec{F}) = \vec{F}\cdot\overrightarrow{AB} = F\times AB\times\cos\alpha$.
      \item Poids : $W_{AB}(\vec{P}) = mg(z_A - z_B)$.
      \item Force perpendiculaire au déplacement : travail nul.
      \end{itemize}
\item \textbf{Écrire le théorème} :
      \[ \frac{1}{2}mv_B^2 - \frac{1}{2}mv_A^2 = \sum W_{AB}(\vec{F}_{\mathrm{ext}}). \]
\item \textbf{Isoler la grandeur inconnue} (vitesse, force, distance) et calculer numériquement en unités SI.
\item \textbf{Vérifier la cohérence} du signe : un travail moteur ($W>0$) augmente l'énergie cinétique ; un travail résistant ($W<0$) la diminue.
\end{enumerate}
\end{methode}

\begin{exoresolu}[Freinage d'un minicar sur l'autoroute Yaoundé--Douala]
Un minicar de masse $m = \SI{2500}{kg}$ roule à $v_A = \SI{90}{km\cdot h^{-1}}$. Le conducteur freine et le véhicule s'arrête après une distance $d = \SI{60}{m}$. On suppose la force de freinage $\vec{f}$ constante, parallèle au déplacement et de sens opposé ; la route est horizontale.
\begin{enumerate}
\item Énoncer le théorème de l'énergie cinétique.
\item Déterminer l'intensité $f$ de la force de freinage.
\end{enumerate}
\tcblower
\textbf{1. Énoncé du théorème.}

Dans un référentiel galiléen, la variation de l'énergie cinétique d'un solide entre deux instants est égale à la somme des travaux des forces extérieures appliquées au solide entre ces deux instants :
\[ \Delta E_c = E_{c,\mathrm{final}} - E_{c,\mathrm{initial}} = \sum W(\vec{F}_{\mathrm{ext}}). \]

\textbf{2. Intensité de la force de freinage.}

\textit{Système :} le minicar. \textit{Référentiel :} terrestre, supposé galiléen.

\textit{États initial et final :} en $A$, $v_A = \dfrac{90}{3,6} = \SI{25}{m\cdot s^{-1}}$ ; en $B$, le véhicule est à l'arrêt : $v_B = 0$.

\textit{Bilan des forces extérieures :}
\begin{itemize}
\item le poids $\vec{P}$, vertical, perpendiculaire au déplacement horizontal : $W_{AB}(\vec{P}) = 0$ ;
\item la réaction normale de la route $\vec{R}_N$, perpendiculaire au déplacement : $W_{AB}(\vec{R}_N) = 0$ ;
\item la force de freinage $\vec{f}$, colinéaire au déplacement et de sens opposé : $W_{AB}(\vec{f}) = -f\times d$ (travail résistant).
\end{itemize}

\textit{Application du théorème :}
\[ \frac{1}{2}mv_B^2 - \frac{1}{2}mv_A^2 = -f\,d. \]

Comme $v_B = 0$ :
\[ -\frac{1}{2}mv_A^2 = -f\,d \quad\Longrightarrow\quad f = \frac{mv_A^2}{2d}. \]

Application numérique :
\[ f = \frac{2500 \times 25^2}{2\times 60} = \frac{2500 \times 625}{120} = \frac{\num{1562500}}{120} \approx \SI{1,3e4}{N}. \]

\textbf{Conclusion :} la force de freinage moyenne nécessaire vaut environ $\SI{1,3e4}{N}$. On vérifie la cohérence : le travail du freinage est résistant ($W<0$), ce qui correspond bien à une diminution de l'énergie cinétique jusqu'à l'arrêt.
\end{exoresolu}

\courssub{Applications directes du théorème}

\begin{methode}[Théorème de l'énergie cinétique et travail du poids]
\begin{enumerate}
\item \textbf{Repérer les altitudes} : choisir un axe $Oz$ vertical ascendant et noter $z_A$ et $z_B$.
\item \textbf{Écrire le travail du poids} : $W_{AB}(\vec{P}) = mg(z_A - z_B)$. Il est \cle{moteur} si le solide descend ($z_A > z_B$) et \cle{résistant} s'il monte.
\item \textbf{Identifier les forces qui ne travaillent pas} : la réaction normale d'un support, perpendiculaire au déplacement, a un travail nul (s'il n'y a pas de frottements).
\item \textbf{Appliquer le théorème} en remplaçant chaque travail par son expression, puis isoler l'inconnue.
\item \textbf{Remarque utile} : sur un plan incliné d'angle $\alpha$, pour un déplacement $AB = L$, la dénivellation vaut $z_A - z_B = L\sin\alpha$, donc $W_{AB}(\vec{P}) = mgL\sin\alpha$ (descente).
\end{enumerate}
\begin{attention}[Le travail du poids ne dépend pas du chemin]
Le travail du poids entre $A$ et $B$ ne dépend que de la \cle{dénivellation} $z_A - z_B$, jamais de la forme de la trajectoire ni de la longueur du chemin parcouru. Ne pas multiplier le poids par la longueur du plan incliné sans le facteur $\sin\alpha$.
\end{attention}
\end{methode}

\begin{exoresolu}[Descente d'un chariot sur une piste]
Un chariot de masse $m = \SI{80}{kg}$ est lâché sans vitesse initiale au sommet $A$ d'une piste inclinée de $\alpha = 20^{\circ}$ par rapport à l'horizontale. Il descend jusqu'en $B$, avec $AB = L = \SI{25}{m}$. Les frottements sont supposés négligeables. On donne $g = \SI{9,8}{N\cdot kg^{-1}}$ et $\sin 20^{\circ} \approx 0,342$.
\begin{enumerate}
\item Faire le bilan des forces et calculer leurs travaux entre $A$ et $B$.
\item En déduire la vitesse $v_B$ du chariot en $B$.
\end{enumerate}
\tcblower
\textbf{1. Bilan des forces et travaux.}

\textit{Système :} le chariot. \textit{Référentiel :} terrestre, supposé galiléen.

Forces extérieures appliquées au chariot :
\begin{itemize}
\item le poids $\vec{P}$ ($P = mg$), vertical vers le bas ;
\item la réaction normale de la piste $\vec{R}_N$, perpendiculaire à la piste (frottements négligés).
\end{itemize}

Travail de la réaction normale : $\vec{R}_N$ est perpendiculaire au déplacement, donc
\[ W_{AB}(\vec{R}_N) = 0. \]

Travail du poids : la dénivellation entre $A$ et $B$ vaut $z_A - z_B = L\sin\alpha$. Le chariot descend, le travail du poids est moteur :
\[ W_{AB}(\vec{P}) = mgL\sin\alpha = 80 \times 9,8 \times 25 \times 0,342. \]

Calcul : $80 \times 9,8 = 784$ ; $784 \times 25 = \num{19600}$ ; $\num{19600} \times 0,342 \approx \SI{6703}{J}$, soit
\[ W_{AB}(\vec{P}) \approx \SI{6,7e3}{J}. \]

\textbf{2. Vitesse en B.}

Application du théorème de l'énergie cinétique entre $A$ et $B$ :
\[ \frac{1}{2}mv_B^2 - \frac{1}{2}mv_A^2 = W_{AB}(\vec{P}) + W_{AB}(\vec{R}_N). \]

Le chariot est lâché sans vitesse initiale : $v_A = 0$. Donc :
\[ \frac{1}{2}mv_B^2 = mgL\sin\alpha \quad\Longrightarrow\quad v_B = \sqrt{2gL\sin\alpha}. \]

Application numérique :
\[ v_B = \sqrt{2\times 9,8 \times 25 \times 0,342} = \sqrt{167,6} \approx \SI{12,9}{m\cdot s^{-1}}. \]

\textbf{Conclusion :} le chariot atteint en $B$ une vitesse d'environ $\SI{13}{m\cdot s^{-1}}$ (soit environ $\SI{47}{km\cdot h^{-1}}$). On remarque que la masse s'élimine du calcul : en l'absence de frottements, la vitesse acquise ne dépend que de la dénivellation, pas de la masse du chariot.
\end{exoresolu}

\begin{methode}[Problème avec frottements : démarche complète]
\begin{enumerate}
\item \textbf{Lire l'énoncé en identifiant les deux phases} éventuelles (montée/descente, lancement/arrêt) et choisir les états initial et final adaptés à la question.
\item \textbf{Modéliser les frottements} par une force $\vec{f}$ d'intensité constante, colinéaire au déplacement et de sens opposé : $W_{AB}(\vec{f}) = -f\times AB$ (toujours résistant).
\item \textbf{Écrire le théorème de l'énergie cinétique} avec tous les travaux : poids, frottements, réaction normale (nulle), force motrice éventuelle.
\item \textbf{Résoudre algébriquement avant l'application numérique} : isoler littéralement l'inconnue, puis substituer les valeurs en unités SI.
\item \textbf{Discuter le résultat} : signe, ordre de grandeur, comparaison avec une situation sans frottements.
\end{enumerate}
\end{methode}

\begin{exoresolu}[La caisse sur le quai de déchargement]
Pour décharger des sacs de cacao au port de Douala, un ouvrier pousse une caisse de masse $m = \SI{50}{kg}$ sur un quai horizontal. Il lui communique une vitesse initiale $v_A = \SI{4,0}{m\cdot s^{-1}}$, puis la laisse glisser. La caisse s'arrête en $B$ après avoir parcouru $d = \SI{5,0}{m}$. On donne $g = \SI{10}{N\cdot kg^{-1}}$.
\begin{enumerate}
\item Déterminer l'intensité $f$ de la force de frottement, supposée constante, exercée par le quai sur la caisse.
\item Quelle distance la caisse parcourrait-elle si sa vitesse initiale était $v_A' = \SI{8,0}{m\cdot s^{-1}}$, avec la même force de frottement ?
\end{enumerate}
\tcblower
\textbf{1. Intensité de la force de frottement.}

\textit{Système :} la caisse. \textit{Référentiel :} terrestre, supposé galiléen.

\textit{Bilan des forces extérieures} entre $A$ et $B$ :
\begin{itemize}
\item le poids $\vec{P}$, perpendiculaire au déplacement horizontal : $W_{AB}(\vec{P}) = 0$ ;
\item la réaction normale $\vec{R}_N$, perpendiculaire au déplacement : $W_{AB}(\vec{R}_N) = 0$ ;
\item la force de frottement $\vec{f}$, opposée au déplacement : $W_{AB}(\vec{f}) = -f\,d$.
\end{itemize}

\textit{Théorème de l'énergie cinétique} entre $A$ (vitesse $v_A$) et $B$ (arrêt, $v_B = 0$) :
\[ 0 - \frac{1}{2}mv_A^2 = -f\,d \quad\Longrightarrow\quad f = \frac{mv_A^2}{2d}. \]

Application numérique :
\[ f = \frac{50 \times (4,0)^2}{2 \times 5,0} = \frac{50 \times 16}{10} = \frac{800}{10} = \SI{80}{N}. \]

\textbf{2. Nouvelle distance d'arrêt.}

Avec la même force de frottement, le même raisonnement donne :
\[ d' = \frac{mv_A'^2}{2f} = \frac{50 \times (8,0)^2}{2 \times 80} = \frac{50 \times 64}{160} = \frac{3200}{160} = \SI{20}{m}. \]

Autre raisonnement (plus rapide) : la distance d'arrêt est proportionnelle au carré de la vitesse initiale. La vitesse ayant doublé, la distance est multipliée par quatre : $d' = 4\times 5,0 = \SI{20}{m}$. Les deux méthodes concordent.

\textbf{Conclusion :} la force de frottement vaut $f = \SI{80}{N}$, et doubler la vitesse initiale quadruple la distance d'arrêt ($d' = \SI{20}{m}$). C'est exactement la raison pour laquelle les distances de freinage des véhicules croissent si vite avec la vitesse.
\end{exoresolu}

\begin{methode}[Théorème de l'énergie cinétique pour un solide en rotation]
\begin{enumerate}
\item \textbf{Exprimer l'énergie cinétique de rotation} aux états initial et final : $E_c = \dfrac{1}{2}J_{\Delta}\omega^2$.
\item \textbf{Faire le bilan des forces} et identifier celles dont le travail est nul : le poids et les forces d'axe dont le point d'application est sur l'axe de rotation ne travaillent pas (bras de levier nul, pas de déplacement du point d'application).
\item \textbf{Écrire le théorème} sous la forme :
      \[ \frac{1}{2}J_{\Delta}\omega_B^2 - \frac{1}{2}J_{\Delta}\omega_A^2 = \sum W_{AB}(\vec{F}_{\mathrm{ext}}). \]
\item \textbf{Convertir systématiquement} les vitesses angulaires en $\mathrm{rad\cdot s^{-1}}$ avant tout calcul.
\item \textbf{Conclure} avec une valeur numérique accompagnée de son unité.
\end{enumerate}
\end{methode}

\begin{exoresolu}[Arrêt d'une meule]
Une meule de potier, assimilée à un disque homogène de moment d'inertie $J_{\Delta} = \SI{2,0}{kg\cdot m^2}$, tourne à $N_A = \SI{120}{tr/min}$. Le potier cesse de l'entraîner ; les frottements de l'axe exercent un couple résistant dont le travail vaut $W = \SI{-4,0}{J}$ par tour.
\begin{enumerate}
\item Calculer l'énergie cinétique initiale de la meule.
\item Déterminer le nombre de tours effectués par la meule avant l'arrêt complet.
\end{enumerate}
\tcblower
\textbf{1. Énergie cinétique initiale.}

Conversion de la vitesse angulaire initiale :
\[ \omega_A = N_A\times\frac{2\pi}{60} = 120\times\frac{2\pi}{60} = 4\pi \approx \SI{12,6}{rad\cdot s^{-1}}. \]

Énergie cinétique initiale :
\[ E_c(A) = \frac{1}{2}J_{\Delta}\omega_A^2 = \frac{1}{2}\times 2,0 \times (4\pi)^2 = 16\pi^2 \approx \SI{1,6e2}{J}. \]

\textbf{2. Nombre de tours avant l'arrêt.}

\textit{Système :} la meule. \textit{Référentiel :} terrestre, supposé galiléen.

\textit{Bilan des travaux :} le poids de la meule et la réaction de l'axe sont appliqués sur l'axe de rotation (ou ont un moment nul par rapport à l'axe) : leur travail est nul. Seul le couple de frottement travaille. Si la meule effectue $n$ tours avant l'arrêt :
\[ W_{\mathrm{total}} = n\times(-4,0)\ \mathrm{J}. \]

\textit{Théorème de l'énergie cinétique} entre l'état initial ($\omega_A$) et l'arrêt ($\omega_B = 0$) :
\[ 0 - \frac{1}{2}J_{\Delta}\omega_A^2 = -4,0\,n. \]

D'où :
\[ n = \frac{\frac{1}{2}J_{\Delta}\omega_A^2}{4,0} = \frac{16\pi^2}{4,0} = 4\pi^2 \approx \SI{39,5}{tours}. \]

\textbf{Conclusion :} l'énergie cinétique initiale de la meule vaut environ $\SI{1,6e2}{J}$ et la meule effectue environ \textbf{40 tours} avant de s'immobiliser. Le travail résistant des frottements a intégralement absorbé l'énergie cinétique de rotation initiale.
\end{exoresolu}

\courssub{Méthode de résolution et synthèse}

\begin{attention}[Les 5 erreurs qui coûtent des points au Probatoire]
\begin{enumerate}
\item \textbf{Oublier les conversions d'unités.} Une vitesse en $\mathrm{km\cdot h^{-1}}$ doit être divisée par $3,6$ pour obtenir des $\mathrm{m\cdot s^{-1}}$ ; une vitesse angulaire en $\mathrm{tr/min}$ doit être multipliée par $\dfrac{2\pi}{60}$. Toute énergie calculée avec des unités non SI est fausse, même si la formule est correcte.
\item \textbf{Écrire $E_c = mv^2$ en oubliant le facteur $\dfrac{1}{2}$}, ou oublier le carré sur la vitesse. La formule est $E_c = \dfrac{1}{2}mv^2$ (translation) et $E_c = \dfrac{1}{2}J_{\Delta}\omega^2$ (rotation). Une vérification dimensionnelle rapide permet de détecter l'oubli.
\item \textbf{Mal rédiger le théorème de l'énergie cinétique.} Il faut écrire une \emph{variation} : $E_c(B) - E_c(A) = \sum W_{AB}(\vec{F}_{\mathrm{ext}})$, dans cet ordre (final moins initial). Inverser les deux membres change tous les signes et fausse le résultat. Préciser aussi le système et le référentiel galiléen : c'est une condition d'application exigée.
\item \textbf{Faire travailler les forces perpendiculaires au déplacement.} Le poids et la réaction normale ne travaillent pas sur un trajet horizontal ; la réaction normale d'un support sans frottement ne travaille jamais. À l'inverse, oublier le signe négatif du travail des frottements ($W = -f\,d$) est une erreur rédhibitoire : un travail résistant doit apparaître avec un signe moins.
\item \textbf{Se tromper dans le travail du poids.} $W_{AB}(\vec{P}) = mg(z_A - z_B)$ : moteur en descente, résistant en montée. Sur un plan incliné, il faut la dénivellation $L\sin\alpha$, et non la longueur $L$ du trajet. Ne jamais écrire $W(\vec{P}) = P\times AB$ sans le facteur $\sin\alpha$.
\end{enumerate}
\end{attention}

\newpage

\coursec{Exercices}

\courssub{Je vérifie mes connaissances}

\begin{exercice}[QCM : notions fondamentales]
Pour chaque question, une seule réponse est exacte. Recopie la lettre correspondante.
\begin{enumerate}
\item L'unité SI de l'énergie cinétique est :
\begin{enumerate}
\item le watt (W) ;
\item le joule (J) ;
\item le newton (N) ;
\item le newton-mètre par seconde (N.m/s).
\end{enumerate}
\item Si la vitesse d'un solide en translation est multipliée par 3, son énergie cinétique est multipliée par :
\begin{enumerate}
\item 3 ;
\item 6 ;
\item 9 ;
\item $\sqrt{3}$.
\end{enumerate}
\item L'énergie cinétique d'un solide de moment d'inertie $J_{\Delta}$ tournant à la vitesse angulaire $\omega$ autour d'un axe fixe $\Delta$ vaut :
\begin{enumerate}
\item $E_c = \dfrac{1}{2}J_{\Delta}\,\omega$ ;
\item $E_c = \dfrac{1}{2}J_{\Delta}\,\omega^2$ ;
\item $E_c = J_{\Delta}\,\omega^2$ ;
\item $E_c = \dfrac{1}{2}m\,\omega^2$.
\end{enumerate}
\item Le théorème de l'énergie cinétique s'écrit, entre deux positions A et B :
\begin{enumerate}
\item $E_{cB} + E_{cA} = \sum W_{A\to B}(\vec{F})$ ;
\item $\Delta E_c = E_{cB} - E_{cA} = \sum W_{A\to B}(\vec{F}_{\text{ext}})$ ;
\item $E_{cB} - E_{cA} = \sum \vec{F}$ ;
\item $E_{cA} = E_{cB}$ toujours.
\end{enumerate}
\end{enumerate}
\end{exercice}

\begin{exercice}[Vrai ou faux, avec justification]
Réponds par vrai ou faux en justifiant chaque réponse par une phrase ou un calcul.
\begin{enumerate}
\item L'énergie cinétique d'un solide peut être négative.
\item Le poids d'un solide qui se déplace horizontalement ne travaille pas.
\item Une roue qui roule sans glisser possède à la fois une énergie cinétique de translation et une énergie cinétique de rotation.
\item Le travail des forces de frottement est toujours nul.
\item Si la somme des travaux des forces extérieures appliquées à un solide est nulle entre deux points, sa vitesse est constante entre ces deux points.
\end{enumerate}
\end{exercice}

\begin{exercice}[Le benskin de Douala]
Une moto-taxi (benskin) a une masse de \SI{180}{kg} et transporte son pilote de \SI{70}{kg}.
\begin{enumerate}
\item Rappeler l'expression de l'énergie cinétique d'un solide en translation en précisant la signification et l'unité de chaque grandeur.
\item Calculer l'énergie cinétique de l'ensemble \{moto + pilote\} roulant à \SI{54}{km/h} (penser à convertir la vitesse en m/s).
\item Calculer l'énergie cinétique du même ensemble roulant à \SI{108}{km/h}.
\item Par quel facteur l'énergie cinétique est-elle multipliée lorsque la vitesse double ? Commenter en une phrase la conséquence pour la sécurité routière.
\end{enumerate}
\end{exercice}

\begin{exercice}[Calculs directs]
\begin{enumerate}
\item Un ballon de football de masse $m = \SI{450}{g}$ est frappé à la vitesse $v = \SI{25}{m/s}$. Calculer son énergie cinétique.
\item Un disque plein de masse $M = \SI{10}{kg}$ et de rayon $R = \SI{20}{cm}$ tourne autour de son axe à la vitesse angulaire $\omega = \SI{50}{rad/s}$. On rappelle que $J_{\Delta} = \dfrac{1}{2}MR^2$. Calculer son moment d'inertie puis son énergie cinétique.
\item Un solide de \SI{2}{kg} voit sa vitesse passer de \SI{3}{m/s} à \SI{7}{m/s} sur un trajet rectiligne. Calculer la variation de son énergie cinétique et en déduire la somme des travaux des forces extérieures qu'il a subies.
\end{enumerate}
\end{exercice}

\courssub{Je m'entraîne}

\begin{exercice}[Le palet sur la piste]
Un palet de masse $m = \SI{2}{kg}$ est lancé avec une vitesse initiale $v_0 = \SI{10}{m/s}$ sur une piste horizontale. Il s'arrête après avoir parcouru une distance $d = \SI{25}{m}$. Les frottements sont assimilés à une force $\vec{f}$ constante, colinéaire au déplacement et de sens opposé.
\begin{enumerate}
\item Faire le bilan des forces extérieures appliquées au palet et préciser lesquelles travaillent.
\item En appliquant le théorème de l'énergie cinétique entre le point de lancement et le point d'arrêt, déterminer l'intensité $f$ de la force de frottement.
\end{enumerate}
\end{exercice}

\begin{exercice}[Le volant du groupe électrogène]
Dans un quartier de Yaoundé, un groupe électrogène possède un volant assimilé à un disque plein homogène de masse $M = \SI{25}{kg}$ et de rayon $R = \SI{40}{cm}$, tournant autour de son axe à $N = \SI{1200}{tr/min}$. On donne $J_{\Delta} = \dfrac{1}{2}MR^2$.
\begin{enumerate}
\item Convertir la fréquence de rotation en rad/s et calculer le moment d'inertie $J_{\Delta}$.
\item Calculer l'énergie cinétique du volant.
\item À la coupure du moteur, le volant n'est plus soumis qu'à un couple résistant de moment constant $\mathcal{M}_r = -\SI{2}{N.m}$. En appliquant le théorème de l'énergie cinétique, déterminer l'angle total $\theta$ (en rad, puis en tours) parcouru par le volant avant l'arrêt.
\end{enumerate}
\end{exercice}

\begin{exercice}[La roue qui roule sans glisser]
Une roue pleine (cylindre homogène) de masse $m = \SI{20}{kg}$ et de rayon $R = \SI{30}{cm}$ roule sans glisser sur une route horizontale ; son centre d'inertie G se déplace à la vitesse $v_G = \SI{8}{m/s}$. On rappelle : $J_{\Delta} = \dfrac{1}{2}mR^2$ et, en cas de roulement sans glissement, $\omega = \dfrac{v_G}{R}$.
\begin{enumerate}
\item Calculer l'énergie cinétique de translation de la roue.
\item Calculer la vitesse angulaire $\omega$ puis l'énergie cinétique de rotation.
\item En déduire l'énergie cinétique totale de la roue.
\item Quel pourcentage de l'énergie cinétique totale est dû à la rotation ? Commenter.
\end{enumerate}
\end{exercice}

\begin{exercice}[Le camion dans la descente]
Un camion de masse $m = \SI{10}{t}$ aborde, à la vitesse $v_A = \SI{60}{km/h}$, une descente rectiligne de longueur $L = \SI{500}{m}$ et de pente $6\,\%$ (la dénivellation $h$ vérifie $h = 0{,}06\,L$). L'ensemble des forces de freinage et de résistance est assimilé à une force constante d'intensité $F = \SI{15}{kN}$, opposée au mouvement. On prend $g = \SI{10}{N/kg}$.
\begin{enumerate}
\item Calculer la dénivellation $h$ entre le haut et le bas de la descente.
\item Exprimer puis calculer le travail du poids et le travail de la force résistante sur la descente.
\item En appliquant le théorème de l'énergie cinétique, déterminer la vitesse $v_B$ du camion en bas de la descente (en m/s puis en km/h).
\item Conclure en une phrase sur la sécurité : le freinage est-il suffisant ?
\end{enumerate}
\end{exercice}

\courssub{Je me prépare au Probatoire}

\begin{exercice}[Plan incliné et piste horizontale (type Probatoire)]
Un solide (S) de masse $m = \SI{500}{g}$, assimilé à un point matériel, est lâché sans vitesse initiale du point A, sommet d'un plan incliné de longueur $AB = L = \SI{2}{m}$ faisant un angle $\alpha = 30^{\circ}$ avec l'horizontale. Sur tout le trajet, les frottements équivalent à une force constante d'intensité $f = \SI{0,8}{N}$, opposée au mouvement. En B, le plan incliné est raccordé à une piste horizontale rugueuse (même force de frottement $f$) sur laquelle le solide s'arrête en C. On prend $g = \SI{10}{N/kg}$ ; $\sin 30^{\circ} = 0{,}5$.
\begin{enumerate}
\item Faire le bilan des forces appliquées au solide sur le plan incliné et les représenter sur un schéma.
\item Exprimer le travail du poids et le travail de la force de frottement entre A et B en fonction de $m$, $g$, $L$, $\alpha$ et $f$.
\item En appliquant le théorème de l'énergie cinétique entre A et B, calculer la vitesse $v_B$ du solide en bas du plan.
\item En appliquant à nouveau le théorème de l'énergie cinétique entre B et C, déterminer la distance $BC$ parcourue sur la partie horizontale avant l'arrêt.
\item Quelle serait la valeur de $v_B$ si les frottements étaient négligeables sur le plan incliné ? Conclure sur le rôle des frottements.
\end{enumerate}
\end{exercice}

\begin{exercice}[Le bus de la ligne Yaoundé--Bafoussam (type Probatoire)]
Un bus de masse $m = \SI{9}{t}$ démarre sur une route horizontale rectiligne et atteint la vitesse $v = \SI{72}{km/h}$ en $\Delta t = \SI{20}{s}$. On suppose la somme des forces de résistance (frottements, air) négligeable pendant cette phase, et le travail de la force motrice constant. On prend $g = \SI{10}{N/kg}$.
\begin{enumerate}
\item Convertir la vitesse $v$ en m/s et calculer l'énergie cinétique du bus à \SI{72}{km/h}.
\item En appliquant le théorème de l'énergie cinétique entre l'instant du départ et l'instant où le bus atteint \SI{72}{km/h}, déterminer le travail $W$ de la force motrice.
\item En déduire la puissance moyenne minimale $P = \dfrac{W}{\Delta t}$ développée par le moteur pendant cette phase. Exprimer le résultat en kW.
\item Le bus est maintenant surchargé de \SI{2}{t} (passagers et bagages supplémentaires). En reprenant le même raisonnement, calculer la nouvelle puissance moyenne minimale nécessaire pour atteindre \SI{72}{km/h} en \SI{20}{s}.
\item En réalité, les forces de résistance ne sont pas négligeables : la puissance réelle du moteur doit-elle être supérieure ou inférieure à la puissance minimale calculée ? Justifier à l'aide du théorème de l'énergie cinétique.
\end{enumerate}
\end{exercice}

\begin{exercice}[Le forage et la chute du marteau (type Probatoire)]
Sur un chantier de forage d'un puits villageois dans la région de l'Extrême-Nord, un marteau de masse $m = \SI{200}{kg}$ est lâché sans vitesse initiale d'une hauteur $h = \SI{5}{m}$ au-dessus d'un pieu. Après avoir frappé le pieu, le marteau et le pieu s'enfoncent ensemble de $d = \SI{10}{cm}$ dans le sol, la force de résistance du sol étant supposée constante d'intensité $F$. On néglige la résistance de l'air et on prend $g = \SI{10}{N/kg}$.
\begin{enumerate}
\item En appliquant le théorème de l'énergie cinétique pendant la chute libre du marteau, calculer sa vitesse $v$ juste avant l'impact sur le pieu.
\item Justifier que, pendant l'enfoncement, le poids de l'ensemble travaille encore, et exprimer son travail sur la distance $d$.
\item En appliquant le théorème de l'énergie cinétique à l'ensemble \{marteau + pieu\} entre l'instant de l'impact (vitesse $v$) et l'arrêt après enfoncement de $d$, établir l'expression de $F$ en fonction de $m$, $g$, $h$ et $d$ (on assimilera l'énergie cinétique de l'ensemble juste après l'impact à celle du marteau juste avant l'impact).
\item Calculer $F$. Comparer cette force au poids du marteau et commenter l'efficacité du procédé.
\end{enumerate}
\end{exercice}

\newpage

\coursec{Corrigés des exercices}

\courssub{Je vérifie mes connaissances}

\begin{corrige}[Exercice 1 — QCM : notions fondamentales]
\begin{enumerate}
\item \textbf{Réponse : b) le joule (J).} \\
L'énergie cinétique est une énergie ; son unité SI est le joule (J). Le watt (W) est l'unité de la puissance, le newton (N) celle de la force, et le newton-mètre par seconde (N$\cdot$m/s) est homogène à un watt.
\item \textbf{Réponse : c) 9.} \\
L'énergie cinétique de translation s'écrit $E_c = \frac{1}{2}mv^2$. Si $v$ est multipliée par 3, $v^2$ est multipliée par $3^2 = 9$, donc $E_c$ est multipliée par 9.
\item \textbf{Réponse : b) $E_c = \dfrac{1}{2}J_{\Delta}\,\omega^2$.} \\
C'est l'expression exacte de l'énergie cinétique de rotation pour un solide tournant autour d'un axe fixe $\Delta$ de moment d'inertie $J_{\Delta}$ à la vitesse angulaire $\omega$. L'option a) est dimensionnellement fausse ($\omega$ au lieu de $\omega^2$), l'option c) manque le facteur $1/2$, l'option d) confond avec la translation ($m$ au lieu de $J_{\Delta}$).
\item \textbf{Réponse : b) $\Delta E_c = E_{cB} - E_{cA} = \sum W_{A\to B}(\vec{F}_{\text{ext}})$.} \\
C'est l'énoncé rigoureux du théorème de l'énergie cinétique : la variation de l'énergie cinétique entre deux positions A et B est égale à la somme des travaux des forces \emph{extérieures} le long de la trajectoire entre A et B. L'option a) a un signe erroné, l'option c) confond travail et somme des forces, l'option d) n'est vrai que si la somme des travaux est nulle.
\end{enumerate}
\end{corrige}

\begin{corrige}[Exercice 2 — Vrai ou faux, avec justification]
\begin{enumerate}
\item \textbf{Faux.} L'énergie cinétique d'un solide est définie par $E_c = \frac{1}{2}mv_G^2 + \frac{1}{2}J_{\Delta}\omega^2$ (somme de termes positifs ou nuls). Elle est toujours \textbf{positive ou nulle}, jamais négative.
\item \textbf{Vrai.} Le poids $\vec{P}=m\vec{g}$ est vertical. Si le déplacement est horizontal, l'angle entre $\vec{P}$ et le vecteur déplacement est de $90^{\circ}$. Le travail $W = \vec{P}\cdot\vec{dr} = P\,dr\cos 90^{\circ} = 0$.
\item \textbf{Vrai.} Une roue qui roule sans glisser est en mouvement composé : son centre d'inertie G a une vitesse de translation $\vec{v}_G$ (énergie cinétique de translation $\frac{1}{2}mv_G^2$) et la roue tourne autour de G avec une vitesse angulaire $\omega$ (énergie cinétique de rotation $\frac{1}{2}J_G\omega^2$). L'énergie cinétique totale est la somme des deux.
\item \textbf{Faux.} Les forces de frottement s'opposent généralement au mouvement relatif (frottements de glissement) : leur travail est alors \textbf{négatif}. Dans certains cas (marche, voiture qui accélère), le frottement statique peut avoir un travail positif. Il n'est \emph{jamais} « toujours nul ».
\item \textbf{Faux.} Si $\sum W_{\text{ext}} = 0$, le théorème donne $\Delta E_c = 0$, donc $E_{cB}=E_{cA}$. La \textbf{vitesse (norme du vecteur vitesse) est constante}, mais le \emph{vecteur vitesse} peut changer de direction (exemple : mouvement circulaire uniforme sous l'action d'une force centripète dont le travail est nul).
\end{enumerate}
\end{corrige}

\begin{corrige}[Exercice 3 — Le benskin de Douala]
\begin{enumerate}
\item L'énergie cinétique d'un solide en translation de masse $m$ (en kg) dont le centre d'inertie se déplace à la vitesse $v$ (en m/s) est :
\[
E_c = \frac{1}{2}\,m\,v^2 \quad \text{(en joules, J)}.
\]
\item Masse totale $m = 180 + 70 = \SI{250}{kg}$.
Vitesse $v_1 = \SI{54}{km/h} = 54 \times \frac{1000}{3600} = \SI{15}{m/s}$.
\[
E_{c1} = \frac{1}{2} \times 250 \times 15^2 = 125 \times 225 = \SI{28125}{J} \approx \SI{28,1}{kJ}.
\]
\item Vitesse $v_2 = \SI{108}{km/h} = \SI{30}{m/s}$.
\[
E_{c2} = \frac{1}{2} \times 250 \times 30^2 = 125 \times 900 = \SI{112500}{J} = \SI{112,5}{kJ}.
\]
\item Le facteur est $\dfrac{E_{c2}}{E_{c1}} = \dfrac{112500}{28125} = 4$.
\textbf{Conséquence :} Lorsque la vitesse double, l'énergie cinétique (et donc le travail de freinage nécessaire pour s'arrêter) est multipliée par 4. La distance de freinage augmente donc très fortement, ce qui aggrave considérablement les risques d'accident.
\end{enumerate}
\end{corrige}

\begin{corrige}[Exercice 4 — Calculs directs]
\begin{enumerate}
\item $m = \SI{450}{g} = \SI{0,45}{kg}$, $v = \SI{25}{m/s}$.
\[
E_c = \frac{1}{2} \times \num{0,45} \times 25^2 = \num{0,225} \times 625 = \SI{140,625}{J} \approx \SI{141}{J}.
\]
\item Disque plein : $M = \SI{10}{kg}$, $R = \SI{20}{cm} = \SI{0,20}{m}$.
Moment d'inertie : $J_{\Delta} = \frac{1}{2}MR^2 = \frac{1}{2} \times 10 \times \num{0,20}^2 = 5 \times \num{0,04} = \SI{0,20}{kg\cdot m^2}$.
$\omega = \SI{50}{rad/s}$.
\[
E_c = \frac{1}{2}J_{\Delta}\omega^2 = \frac{1}{2} \times \num{0,20} \times 50^2 = \num{0,10} \times 2500 = \SI{250}{J}.
\]
\item $m = \SI{2}{kg}$, $v_A = \SI{3}{m/s}$, $v_B = \SI{7}{m/s}$.
\[
\Delta E_c = \frac{1}{2}m(v_B^2 - v_A^2) = \frac{1}{2} \times 2 \times (7^2 - 3^2) = 1 \times (49 - 9) = \SI{40}{J}.
\]
D'après le théorème de l'énergie cinétique, $\Delta E_c = \sum W_{\text{ext}}$. La somme des travaux des forces extérieures est donc \textbf{$\SI{40}{J}$}.
\end{enumerate}
\end{corrige}

\courssub{Je m'entraîne}

\begin{corrige}[Exercice 5 — Le palet sur la piste]
\begin{enumerate}
\item \textbf{Bilan des forces :}
\begin{itemize}
\item Poids $\vec{P} = m\vec{g}$ (vertical, vers le bas).
\item Réaction normale du sol $\vec{N}$ (verticale, vers le haut).
\item Force de frottement $\vec{f}$ (horizontale, opposée au mouvement).
\end{itemize}
\textbf{Forces qui travaillent :} Seule la force de frottement $\vec{f}$ travaille car elle est colinéaire au déplacement. Le poids et la réaction normale sont perpendiculaires au déplacement : leur travail est nul.
\item Application du théorème de l'énergie cinétique entre le lancement (A, $v_A = v_0$) et l'arrêt (B, $v_B = 0$) :
\[
E_{cB} - E_{cA} = W_{A\to B}(\vec{f}) + \underbrace{W(\vec{P})}_{=0} + \underbrace{W(\vec{N})}_{=0}.
\]
\[
0 - \frac{1}{2}mv_0^2 = -f \times d \quad (\text{car } \vec{f} \text{ s'oppose au déplacement}).
\]
\[
f = \frac{mv_0^2}{2d} = \frac{2 \times 10^2}{2 \times 25} = \frac{200}{50} = \SI{4}{N}.
\]
L'intensité de la force de frottement est \textbf{$\SI{4}{N}$}.
\end{enumerate}
\end{corrige}

\begin{corrige}[Exercice 6 — Le volant du groupe électrogène]
\begin{enumerate}
\item Conversion : $N = \SI{1200}{tr/min}$.
\[
\omega = 1200 \times \frac{2\pi}{60} = 20 \times 2\pi = 40\pi \approx \SI{125,7}{rad/s}.
\]
Moment d'inertie (disque plein, axe $\Delta$ central) : $J_{\Delta} = \frac{1}{2}MR^2$.
$M = \SI{25}{kg}$, $R = \SI{40}{cm} = \SI{0,40}{m}$.
\[
J_{\Delta} = \frac{1}{2} \times 25 \times \num{0,40}^2 = 12,5 \times \num{0,16} = \SI{2,0}{kg\cdot m^2}.
\]
\item Énergie cinétique de rotation :
\[
E_c = \frac{1}{2}J_{\Delta}\omega^2 = \frac{1}{2} \times \num{2,0} \times (40\pi)^2 = 1 \times 1600\pi^2 = 1600\pi^2 \approx \SI{15791}{J} \approx \SI{15,8}{kJ}.
\]
\item Après coupure, seul le couple résistant $\mathcal{M}_r = -\SI{2}{N\cdot m}$ travaille (poids et réaction de l'axe ont un travail nul car leurs points d'application ne se déplacent pas ou sont perpendiculaires à l'axe).
Théorème de l'énergie cinétique (forme rotation) : $\Delta E_c = W_{\text{ext}}$.
\[
0 - E_c = \mathcal{M}_r \times \theta \quad (\theta \text{ en radians}).
\]
\[
-1600\pi^2 = -2 \times \theta \quad \Rightarrow \quad \theta = \frac{1600\pi^2}{2} = 800\pi^2 \approx \SI{7896}{rad}.
\]
Nombre de tours : $n = \dfrac{\theta}{2\pi} = \dfrac{800\pi^2}{2\pi} = 400\pi \approx \SI{1257}{tr}$.
Le volant effectue environ \textbf{1257 tours} avant de s'arrêter.
\end{enumerate}
\end{corrige}

\begin{corrige}[Exercice 7 — La roue qui roule sans glisser]
\begin{enumerate}
\item Énergie cinétique de translation :
$m = \SI{20}{kg}$, $v_G = \SI{8}{m/s}$.
\[
E_{c,\text{trans}} = \frac{1}{2}mv_G^2 = \frac{1}{2} \times 20 \times 8^2 = 10 \times 64 = \SI{640}{J}.
\]
\item Vitesse angulaire (roulement sans glissement) : $\omega = \dfrac{v_G}{R} = \dfrac{8}{0,30} = \dfrac{80}{3} \approx \SI{26,7}{rad/s}$.
Moment d'inertie (cylindre plein) : $J_{\Delta} = \frac{1}{2}mR^2 = \frac{1}{2} \times 20 \times \num{0,30}^2 = 10 \times \num{0,09} = \SI{0,90}{kg\cdot m^2}$.
Énergie cinétique de rotation :
\[
E_{c,\text{rot}} = \frac{1}{2}J_{\Delta}\omega^2 = \frac{1}{2} \times \num{0,90} \times \left(\frac{80}{3}\right)^2 = \num{0,45} \times \frac{6400}{9} = \num{0,45} \times 711,\overline{1} = \SI{320}{J}.
\]
(On peut aussi utiliser la relation directe pour un cylindre plein : $E_{c,\text{rot}} = \frac{1}{4}mv_G^2 = \frac{1}{4} \times 20 \times 64 = \SI{320}{J}$.)
\item Énergie cinétique totale :
\[
E_{c,\text{tot}} = E_{c,\text{trans}} + E_{c,\text{rot}} = 640 + 320 = \SI{960}{J}.
\]
\item Pourcentage dû à la rotation :
\[
\frac{E_{c,\text{rot}}}{E_{c,\text{tot}}} \times 100 = \frac{320}{960} \times 100 = \frac{1}{3} \times 100 \approx \SI{33,3}{\%}.
\]
\textbf{Commentaire :} Pour un cylindre (ou disque) plein homogène roulant sans glisser, l'énergie cinétique de rotation représente toujours exactement 1/3 de l'énergie cinétique totale (et 1/2 de l'énergie de translation), indépendamment de la vitesse ou de la masse. C'est une conséquence du rapport $J_{\Delta}/(mR^2) = 1/2$.
\end{enumerate}
\end{corrige}

\begin{corrige}[Exercice 8 — Le camion dans la descente]
\begin{enumerate}
\item Dénivellation : $h = 0,06 \times L = 0,06 \times 500 = \SI{30}{m}$.
\item Travail du poids : $W_P = mgh = 10000 \times 10 \times 30 = \SI{3,0 \times 10^6}{J} = \SI{3,0}{MJ}$ (moteur, positif).
Travail de la force résistante : $W_F = -F \times L = -15000 \times 500 = \SI{-7,5 \times 10^6}{J} = \SI{-7,5}{MJ}$ (résistant, négatif).
\item Théorème de l'énergie cinétique entre A (haut) et B (bas) :
$v_A = \SI{60}{km/h} = \frac{60}{3,6} = \SI{16,67}{m/s}$ (exactement $50/3$ m/s).
\[
E_{cA} = \frac{1}{2}mv_A^2 = \frac{1}{2} \times 10000 \times \left(\frac{50}{3}\right)^2 = 5000 \times \frac{2500}{9} = \frac{12,5 \times 10^6}{9} \approx \SI{1,389 \times 10^6}{J}.
\]
\[
\Delta E_c = E_{cB} - E_{cA} = W_P + W_F = 3,0 \times 10^6 - 7,5 \times 10^6 = -4,5 \times 10^6 \text{ J}.
\]
\[
E_{cB} = E_{cA} - 4,5 \times 10^6 \approx 1,389 \times 10^6 - 4,5 \times 10^6 = -3,111 \times 10^6 \text{ J}.
\]
L'énergie cinétique ne peut pas être négative. Cela signifie que le camion \textbf{s'arrête avant le bas de la descente}.
Calcul de la distance d'arrêt $x$ (mesurée depuis le haut) :
$W_P(x) = mg(0,06x) = 6000x$, $W_F(x) = -15000x$.
$\Delta E_c = -E_{cA} = 6000x - 15000x = -9000x$.
$x = \frac{E_{cA}}{9000} = \frac{1,389 \times 10^6}{9000} \approx \SI{154}{m} < \SI{500}{m}$.
Le camion s'arrête après environ 154 m. Sa vitesse en bas de la descente (s'il y arrivait) serait nulle, mais physiquement il est déjà arrêté.
\textbf{Vitesse $v_B = \SI{0}{m/s}$ (arrêt avant le bas).}
\item \textbf{Conclusion sécurité :} Le freinage est \textbf{plus que suffisant} pour arrêter le camion bien avant la fin de la descente (au bout de 154 m sur 500 m). Le conducteur n'a pas besoin de freiner davantage ; il doit cependant veiller à ne pas laisser le camion redémarrer dans la pente après l'arrêt.
\end{enumerate}
\end{corrige}

\courssub{Je me prépare au Probatoire}

\begin{corrige}[Exercice 9 — Plan incliné et piste horizontale (type Probatoire)]
\begin{enumerate}
\item \textbf{Bilan des forces sur le plan incliné (schéma à faire par l'élève) :}
\begin{itemize}
\item Poids $\vec{P} = m\vec{g}$ (vertical vers le bas).
\item Réaction normale $\vec{N}$ (perpendiculaire au plan, vers le haut).
\item Force de frottement $\vec{f}$ (parallèle au plan, vers le haut, opposée au mouvement).
\end{itemize}
Sur un schéma : repère $(x,y)$ avec $x$ selon la plus grande pente vers le bas. $\vec{P}$ décomposé en $P_x = mg\sin\alpha$ (moteur) et $P_y = -mg\cos\alpha$. $\vec{N}$ sur $Oy$ positif. $\vec{f}$ sur $Ox$ négatif.
\item Travaux entre A et B (longueur $L$) :
\begin{itemize}
\item Poids : $W_{AB}(\vec{P}) = mg \times \text{dénivellation} = mgL\sin\alpha$. (Moteur, positif).
\item Frottement : $W_{AB}(\vec{f}) = -f \times L$. (Résistant, négatif).
\item Réaction normale : travail nul (perpendiculaire au déplacement).
\end{itemize}
\item Théorème de l'énergie cinétique entre A ($v_A=0$) et B ($v_B$) :
\[
\frac{1}{2}mv_B^2 - 0 = mgL\sin\alpha - fL.
\]
Données : $m = \SI{0,5}{kg}$, $g = \SI{10}{N/kg}$, $L = \SI{2}{m}$, $\sin\alpha = 0,5$, $f = \SI{0,8}{N}$.
\[
mgL\sin\alpha = \num{0,5} \times 10 \times 2 \times \num{0,5} = \SI{5}{J}.
\]
\[
fL = \num{0,8} \times 2 = \SI{1,6}{J}.
\]
\[
\frac{1}{2} \times \num{0,5} \times v_B^2 = 5 - 1,6 = 3,4 \quad \Rightarrow \quad \num{0,25} v_B^2 = 3,4 \quad \Rightarrow \quad v_B^2 = 13,6.
\]
\[
v_B = \sqrt{13,6} \approx \SI{3,69}{m/s}.
\]
\item Théorème entre B ($v_B$) et C ($v_C=0$) sur l'horizontale (distance $BC = x$) :
Forces : Poids et Normale (travail nul), Frottement $\vec{f}$ (même intensité, opposé au déplacement).
\[
0 - \frac{1}{2}mv_B^2 = -f x \quad \Rightarrow \quad x = \frac{mv_B^2}{2f} = \frac{2 \times 3,4}{2 \times \num{0,8}} = \frac{3,4}{\num{0,8}} = \SI{4,25}{m}.
\]
(On a utilisé $\frac{1}{2}mv_B^2 = 3,4$ J trouvé à la question précédente).
\item Sans frottement sur le plan ($f=0$) :
\[
\frac{1}{2}mv_B'^2 = mgL\sin\alpha = 5 \text{ J} \quad \Rightarrow \quad v_B'^2 = \frac{10}{\num{0,5}} = 20 \quad \Rightarrow \quad v_B' = \sqrt{20} \approx \SI{4,47}{m/s}.
\]
\textbf{Conclusion :} Les frottements dissipent une partie de l'énergie potentielle de pesanteur en chaleur (travail négatif), ce qui réduit l'énergie cinétique acquise et donc la vitesse en bas du plan (de 4,47 à 3,69 m/s ici). Ils réduisent aussi la distance d'arrêt sur l'horizontale.
\end{enumerate}
\textbf{Barème indicatif (sur 6-7 pts) :} Bilan des forces + schéma (1,5 pt) ; Expression travaux (1 pt) ; Calcul $v_B$ (1,5 pt) ; Calcul $BC$ (1,5 pt) ; Cas sans frottement + conclusion (1 pt).
\textbf{Points de vigilance :} Ne pas oublier de convertir $m$ en kg. Bien distinguer $\sin\alpha$ et $\cos\alpha$ pour le travail du poids. Citer explicitement le théorème $\Delta E_c = \sum W$. Vérifier les signes des travaux.
\end{corrige}

\begin{corrige}[Exercice 10 — Le bus de la ligne Yaoundé–Bafoussam (type Probatoire)]
\begin{enumerate}
\item $v = \SI{72}{km/h} = \frac{72}{3,6} = \SI{20}{m/s}$.
$m = \SI{9}{t} = \SI{9000}{kg}$.
\[
E_c = \frac{1}{2}mv^2 = \frac{1}{2} \times 9000 \times 20^2 = 4500 \times 400 = \SI{1,8 \times 10^6}{J} = \SI{1,8}{MJ}.
\]
\item Théorème de l'énergie cinétique (départ $v=0$ à $v=20$ m/s) :
Forces : Force motrice $\vec{F}_m$ (moteur), résistances négligées, Poids et Normale (travail nul).
\[
\Delta E_c = E_c - 0 = W(\vec{F}_m) \quad \Rightarrow \quad W = \SI{1,8 \times 10^6}{J}.
\]
\item Puissance moyenne minimale :
\[
P = \frac{W}{\Delta t} = \frac{1,8 \times 10^6}{20} = \SI{90000}{W} = \SI{90}{kW}.
\]
\item Nouveau masse $m' = 9 + 2 = \SI{11}{t} = \SI{11000}{kg}$.
Nouvelle énergie cinétique à 20 m/s : $E_c' = \frac{1}{2} \times 11000 \times 400 = \SI{2,2 \times 10^6}{J}$.
Travail moteur nécessaire $W' = \SI{2,2 \times 10^6}{J}$.
Puissance moyenne minimale $P' = \frac{2,2 \times 10^6}{20} = \SI{110000}{W} = \SI{110}{kW}$.
\item \textbf{Supérieure.}
Justification : En réalité, les forces de résistance (frottements, air) effectuent un travail négatif $W_{\text{rés}} < 0$. Le théorème s'écrit : $W_{\text{moteur}} + W_{\text{rés}} = \Delta E_c$. Donc $W_{\text{moteur}} = \Delta E_c - W_{\text{rés}} > \Delta E_c$. La puissance réelle $P_{\text{réelle}} = W_{\text{moteur}}/\Delta t$ doit donc être \textbf{strictement supérieure} à la puissance minimale calculée sans résistances.
\end{enumerate}
\textbf{Barème indicatif (sur 6-7 pts) :} Conversion vitesse + $E_c$ (1 pt) ; Travail moteur (1 pt) ; Puissance moyenne (1 pt) ; Cas surchargé (1,5 pt) ; Justification résistances (1,5 pt).
\textbf{Points de vigilance :} Conversion km/h $\to$ m/s obligatoire ($/3,6$). Unités : tonnes $\to$ kg. Préciser « puissance \emph{minimale} » car on a négligé les résistances. Pour la dernière question, utiliser l'équation du théorème avec les signes.
\end{corrige}

\begin{corrige}[Exercice 11 — Le forage et la chute du marteau (type Probatoire)]
\begin{enumerate}
\item Chute libre (hauteur $h = \SI{5}{m}$, $v_0 = 0$). Forces : Poids (moteur), résistance de l'air négligée.
Théorème : $\frac{1}{2}mv^2 - 0 = mgh$.
\[
v = \sqrt{2gh} = \sqrt{2 \times 10 \times 5} = \sqrt{100} = \SI{10}{m/s}.
\]
Énergie cinétique avant impact : $E_c = mgh = 200 \times 10 \times 5 = \SI{10000}{J}$.
\item Pendant l'enfoncement (distance $d = \SI{0,10}{m}$), le marteau (et le pieu) descendent encore. Le poids $\vec{P} = m\vec{g}$ reste vertical et le déplacement est vertical vers le bas : l'angle est $0^{\circ}$, le travail est \textbf{positif}.
Travail du poids sur $d$ : $W_P = mgd = 200 \times 10 \times \num{0,10} = \SI{200}{J}$.
\item Enfoncement : Système \{Marteau + Pieu\}. On assimile $E_c(\text{juste après choc}) \approx E_c(\text{marteau avant choc}) = \SI{10000}{J}$.
Forces extérieures travaillant : Poids de l'ensemble (on prend la masse $m$ du marteau, le pieu étant au repos initialement et sa masse non donnée, l'énoncé dit « le poids de l'ensemble travaille encore » en donnant $m$ pour le marteau, on utilise $m$) $\to$ travail $+mgd$. Force de résistance du sol $\vec{F}$ (vers le haut, opposée au déplacement) $\to$ travail $-Fd$.
Arrêt final : $v_f = 0$.
Théorème :
\[
0 - E_c = mgd - Fd.
\]
\[
Fd = E_c + mgd = mgh + mgd = mg(h+d).
\]
\[
F = \frac{mg(h+d)}{d} = mg\left(1 + \frac{h}{d}\right).
\]
\item Calcul numérique :
\[
F = 200 \times 10 \times \left(1 + \frac{5}{\num{0,10}}\right) = 2000 \times (1 + 50) = 2000 \times 51 = \SI{102000}{N} = \SI{102}{kN}.
\]
Poids du marteau : $P = mg = \SI{2000}{N}$.
Rapport : $\frac{F}{P} = 51$.
\textbf{Commentaire :} La force de résistance du sol (et donc la force exercée sur le pieu) est \textbf{51 fois plus grande} que le poids du marteau. C'est l'efficacité du principe du marteau-piqueur : l'énergie potentielle accumulée pendant la chute libre (sur une grande distance $h$) est restituée sur une très petite distance $d$, multipliant l'effort par le rapport $h/d$ (ici 50), augmenté de 1 pour le travail du poids pendant l'enfoncement.
\end{enumerate}
\textbf{Barème indicatif (sur 7-8 pts) :} Chute : théorème + calcul $v$ (1,5 pt) ; Enfoncement : justification travail poids + expression (1,5 pt) ; Théorème enfoncement + expression $F$ (2 pts) ; Calcul $F$ + comparaison + commentaire (2 pts).
\textbf{Points de vigilance :} Ne pas oublier le travail du poids \emph{pendant} l'enfoncement (souvent oublié). Bien écrire $h+d$ au numérateur. Unités : $d$ en mètres (\num{0,10} m). Dire explicitement « on néglige la masse du pieu / on assimile l'énergie cinétique post-choc à celle pré-choc » comme le demande l'énoncé.
\end{corrige}

\newpage

\coursec{Fiche bilan}

\begin{popfigure}
\begin{tikzpicture}[mindmap, grow cyclic, every node/.style={concept, text=white, font=\small\bfseries},
root concept/.append style={concept color=popink, font=\large\bfseries},
level 1/.append style={level distance=3.6cm, sibling angle=60, minimum size=2.1cm},
level 2/.append style={level distance=2.4cm, sibling angle=45, minimum size=1.7cm, font=\scriptsize}]
\node[root concept]{Énergie cinétique et théorème de l'énergie cinétique}
  child[concept color=popblue]{ node[align=center]{Translation \\ $E_c=\dfrac{1}{2}mv^2$}
    child[concept color=popblueL, text=popdark]{ node{$v$ en m/s}}
    child[concept color=popblueL, text=popdark]{ node{$E_c$ en joule (J)}}
  }
  child[concept color=poporange]{ node[align=center]{Rotation \\ $E_c=\dfrac{1}{2}J_{\Delta}\omega^2$}
    child[concept color=poporangeL, text=popdark]{ node{$\omega$ en rad/s}}
    child[concept color=poporangeL, text=popdark]{ node{$J_\Delta$ en kg.m$^2$}}
  }
  child[concept color=popgreen]{ node[align=center]{Mouvement combiné \\ $E_c=\dfrac{1}{2}mv_G^2+\dfrac{1}{2}J_\Delta\omega^2$}
    child[concept color=popgreenL, text=popdark]{ node{Th. de König}}
  }
  child[concept color=popgold]{ node[align=center]{Théorème \\ $\Delta E_c=\sum W(\vec{F})$}
    child[concept color=popgoldL, text=popdark]{ node{Travail du poids $mgh$}}
    child[concept color=popgoldL, text=popdark]{ node{Forces $\perp$ : $W=0$}}
  }
  child[concept color=poppurple]{ node{Applications}
    child[concept color=poppurpleL, text=popdark]{ node{Vitesse, freinage}}
    child[concept color=poppurpleL, text=popdark]{ node{Distance, force}}
  }
  child[concept color=poppink]{ node{Méthode}
    child[concept color=poppinkL, text=popdark]{ node{Bilan des forces}}
    child[concept color=poppinkL, text=popdark]{ node{Schéma + unités SI}}
  };
\end{tikzpicture}
\legende{Carte mentale de la leçon : énergie cinétique et théorème de l'énergie cinétique.}
\end{popfigure}

\begin{bilan}
\textbf{Définitions clés}
\begin{itemize}
\item L'\cle{énergie cinétique} d'un solide est l'énergie qu'il possède du fait de son mouvement. Elle s'exprime en \cle{joule} (J) : $1~\mathrm{J}=1~\mathrm{kg}\cdot\mathrm{m}^2\cdot\mathrm{s}^{-2}$.
\item Solide ponctuel ou en \cle{translation} : tous les points ont la même vitesse $v$.
\item Solide en \cle{rotation} autour d'un axe fixe $\Delta$ : tous les points ont la même vitesse angulaire $\omega$.
\item Le \cle{moment d'inertie} $J_\Delta$ mesure la résistance du solide à la mise en rotation ; il dépend de la masse et de sa répartition autour de l'axe.
\end{itemize}

\textbf{Formules à connaître par c\oe ur}
\begin{center}
\begin{tabularx}{\linewidth}{|l|X|X|}
\hline
\rowcolor{popblueL} \textbf{Grandeur} & \textbf{Formule} & \textbf{Unités SI} \\
\hline
$E_c$ (translation) & $E_c=\dfrac{1}{2}mv^2$ & $m$ en kg ; $v$ en m/s ; $E_c$ en J \\
\hline
$E_c$ (rotation) & $E_c=\dfrac{1}{2}J_\Delta\,\omega^2$ & $J_\Delta$ en kg$\cdot$m$^2$ ; $\omega$ en rad/s \\
\hline
$E_c$ (mouvement combiné) & $E_c=\dfrac{1}{2}mv_G^2+\dfrac{1}{2}J_\Delta\,\omega^2$ & $v_G$ : vitesse du centre d'inertie \\
\hline
Théorème de l'énergie cinétique & $E_{c_B}-E_{c_A}=\displaystyle\sum W_{A\to B}(\vec{F})$ & travaux en J \\
\hline
Travail du poids & $W(\vec{P})=mg(z_A-z_B)=mgh$ & $h$ : dénivelé (positif en descente) \\
\hline
Travail d'une force constante & $W(\vec{F})=F\cdot AB\cdot\cos\alpha$ & $\alpha$ : angle entre $\vec{F}$ et le déplacement \\
\hline
\end{tabularx}
\end{center}

\textbf{Propriétés essentielles}
\begin{itemize}
\item Une force \cle{perpendiculaire} au déplacement (réaction normale, tension d'un fil inextensible en rotation) ne travaille pas : $W=0$.
\item Le théorème s'applique à \emph{tout} solide, entre deux positions $A$ et $B$, dans un référentiel galiléen, quelles que soient les forces (conservatives ou non, frottements compris).
\item Si $\sum W>0$ (travail moteur) : la vitesse augmente ; si $\sum W<0$ (travail résistant) : la vitesse diminue ; si $\sum W=0$ : la vitesse est constante en norme.
\item Roulement sans glissement : $v_G=R\,\omega$.
\end{itemize}

\textbf{Méthode express (4 étapes)}
\begin{enumerate}
\item \textbf{Choisir} le système, le référentiel galiléen, les états initial $A$ et final $B$ ; faire un schéma avec le bilan des forces extérieures.
\item \textbf{Convertir} toutes les données en unités SI (km/h $\rightarrow$ m/s : diviser par $3{,}6$ ; tr/min $\rightarrow$ rad/s : multiplier par $\dfrac{2\pi}{60}$).
\item \textbf{Calculer} le travail de chaque force entre $A$ et $B$ (repérer les forces qui ne travaillent pas).
\item \textbf{Appliquer} $\Delta E_c=\sum W(\vec{F})$ et isoler la grandeur cherchée ($v$, $d$, $F$, $h$\dots) ; vérifier le signe et l'ordre de grandeur.
\end{enumerate}
\end{bilan}

\begin{aretenir}[Checklist avant le Probatoire]
\begin{itemize}
\item[$\square$] Je sais définir l'énergie cinétique et citer son unité SI (le joule).
\item[$\square$] Je sais écrire $E_c=\dfrac{1}{2}mv^2$ pour un solide en translation et vérifier l'homogénéité : $\mathrm{kg}\cdot(\mathrm{m/s})^2=\mathrm{J}$.
\item[$\square$] Je sais écrire $E_c=\dfrac{1}{2}J_\Delta\omega^2$ pour un solide en rotation et préciser le rôle du moment d'inertie.
\item[$\square$] Je sais écrire l'énergie cinétique d'un solide en mouvement combiné : $E_c=\dfrac{1}{2}mv_G^2+\dfrac{1}{2}J_\Delta\omega^2$.
\item[$\square$] Je convertis sans erreur : km/h en m/s (diviser par $3{,}6$) et tr/min en rad/s.
\item[$\square$] Je sais énoncer le théorème de l'énergie cinétique : $\Delta E_c=\sum W(\vec{F}_{\mathrm{ext}})$.
\item[$\square$] Je sais calculer le travail du poids $W(\vec{P})=mgh$ avec le bon signe selon la montée ou la descente.
\item[$\square$] Je repère immédiatement les forces qui ne travaillent pas (forces perpendiculaires au déplacement).
\item[$\square$] Je sais exploiter le théorème pour trouver une vitesse, une distance de freinage, une force de frottement ou une hauteur.
\item[$\square$] Je rédige proprement : bilan des forces, expression littérale avant l'application numérique, résultat avec unité et chiffres significatifs.
\item[$\square$] Je vérifie la cohérence physique du résultat (signe du travail, ordre de grandeur de la vitesse).
\end{itemize}
\end{aretenir}

\findecours

\end{document}
